\documentclass[11pt,a4paper]{article}

% ==================== TEMEL PAKETLER (pdflatex) ====================
\usepackage[T1]{fontenc}
\usepackage[utf8]{inputenc}
\usepackage{lmodern}
\usepackage[a4paper,margin=2cm,top=2.4cm,bottom=2.4cm]{geometry}
\usepackage{xcolor}
\usepackage{colortbl}
\usepackage{array}
\usepackage{booktabs}
\usepackage{longtable}
\usepackage{tabularx}
\usepackage{listings}
\usepackage{fancyhdr}
\usepackage{amsmath}
\usepackage{amssymb}
\usepackage{mathtools}
\usepackage{hyperref}

% ==================== RENK PALETİ ====================
\definecolor{anaMavi}{HTML}{1B3A5B}
\definecolor{vurguMavi}{HTML}{2E6DA4}
\definecolor{acikMavi}{HTML}{EAF2F9}
\definecolor{turuncu}{HTML}{D9701E}
\definecolor{yesil}{HTML}{2E7D32}
\definecolor{acikYesil}{HTML}{E8F5E9}
\definecolor{kirmizi}{HTML}{C0392B}
\definecolor{mor}{HTML}{6C3483}
\definecolor{acikMor}{HTML}{F3E9F7}
\definecolor{acikGri}{HTML}{F4F5F7}
\definecolor{ortaGri}{HTML}{6B7280}
\definecolor{koyuGri}{HTML}{2D2D2D}
\definecolor{kodArka}{HTML}{FBFBFA}

% ==================== HYPERREF ====================
\hypersetup{
  colorlinks=true,
  linkcolor={[HTML]{2E6DA4}},
  urlcolor={[HTML]{2E6DA4}},
  pdftitle={C++ ile Veri Yapıları ve Algoritmalar Rehberi},
  pdfauthor={Algoritma Rehberi},
  bookmarksnumbered=true
}

% ==================== BAŞLIK STİLLERİ (titlesec olmadan) ====================
\setcounter{secnumdepth}{0}
\newcounter{secc}
\newcounter{subsecc}[secc]
\newcounter{subsubsecc}[subsecc]
\newcommand{\gsec}[1]{\par\phantomsection\stepcounter{secc}%
  \setcounter{subsecc}{0}%
  \vspace{18pt}{\color{anaMavi}\Large\bfseries\thesecc.\hspace{0.5em}#1}\par
  \vspace{2pt}{\color{turuncu}\rule{\linewidth}{1.6pt}}\par\vspace{8pt}%
  \addcontentsline{toc}{section}{\protect\numberline{\thesecc}#1}\markboth{#1}{#1}}
\newcommand{\gsub}[1]{\par\phantomsection\stepcounter{subsecc}%
  \setcounter{subsubsecc}{0}%
  \vspace{11pt}{\color{vurguMavi}\large\bfseries\thesecc.\thesubsecc\hspace{0.5em}#1}\par\vspace{4pt}%
  \addcontentsline{toc}{subsection}{\protect\numberline{\thesecc.\thesubsecc}#1}}
\newcommand{\gsubsub}[1]{\par\phantomsection\stepcounter{subsubsecc}%
  \vspace{7pt}{\color{koyuGri}\normalsize\bfseries #1}\par\vspace{2pt}}

% Kısım (Part) başlığı
\newcommand{\gpart}[2]{\clearpage
  \begin{center}\vspace*{1.5cm}
  {\color{turuncu}\rule{0.7\linewidth}{1.5pt}}\\[10pt]
  {\color{ortaGri}\Large KISIM #1}\\[8pt]
  {\color{anaMavi}\Huge\bfseries #2}\\[10pt]
  {\color{turuncu}\rule{0.7\linewidth}{1.5pt}}
  \end{center}\vspace{1cm}}

% ==================== BAŞLIK / ALTBİLGİ ====================
\setlength{\headheight}{14pt}
\pagestyle{fancy}
\fancyhf{}
\renewcommand{\headrulewidth}{0.4pt}
\renewcommand{\footrulewidth}{0.4pt}
\fancyhead[L]{\small\color{ortaGri}C++ Veri Yapıları \& Algoritmalar}
\fancyhead[R]{\small\color{ortaGri}\leftmark}
\fancyfoot[C]{\small\color{ortaGri}\thepage}

% ==================== KOD LİSTELEME (C++) ====================
\lstdefinestyle{cppstyle}{
  language=C++,
  basicstyle=\ttfamily\small\color{koyuGri},
  keywordstyle=\color{vurguMavi}\bfseries,
  commentstyle=\color{yesil}\itshape,
  stringstyle=\color{turuncu},
  numberstyle=\tiny\color{ortaGri},
  numbers=left,
  numbersep=8pt,
  backgroundcolor=\color{kodArka},
  frame=single,
  framesep=6pt,
  rulecolor=\color{ortaGri!40},
  breaklines=true,
  showstringspaces=false,
  tabsize=4,
  columns=fullflexible,
  keepspaces=true,
  extendedchars=true,
  literate={ı}{{\i}}1{İ}{{\.I}}1{ş}{{\c s}}1{Ş}{{\c S}}1{ğ}{{\u g}}1{Ğ}{{\u G}}1{ç}{{\c c}}1{Ç}{{\c C}}1{ö}{{\"o}}1{Ö}{{\"O}}1{ü}{{\"u}}1{Ü}{{\"U}}1,
  morekeywords={size_t,uint8_t,uint16_t,uint32_t,uint64_t,int8_t,int16_t,
    int32_t,int64_t,nullptr,constexpr,noexcept,override,final,auto,
    vector,string,pair,map,set,unordered_map,unordered_set,priority_queue,
    queue,stack,deque,array,list,tuple,optional,shared_ptr,unique_ptr,
    template,typename,class,public,private,protected,virtual,namespace,
    using,std,ll,int64,move,make_pair,make_shared,make_unique}
}
\lstset{style=cppstyle}

% Kabuk/terminal stili
\lstdefinestyle{shstyle}{
  basicstyle=\ttfamily\small\color{koyuGri},
  backgroundcolor=\color{koyuGri!6},
  frame=single,framesep=6pt,rulecolor=\color{ortaGri!40},
  breaklines=true,showstringspaces=false,tabsize=4,columns=fullflexible,
  keepspaces=true,
  literate={ı}{{\i}}1{İ}{{\.I}}1{ş}{{\c s}}1{Ş}{{\c S}}1{ğ}{{\u g}}1{Ğ}{{\u G}}1{ç}{{\c c}}1{Ç}{{\c C}}1{ö}{{\"o}}1{Ö}{{\"O}}1{ü}{{\"u}}1{Ü}{{\"U}}1,
}

% ==================== ÖZEL KUTULAR (tcolorbox olmadan) ====================
\newcommand{\callout}[4]{\par\smallskip\noindent\begingroup
  \setlength{\fboxsep}{8pt}\setlength{\fboxrule}{0.8pt}%
  \fcolorbox{#1}{#2}{\parbox{\dimexpr\linewidth-2\fboxsep-2\fboxrule\relax}{%
    {\bfseries\color{#1}#3}\par\smallskip #4}}\endgroup\par\smallskip}
\newcommand{\bilgi}[1]{\callout{vurguMavi}{acikMavi}{Bilgi}{#1}}
\newcommand{\ipucu}[1]{\callout{yesil}{acikYesil}{İpucu}{#1}}
\newcommand{\uyari}[1]{\callout{kirmizi}{kirmizi!7}{Dikkat}{#1}}
\newcommand{\ornek}[2][Örnek]{\callout{ortaGri!90!black}{acikGri}{#1}{#2}}
\newcommand{\notk}[2][Not]{\callout{mor}{acikMor}{#1}{#2}}

% Soru / Çözüm sistemi (alıştırmalar için)
\newcounter{soruc}
\newcommand{\soru}[1]{\par\smallskip\refstepcounter{soruc}%
  \noindent\begingroup\setlength{\fboxsep}{8pt}\setlength{\fboxrule}{0.8pt}%
  \fcolorbox{turuncu}{turuncu!6}{\parbox{\dimexpr\linewidth-2\fboxsep-2\fboxrule\relax}{%
  {\bfseries\color{turuncu}Soru \thesoruc}\par\smallskip #1}}\endgroup\par\smallskip}
\newcommand{\cozum}{\par\noindent{\bfseries\color{yesil}Çözüm.}\ }

% Yardımcı komutlar
\newcommand{\code}[1]{\texttt{\color{koyuGri}#1}}
\newcommand{\bigO}[1]{\ensuremath{\mathcal{O}(#1)}}
% Matematik sembollerini metin modunda da güvenli kıl (ensuremath sarmalayıcı)
\let\svtimes\times   \renewcommand{\times}{\ensuremath{\svtimes}}
\let\svcdot\cdot     \renewcommand{\cdot}{\ensuremath{\svcdot}}
\let\svto\to         \renewcommand{\to}{\ensuremath{\svto}}
\let\svle\le         \renewcommand{\le}{\ensuremath{\svle}}
\let\svge\ge         \renewcommand{\ge}{\ensuremath{\svge}}
\renewcommand{\arraystretch}{1.35}

% ==================== BAŞLANGIÇ ====================
\begin{document}

% ---------- KAPAK ----------
\begin{titlepage}
\centering
\vspace*{1.4cm}
{\color{turuncu}\rule{\textwidth}{2pt}}\\[8pt]
{\color{anaMavi}\Huge\bfseries C++ ile Veri Yapıları}\\[10pt]
{\color{vurguMavi}\LARGE ve Algoritmalar Rehberi}\\[8pt]
{\color{ortaGri}\large Temellerden İleri Tekniklere Kapsamlı Başvuru}\\[6pt]
{\color{turuncu}\rule{\textwidth}{2pt}}\\[1.4cm]

\begin{center}
\fcolorbox{vurguMavi}{acikMavi}{\parbox{0.86\textwidth}{\centering\large
Bu belge, veri yapılarını ve algoritmaları modern C++ (C++17) ile,
her konuyu \textbf{çalışan kod örnekleriyle}, karmaşıklık analizleriyle
ve çizimlerle adım adım anlatır. Temel karmaşıklık analizinden dinamik
programlamaya kadar bir bilgisayar bilimi müfredatının tamamını kapsar.\vspace{4pt}}}
\end{center}

\vspace{0.8cm}
\begin{center}
\fcolorbox{ortaGri!50}{acikGri}{\parbox{0.86\textwidth}{\small
\textbf{Kapsam --- Kısım 1 (Temeller):} algoritma analizi \textbullet\ asimptotik
gösterim (Big-O) \textbullet\ karmaşıklık sınıfları \textbullet\ C++ \& STL araçları.\vspace{4pt}

\textbf{Kapsam --- Kısım 2 (Temel Veri Yapıları):} diziler \& \code{vector}
\textbullet\ bağlı listeler \textbullet\ yığın (stack) \textbullet\ kuyruk \& deque
\textbullet\ hash tabloları.\vspace{4pt}

\textbf{Kapsam --- Kısım 3 (Ağaçlar):} ikili ağaçlar \textbullet\ BST \textbullet\ AVL
\& kırmızı-siyah ağaç \textbullet\ heap \& öncelik kuyruğu \textbullet\ trie
\textbullet\ segment ağacı \& Fenwick.\vspace{4pt}

\textbf{Kapsam --- Kısım 4 (Graflar):} gösterim \textbullet\ BFS/DFS
\textbullet\ en kısa yol \textbullet\ MST \textbullet\ topolojik sıralama
\textbullet\ birleşim-bulma (DSU).\vspace{4pt}

\textbf{Kapsam --- Kısım 5 (Algoritma Teknikleri):} sıralama \textbullet\ arama
\textbullet\ böl-yönet \textbullet\ açgözlü \textbullet\ \textbf{dinamik programlama
(detaylı)} \textbullet\ geri izleme \textbullet\ string \& bit algoritmaları.\vspace{4pt}}}
\end{center}

\vfill
{\color{ortaGri}\large Hazırlanma Tarihi: 12 Temmuz 2026}\\[4pt]
{\color{ortaGri} C++17 \textbullet\ g++ / clang++ \textbullet\ STL}
\vspace*{0.6cm}
\end{titlepage}

% ---------- İÇİNDEKİLER ----------
\tableofcontents
\thispagestyle{empty}
\clearpage

% #####################################################################
\gpart{1}{Temeller ve Algoritma Analizi}
% #####################################################################

% =====================================================================
\gsec{Algoritma Analizi ve Asimptotik Gösterim}
% =====================================================================

Bir algoritmayı değerlendirirken iki temel kaynağı ölçeriz: \textbf{zaman}
(kaç işlem yapılır) ve \textbf{bellek} (ne kadar yer tutulur). Bir programı
saniye cinsinden ölçmek makineye, derleyiciye ve o anki yüke bağlı olduğundan
güvenilir değildir. Bunun yerine, girdi boyutu $n$ büyürken çalışma süresinin
\emph{nasıl büyüdüğünü} tarif ederiz. Bu, asimptotik analizdir.

\gsub{Neden Asimptotik Analiz?}

İki algoritmayı karşılaştırmak istediğimizde, küçük girdilerdeki farklar
genelde önemsizdir; asıl mesele girdi büyüdüğünde ne olacağıdır. Örneğin
$100n$ işlem yapan bir algoritma, $n^2$ işlem yapan bir algoritmadan
$n > 100$ için daha hızlıdır ve $n$ büyüdükçe fark uçurumlaşır.

\begin{center}
\begin{tabularx}{\textwidth}{@{}>{\bfseries\color{anaMavi}}l c c c c@{}}
\toprule
\rowcolor{acikMavi}\color{anaMavi}\textbf{Karmaşıklık} & \color{anaMavi}$n=10$ & \color{anaMavi}$n=100$ & \color{anaMavi}$n=10^3$ & \color{anaMavi}$n=10^6$\\
\midrule
\bigO{1} & 1 & 1 & 1 & 1\\
\rowcolor{acikGri}\bigO{\log n} & 3 & 7 & 10 & 20\\
\bigO{n} & 10 & 100 & 1000 & $10^6$\\
\rowcolor{acikGri}\bigO{n\log n} & 33 & 664 & $10^4$ & $2\!\times\!10^7$\\
\bigO{n^2} & 100 & $10^4$ & $10^6$ & $10^{12}$\\
\rowcolor{acikGri}\bigO{2^n} & 1024 & $10^{30}$ & --- & ---\\
\bottomrule
\end{tabularx}
\end{center}

\bilgi{Rekabetçi programlamada pratik bir kural: modern bir işlemci saniyede
kabaca $10^8$--$10^9$ basit işlem yapar. Girdi $n \le 10^6$ ise \bigO{n} veya
\bigO{n\log n} hedeflenir; $n \le 10^3$ ise \bigO{n^2} kabul edilebilir;
$n \le 20$ ise \bigO{2^n} üstel çözümler bile geçer.}

\gsub{Big-O, Big-Omega ve Big-Theta}

Üç temel gösterim, bir fonksiyonun büyüme hızını sınırlar:

\begin{center}
\begin{tabularx}{\textwidth}{@{}>{\bfseries}c >{\color{anaMavi}\bfseries}l X@{}}
\toprule
\rowcolor{acikMavi}\textbf{Gösterim} & \color{anaMavi}\textbf{Ad} & \color{anaMavi}\textbf{Anlamı}\\
\midrule
\bigO{f(n)} & Üst sınır & En kötü durumda çalışma süresi \emph{en fazla} bu kadar büyür.\\
\rowcolor{acikGri}$\Omega(f(n))$ & Alt sınır & Çalışma süresi \emph{en az} bu kadardır.\\
$\Theta(f(n))$ & Sıkı sınır & Hem üst hem alt sınır; büyüme \emph{tam olarak} bu mertebede.\\
\bottomrule
\end{tabularx}
\end{center}

Pratikte en çok Big-O kullanılır çünkü genelde ``en kötü durumda ne kadar
kötü olabilir?'' sorusunu sorarız. Resmî tanım: $f(n) = O(g(n))$ demek,
öyle pozitif sabitler $c$ ve $n_0$ vardır ki her $n \ge n_0$ için
$f(n) \le c\cdot g(n)$ olur. Yani yeterince büyük $n$ için $f$, $g$'nin
sabit katından fazla büyümez.

\gsub{Sadeleştirme Kuralları}

Big-O ifadelerini sadeleştirirken iki kural yeterlidir:

\begin{center}
\begin{tabularx}{\textwidth}{@{}>{\bfseries\color{anaMavi}}l X@{}}
\toprule
\rowcolor{acikMavi}\color{anaMavi}\textbf{Kural} & \color{anaMavi}\textbf{Açıklama ve örnek}\\
\midrule
Sabitleri at & $O(5n) \to O(n)$, \quad $O(n/2) \to O(n)$.\\
\rowcolor{acikGri}Baskın terimi tut & $O(n^2 + n) \to O(n^2)$, \quad $O(n^2 + n\log n) \to O(n^2)$.\\
\bottomrule
\end{tabularx}
\end{center}

\ornek[İç içe döngüleri analiz etme]{Aşağıdaki üç örnek, döngü yapısından
karmaşıklığın nasıl okunacağını gösterir.}
\begin{lstlisting}
// (1) Tek döngü -> O(n)
for (int i = 0; i < n; ++i)
    sum += a[i];

// (2) İç içe döngü -> O(n^2)
for (int i = 0; i < n; ++i)
    for (int j = 0; j < n; ++j)
        if (a[i] + a[j] == target) /* ... */;

// (3) Yarıya bölerek ilerleyen döngü -> O(log n)
for (int i = 1; i < n; i *= 2)
    process(i);   // i: 1, 2, 4, 8, ... n  =>  log2(n) adım
\end{lstlisting}

\gsub{Zaman Ölçme: Teoriyi Doğrulama}

Teorik analizi pratikte doğrulamak için \code{<chrono>} kütüphanesiyle
gerçek süreyi ölçebiliriz.

\begin{lstlisting}
#include <chrono>
#include <iostream>
using namespace std;

int main() {
    const int n = 100000;
    auto start = chrono::high_resolution_clock::now();

    long long sum = 0;
    for (int i = 0; i < n; ++i)
        for (int j = 0; j < n; ++j)   // O(n^2): 10^10 islem -> yavas!
            sum += 1;

    auto end = chrono::high_resolution_clock::now();
    auto ms = chrono::duration_cast<chrono::milliseconds>(end - start);
    cout << "Sure: " << ms.count() << " ms\n";
    return 0;
}
\end{lstlisting}

\gsub{Amortize (Amortized) Analiz}

Bazı işlemler tek tek pahalı görünür ama bir dizi işlem üzerinden ortalaması
ucuzdur. Klasik örnek \code{std::vector}'ün \code{push\_back}'idir: dolduğunda
kapasitesini iki katına çıkarıp tüm elemanları kopyalar (\bigO{n}), ama bu
nadiren olur. $n$ ekleme toplamda \bigO{n} sürer, dolayısıyla ekleme başına
\emph{amortize} maliyet \bigO{1}'dir.

\uyari{\bigO{} \emph{en kötü durumu} tarif eder ama her zaman gerçekçi
maliyeti yansıtmayabilir. Hash tablosunda arama en kötü \bigO{n}'dir (tüm
anahtarlar çakışırsa) ama \emph{ortalama} \bigO{1}'dir. Hangi durumdan
bahsettiğini her zaman belirt.}

% =====================================================================
\gsec{C++ ve STL: Algoritma Araç Kutusu}
% =====================================================================

Modern C++, veri yapıları ve algoritmalar için zengin bir standart kütüphane
(STL) sunar. Kendi veri yapılarını yazmayı öğrenmek eğitseldir, ancak gerçek
kodda çoğu zaman STL'in test edilmiş, optimize edilmiş bileşenlerini
kullanırsın. Bu bölüm, rehber boyunca kullanacağımız araçları tanıtır.

\gsub{STL Konteynerlerine Genel Bakış}

\begin{center}
\begin{tabularx}{\textwidth}{@{}>{\ttfamily\bfseries\color{anaMavi}}l >{\raggedright\arraybackslash}X c@{}}
\toprule
\rowcolor{acikMavi}\color{anaMavi}\textbf{Konteyner} & \color{anaMavi}\textbf{Kullanım} & \color{anaMavi}\textbf{Erişim}\\
\midrule
vector & Dinamik dizi; en çok kullanılan & \bigO{1} indeks\\
\rowcolor{acikGri} array & Sabit boyutlu dizi & \bigO{1} indeks\\
deque & Çift uçlu kuyruk & \bigO{1} uçlar\\
\rowcolor{acikGri} list & Çift yönlü bağlı liste & \bigO{n} arama\\
map / set & Sıralı (kırmızı-siyah ağaç) & \bigO{\log n}\\
\rowcolor{acikGri} unordered\_map / set & Hash tablosu & \bigO{1} ort.\\
stack / queue & Yığın / kuyruk (adaptör) & \bigO{1} uç\\
\rowcolor{acikGri} priority\_queue & İkili heap & \bigO{\log n}\\
\bottomrule
\end{tabularx}
\end{center}

\gsub{\code{std::vector}: En Sık Kullanılan Araç}

\begin{lstlisting}
#include <vector>
#include <iostream>
using namespace std;

int main() {
    vector<int> v = {5, 3, 8, 1};   // liste ilklendirme
    v.push_back(9);                 // sona ekle -> {5,3,8,1,9}
    v.pop_back();                   // sondan sil -> {5,3,8,1}

    cout << v[0] << " " << v.at(2) << "\n";  // 5 8  (.at sinir kontrolu yapar)
    cout << "Boyut: " << v.size() << "\n";

    // Aralik tabanli dongu
    for (int x : v) cout << x << " ";
    cout << "\n";

    // 2 boyutlu vector (n x m, sifirla dolu)
    int n = 3, m = 4;
    vector<vector<int>> grid(n, vector<int>(m, 0));
    grid[1][2] = 7;
    return 0;
}
\end{lstlisting}

\gsub{Yararlı STL Algoritmaları (\code{<algorithm>})}

\begin{lstlisting}
#include <algorithm>
#include <numeric>
#include <vector>
using namespace std;

vector<int> v = {5, 3, 8, 1, 3};

sort(v.begin(), v.end());                 // artan: {1,3,3,5,8}
sort(v.begin(), v.end(), greater<int>()); // azalan: {8,5,3,3,1}
reverse(v.begin(), v.end());              // ters cevir

int maxVal = *max_element(v.begin(), v.end());   // en buyuk
int minVal = *min_element(v.begin(), v.end());   // en kucuk
int sum = accumulate(v.begin(), v.end(), 0);  // toplam

// Ikili arama (dizi SIRALI olmali)
bool found = binary_search(v.begin(), v.end(), 5);
auto it = lower_bound(v.begin(), v.end(), 3); // >= 3 olan ilk konum

// Tekrarlari kaldir (once sirala): unique + erase idiomu
sort(v.begin(), v.end());
v.erase(unique(v.begin(), v.end()), v.end());
\end{lstlisting}

\ipucu{Rekabetçi programlama ve hızlı prototipleme için sık kullanılan
kısaltmalar: \code{using ll = long long;} (taşmayı önlemek için),
\code{\#define pb push\_back}, ve girdi/çıktı hızı için \code{main}
başında \code{ios::sync\_with\_stdio(false); cin.tie(nullptr);}.}

\gsub{Şablonlar (Templates): Jenerik Veri Yapıları}

Kendi veri yapılarımızı herhangi bir tip için çalışacak şekilde yazmak için
şablonları kullanırız. Rehber boyunca yapıları çoğunlukla \code{int} ile
göstereceğiz ama gerçekte şu kalıpla jenerikleştirilirler:

\begin{lstlisting}
template <typename T>
class Stack {
    vector<T> data;
public:
    void push(const T& x) { data.push_back(x); }
    void pop()            { data.pop_back(); }
    T& top()              { return data.back(); }
    bool empty() const    { return data.empty(); }
    size_t size() const   { return data.size(); }
};

// Kullanim: her tip icin ayni kod
Stack<int> intStack;      intStack.push(5);
Stack<string> strStack;   strStack.push("merhaba");
\end{lstlisting}

\bilgi{Bu rehberdeki tüm kod örnekleri C++17 ile derlenir. Önerilen komut:
\code{g++ -std=c++17 -Wall -Wextra -O2 program.cpp -o program}. Örnekleri
kendin derleyip çalıştırarak, değiştirerek denemek öğrenmenin en etkili
yoludur.}

\gsec{Temeller --- Alıştırmalar ve Çözümler}

Bu bölüm, KISIM 1'de işlenen algoritma analizi (Big-O), karmaşıklık muhakemesi ve C++ STL konularını pekiştirmek için hazırlanmış çözümlü sorular içerir. Her çözümde ilgili kod parçasının zaman ve gerektiğinde uzay karmaşıklığı açıkça belirtilmiştir.

\gsub{Karmaşıklık Analizi}

\soru{Aşağıdaki kod parçasının zaman karmaşıklığını Big-O cinsinden bulun ve nedenini açıklayın.
}
\begin{lstlisting}
int count = 0;
for (int i = 0; i < n; i++)
    for (int j = 0; j < n; j++)
        count++;
\end{lstlisting}

\cozum
Dıştaki döngü \code{i} için tam olarak $n$ kez, içteki döngü \code{j} için her \code{i} değerinde $n$ kez çalışır. Toplam yineleme sayısı $n \times n = n^2$ olur. \code{count++} işlemi sabit zamanlı (\bigO{1}) olduğundan toplam maliyet $n^2$ ile orantılıdır.

\begin{lstlisting}
// Toplam yineleme: n * n = n^2
// Zaman karmasikligi: O(n^2)
int countPairs(int n) {
    int count = 0;
    for (int i = 0; i < n; i++)   // n kez
        for (int j = 0; j < n; j++)  // her i icin n kez
            count++;              // O(1) islem
    return count;
}
\end{lstlisting}

\bilgi{İç içe iki bağımsız döngü, her ikisi de $n$'e kadar sayıyorsa çarpım kuralıyla \bigO{n^2} verir. Karmaşıklık $n$'e bağlı toplam yineleme sayısıyla belirlenir.}

\soru{Aşağıdaki döngünün zaman karmaşıklığı nedir? \code{i} değişkeninin nasıl değiştiğine dikkat edin.
}
\begin{lstlisting}
for (int i = n; i >= 1; i /= 2)
    doWork();  // O(1)
\end{lstlisting}

\cozum
\code{i} her adımda yarıya iniyor: $n, n/2, n/4, \dots, 1$. Kaç adımda 1'e ulaşırız? $n$'i kaç kez 2'ye böleriz sorusunun cevabı $\log_2 n$'dir. Dolayısıyla döngü yaklaşık $\log_2 n$ kez çalışır ve her adım \bigO{1} olduğundan toplam karmaşıklık \bigO{\log n} olur.

\begin{lstlisting}
// i: n -> n/2 -> n/4 -> ... -> 1
// Adim sayisi: log2(n)
// Zaman karmasikligi: O(log n)
int countSteps(int n) {
    int steps = 0;
    for (int i = n; i >= 1; i /= 2)
        steps++;   // O(1)
    return steps;  // yaklasik floor(log2(n)) + 1
}
\end{lstlisting}

\ipucu{Bir döngü değişkeni her adımda sabit bir çarpanla (2, 3, ...) çarpılıyor veya bölünüyorsa, döngü \bigO{\log n} çalışır. Sabit bir sayı eklenip çıkarılıyorsa \bigO{n} çalışır.}

\soru{İki ayrı (iç içe olmayan) döngüden oluşan aşağıdaki kodun karmaşıklığını bulun.
}
\begin{lstlisting}
for (int i = 0; i < n; i++)
    a();          // O(1)
for (int j = 0; j < n * n; j++)
    b();          // O(1)
\end{lstlisting}

\cozum
Birinci döngü $n$ kez, ikinci döngü $n^2$ kez çalışır. Ardışık (peş peşe) çalışan bloklarda maliyetler toplanır: $n + n^2$. Big-O gösteriminde en baskın terim alınır ve sabit katsayılar ihmal edilir; büyük $n$ için $n^2$ terimi $n$ terimini domine eder. Sonuç \bigO{n^2}'dir.

\begin{lstlisting}
// Birinci dongu: n islem
// Ikinci dongu:  n*n islem
// Toplam: n + n^2  =>  baskin terim n^2
// Zaman karmasikligi: O(n^2)
long long total(int n) {
    long long ops = 0;
    for (int i = 0; i < n; i++) ops++;        // n
    for (int j = 0; j < (long long)n * n; j++) ops++;  // n^2
    return ops;
}
\end{lstlisting}

\notk[Toplam kuralı]{Ardışık kod bloklarında karmaşıklıklar \emph{toplanır}; iç içe döngülerde \emph{çarpılır}. Toplamda yalnızca en hızlı büyüyen terim korunur: \bigO{n} + \bigO{n^2} = \bigO{n^2}.}

\soru{Bir dizide bir elemanı aramak için (a) doğrusal arama ve (b) önceden sıralanmış dizide ikili arama kullanılıyor. Her iki yaklaşımın karmaşıklığını karşılaştırın; hangisi ne zaman tercih edilmelidir?}

\cozum
Doğrusal arama en kötü durumda diziyi baştan sona tarar: \bigO{n}. İkili arama her adımda arama uzayını yarıya böldüğü için \bigO{\log n}'dir, ancak dizinin \emph{sıralı} olmasını gerektirir. Diziyi bir kez sıralamak \bigO{n \log n} maliyetlidir. Tek bir arama için doğrusal arama yeterlidir; ancak aynı dizide çok sayıda ($q$ adet) arama yapılacaksa, bir kez sıralayıp her sorguyu \bigO{\log n} yapmak (\bigO{n \log n + q \log n}) doğrusal aramanın \bigO{q \cdot n} maliyetinden çok daha iyidir.

\begin{lstlisting}
#include <algorithm>
#include <vector>
using namespace std;

// Dogrusal arama: O(n), sirali olma sarti yok
int linearSearch(const vector<int>& v, int target) {
    for (int i = 0; i < (int)v.size(); i++)
        if (v[i] == target) return i;  // en kotu durum: n kar.
    return -1;
}

// Ikili arama: O(log n), dizi SIRALI olmali
bool binarySearch(const vector<int>& sorted, int target) {
    int lo = 0, hi = (int)sorted.size() - 1;
    while (lo <= hi) {
        int mid = lo + (hi - lo) / 2;  // tasmayi onler
        if (sorted[mid] == target) return true;
        if (sorted[mid] < target) lo = mid + 1; // sag yari
        else hi = mid - 1;                       // sol yari
    }
    return false;
}
\end{lstlisting}

\uyari{İkili arama yalnızca sıralı dizide doğru çalışır. Sıralama maliyetini de hesaba katın: yalnızca birkaç arama yapacaksanız sıralama zahmetine değmez.}

\soru{\code{std::vector::push\_back} işleminin \emph{amortize} \bigO{1} olduğu söylenir, oysa bazı çağrılar tüm elemanları yeni belleğe kopyalar. Bu nasıl mümkün olur? $n$ eleman eklemenin toplam maliyetini açıklayın.}

\cozum
\code{vector} dolduğunda kapasitesini genellikle iki katına çıkarır ve mevcut elemanları yeni belleğe kopyalar. Bu kopyalama tek bir \code{push\_back} için \bigO{n} olabilir. Ancak yeniden boyutlandırma seyrek gerçekleşir: kapasite $1, 2, 4, 8, \dots, n$ şeklinde büyür. $n$ ekleme boyunca toplam kopyalama sayısı $1 + 2 + 4 + \dots + n < 2n$ olur; bu geometrik seri $\bigO{n}$ toplam maliyet demektir. Toplam maliyeti işlem sayısına (\bigO{n}/n) bölünce her işlem başına \emph{amortize} \bigO{1} elde edilir. Bu, "her tek çağrı \bigO{1}" değil, "$n$ çağrının ortalaması \bigO{1}" anlamına gelir.

\begin{lstlisting}
#include <vector>
#include <iostream>
using namespace std;

int main() {
    vector<int> v;
    // n push_back cagrisinin TOPLAM maliyeti O(n)
    // => cagri basina amortize O(1)
    for (int i = 0; i < 16; i++) {
        v.push_back(i);
        // Kapasite dolunca ikiye katlanir; kopyalama seyrek olur.
        cout << "size=" << v.size()
             << " capacity=" << v.capacity() << "\n";
    }
    return 0;
}
\end{lstlisting}

\bilgi{Amortize analiz, pahalı ama seyrek işlemlerin maliyetini ucuz sık işlemlere yayar. \code{reserve(n)} ile kapasiteyi önceden ayırırsanız yeniden tahsisleri tamamen ortadan kaldırıp performansı artırabilirsiniz.}

\soru{Aşağıdaki iki fonksiyonun \emph{uzay} (bellek) karmaşıklığını karşılaştırın. Çıktı bir dizi döndürülmesi için ayrılan bellek dahil edilmelidir.
}
\begin{lstlisting}
long long sumA(int n) {                 // A
    long long s = 0;
    for (int i = 1; i <= n; i++) s += i;
    return s;
}
std::vector<int> sumB(int n) {          // B
    std::vector<int> v;
    for (int i = 1; i <= n; i++) v.push_back(i);
    return v;
}
\end{lstlisting}

\cozum
\code{sumA} yalnızca sabit sayıda değişken (\code{s}, \code{i}) kullanır; girdi büyüklüğünden bağımsız sabit bellek harcar: \bigO{1} uzay. \code{sumB} ise $n$ elemanlı bir \code{vector} oluşturur ve döndürür; kullandığı bellek $n$ ile doğrusal artar: \bigO{n} uzay. Her ikisinin de zaman karmaşıklığı \bigO{n}'dir, ancak uzay maliyetleri farklıdır. Genel kural: sonuç için gereken ek belleği (yardımcı yapılar, döndürülen koleksiyon) sayarız; girdinin kendisini çoğu zaman saymayız.

\begin{lstlisting}
// sumA: yardimci bellek sabit -> uzay O(1), zaman O(n)
// sumB: n elemanli vektor    -> uzay O(n), zaman O(n)
\end{lstlisting}

\ipucu{Uzay karmaşıklığında yalnızca \emph{ek} (auxiliary) bellek sayılır: yeni diziler, rekürsiyon yığını, hash tabloları. Sabit sayıda skaler değişken her zaman \bigO{1}'dir.}

\gsub{Big-O Muhakemesi}

\soru{Bir yarışma probleminde $n \le 10^6$ veriliyor ve zaman sınırı 1 saniye. \bigO{n^2} bir çözüm kabul edilir mi? \bigO{n \log n} kabul edilir mi? Muhakemenizi açıklayın.}

\cozum
Kaba bir kural olarak modern bir işlemci 1 saniyede yaklaşık $10^8$ ile $10^9$ temel işlem yapar. $n = 10^6$ için:
\begin{itemize}
\item \bigO{n^2} $= (10^6)^2 = 10^{12}$ işlem. Bu $10^8$--$10^9$ bütçesinin çok üstündedir; saniyeler-dakikalar sürer, \uyari{kesinlikle zaman aşımına uğrar.}
\item \bigO{n \log n} $\approx 10^6 \times 20 = 2 \times 10^7$ işlem. Bu bütçenin rahatça altındadır; \emph{kabul edilir.}
\end{itemize}
Dolayısıyla $n = 10^6$ ölçeğinde \bigO{n^2} çözümler pratik değildir; \bigO{n \log n} veya \bigO{n} hedeflenmelidir.

\begin{lstlisting}
// n = 10^6 icin islem tahminleri:
//   O(n)        ~ 1e6      -> cok hizli
//   O(n log n)  ~ 2e7      -> uygun (< ~1 sn)
//   O(n^2)      ~ 1e12     -> COK yavas (zaman asimi)
// Kural: 1 sn'de ~1e8-1e9 basit islem varsayilir.
\end{lstlisting}

\notk[Ölçeğe göre hedef karmaşıklık]{$n \le 10^6$--$10^7$ için \bigO{n} veya \bigO{n \log n}; $n \le 10^5$ için \bigO{n \log n}; $n \le 10^3$--$10^4$ için \bigO{n^2} genellikle uygundur; $n \le 20$--$25$ için \bigO{2^n} gibi üstel çözümler bile kabul edilebilir.}

\soru{Aşağıdaki üç problem boyutundan her biri için uygun olabilecek \emph{en yavaş} (ama yeterince hızlı) karmaşıklık sınıfını söyleyin ve neden onu seçtiğinizi gerekçelendirin: (a) $n \le 20$, (b) $n \le 10^3$, (c) $n \le 10^6$.}

\cozum
\begin{tabularx}{\textwidth}{@{}lX@{}}
\textbf{Sınır} & \textbf{Uygun en yavaş karmaşıklık ve gerekçe} \\
$n \le 20$ & \bigO{2^n} veya \bigO{n!} kabul edilebilir. $2^{20} \approx 10^6$ olduğundan tüm alt kümeleri/permütasyonları denemek (brute force, bitmask) mümkündür. \\
$n \le 10^3$ & \bigO{n^2} uygundur: $(10^3)^2 = 10^6$ işlem. İç içe iki döngü rahatça çalışır; hatta $n^2 \log n$ da sınırda geçebilir. \\
$n \le 10^6$ & \bigO{n} veya \bigO{n \log n} gerekir: $10^6$ veya $\sim 2\times 10^7$ işlem. \bigO{n^2}$=10^{12}$ olur ve zaman aşımına uğrar. \\
\end{tabularx}

\begin{lstlisting}
// n <= 20    : O(2^n) ~ 1e6  (bitmask brute force uygun)
// n <= 1000  : O(n^2) ~ 1e6  (ic ice donguler uygun)
// n <= 1e6   : O(n) / O(n log n)  (n^2 = 1e12 -> yasak)
\end{lstlisting}

\bilgi{$n$'in üst sınırı, hedeflenmesi gereken karmaşıklık için güçlü bir ipucudur. Küçük $n$ (örneğin $\le 25$) çoğu zaman üstel bir çözümün beklendiğine işaret eder.}

\soru{Bir öğrenci "\bigO{n \log n} ile \bigO{n^2} arasındaki fark sadece teorik, pratikte önemsiz" diyor. $n = 10^6$ için bu iki karmaşıklığın işlem sayısını karşılaştırarak bu iddiayı değerlendirin.}

\cozum
$n = 10^6$ için:
\begin{itemize}
\item \bigO{n \log n} $\approx 10^6 \times \log_2(10^6) \approx 10^6 \times 20 = 2 \times 10^7$.
\item \bigO{n^2} $= (10^6)^2 = 10^{12}$.
\end{itemize}
Oran $10^{12} / (2 \times 10^7) = 5 \times 10^4$; yani \bigO{n^2} çözüm yaklaşık \emph{50.000 kat} daha fazla işlem yapar. \bigO{n \log n} çözüm milisaniyeler sürerken \bigO{n^2} çözüm on dakikalar sürebilir. İddia yanlıştır: girdi büyüdükçe iki sınıf arasındaki fark astronomik boyutlara ulaşır. Fark, sabit katsayıların gizlediği küçük $n$ değerlerinde önemsiz görünse de büyük $n$'de belirleyicidir.

\begin{lstlisting}
#include <cmath>
#include <iostream>
using namespace std;

int main() {
    double n = 1e6;
    double nlogn = n * log2(n);   // ~2e7
    double nsq   = n * n;         // 1e12
    cout << "n log n = " << nlogn << "\n";
    cout << "n^2     = " << nsq   << "\n";
    cout << "oran    = " << nsq / nlogn << "\n"; // ~5e4
    return 0;
}
\end{lstlisting}

\uyari{Big-O sabit katsayıları gizler; bu yüzden küçük $n$'de \bigO{n^2} bir algoritma \bigO{n \log n}'den hızlı olabilir. Fark, girdi büyüdükçe kesinleşir; asimptotik davranış büyük ölçekte kazanır.}

\gsub{STL ile Problem Çözme}

\soru{Bir \code{struct Student \{ string name; int score; \}} vektörünü \emph{puana göre azalan}, eşit puanlarda \emph{isme göre artan} sıralayın. Bunu bir lambda karşılaştırıcı ile yapın ve karmaşıklığı belirtin.}

\cozum
\code{std::sort} karşılaştırmalı sıralama yapar ve \bigO{n \log n} zaman karmaşıklığına sahiptir. Özel sıralama ölçütü için üçüncü argüman olarak bir lambda comparator veririz. Comparator, "birinci eleman ikinciden \emph{önce} gelmeli mi?" sorusuna \code{true}/\code{false} döndürmelidir. Azalan puan için \code{a.score > b.score}, eşitlikte artan isim için \code{a.name < b.name} kullanırız.

\begin{lstlisting}
#include <algorithm>
#include <string>
#include <vector>
#include <iostream>
using namespace std;

struct Student { string name; int score; };

int main() {
    vector<Student> students = {
        {"Ali", 90}, {"Zeynep", 90}, {"Mehmet", 75}
    };
    // O(n log n): once puana gore azalan, esitlikte isme gore artan
    sort(students.begin(), students.end(),
         [](const Student& a, const Student& b) {
             if (a.score != b.score)
                 return a.score > b.score;  // azalan puan
             return a.name < b.name;        // artan isim
         });
    for (const auto& s : students)
        cout << s.name << " " << s.score << "\n";
    return 0;
}
\end{lstlisting}

\bilgi{Comparator "kesin küçüktür" (strict weak ordering) tanımlamalıdır: eşit elemanlar için \code{false} döndürmelidir. \code{<=} kullanmak tanımsız davranışa yol açar.}

\soru{Sıralı bir tam sayı dizisinde, belirli bir $[L, R]$ değer aralığına düşen eleman sayısını \code{lower\_bound} ve \code{upper\_bound} kullanarak \bigO{\log n} zamanda bulun.}

\cozum
Sıralı bir dizide \code{lower\_bound(L)}, değeri $\ge L$ olan ilk elemana; \code{upper\_bound(R)}, değeri $> R$ olan ilk elemana işaret eder. İki iteratör arasındaki mesafe (\code{upper - lower}), $[L, R]$ aralığındaki eleman sayısını verir. Her iki fonksiyon da ikili arama yaptığı için \bigO{\log n} çalışır; dizinin sıralı olması şarttır.

\begin{lstlisting}
#include <algorithm>
#include <vector>
#include <iostream>
using namespace std;

int main() {
    vector<int> v = {1, 3, 3, 5, 8, 8, 8, 10};  // SIRALI olmali
    int L = 3, R = 8;
    // ilk >= L   ve  ilk > R  konumlari (her biri O(log n))
    auto lo = lower_bound(v.begin(), v.end(), L);
    auto hi = upper_bound(v.begin(), v.end(), R);
    long long countInRange = hi - lo;  // [L, R] icindeki eleman sayisi
    cout << countInRange << "\n";      // 6 (3,3,5,8,8,8)
    return 0;
}
\end{lstlisting}

\ipucu{\code{lower\_bound} $\ge$ değeri, \code{upper\_bound} $>$ değeri arar. Bir değerin dizideki kopya sayısı \code{upper\_bound(x) - lower\_bound(x)} ile \bigO{\log n} bulunur.}

\soru{Bir metindeki (kelime vektörü) her kelimenin kaç kez geçtiğini sayın ve en çok geçen kelimeyi bulun. \code{unordered\_map} kullanarak ortalama \bigO{n} bir çözüm yazın.}

\cozum
Frekans sayımı için hash tablosu olan \code{unordered\_map<string,int>} idealdir: ekleme ve erişim ortalama \bigO{1}'dir, dolayısıyla $n$ kelimeyi saymak ortalama \bigO{n} olur. \code{freq[word]++} ifadesi kelime yoksa 0'dan başlatıp artırır (\code{operator[]} varsayılan değer ekler). Sonra map üzerinde tek geçişle (\bigO{k}, $k$ = tekil kelime sayısı) en yüksek frekanslıyı buluruz. Not: \code{unordered\_map} sırasız ve ortalama \bigO{1}; sıralı gezinme gerekiyorsa \code{map} (\bigO{\log n} işlem, kırmızı-siyah ağaç) tercih edilir.

\begin{lstlisting}
#include <unordered_map>
#include <string>
#include <vector>
#include <iostream>
using namespace std;

int main() {
    vector<string> words = {"elma","armut","elma","kiraz","elma","armut"};
    unordered_map<string,int> freq;   // ortalama O(1) ekleme/erisim
    for (const string& w : words)
        freq[w]++;                    // yoksa 0'dan baslar -> toplam O(n)

    string best; int bestCount = 0;
    for (const auto& [word, cnt] : freq)  // O(k), k = tekil sayisi
        if (cnt > bestCount) { bestCount = cnt; best = word; }

    cout << best << " -> " << bestCount << "\n"; // elma -> 3
    return 0;
}
\end{lstlisting}

\uyari{\code{unordered\_map} ortalama \bigO{1}, ancak kötü hash çakışmalarında en kötü durum \bigO{n} olabilir. Deterministik \bigO{\log n} garanti veya sıralı iterasyon gerekiyorsa \code{std::map} kullanın.}

\soru{Bir tam sayı vektöründeki \emph{benzersiz} (tekil) elemanları bulup küçükten büyüğe sıralı yazdırın. Bunu \code{std::set} ile yapın ve karmaşıklığı belirtin.}

\cozum
\code{std::set} elemanları hem tekilleştirir hem de sıralı tutar (dengeli ikili arama ağacı). $n$ eleman eklemek, ekleme başına \bigO{\log n} olduğundan toplam \bigO{n \log n} maliyetlidir. Set üzerinde gezinmek elemanları otomatik olarak artan sırada verir; ek bir sıralamaya gerek kalmaz. Yalnızca tekilleştirme yeterliyse ve sıra gerekmiyorsa \code{unordered\_set} ortalama \bigO{n} ile daha hızlıdır.

\begin{lstlisting}
#include <set>
#include <vector>
#include <iostream>
using namespace std;

int main() {
    vector<int> v = {5, 3, 5, 1, 3, 8, 1};
    set<int> unique(v.begin(), v.end());  // tekillestir + sirala: O(n log n)
    for (int x : unique)                  // artan sirada gezinir
        cout << x << " ";                 // 1 3 5 8
    cout << "\n";
    return 0;
}
\end{lstlisting}

\ipucu{Sıra önemliyse \code{set} (\bigO{n \log n}), önemsizse \code{unordered\_set} (ortalama \bigO{n}). Vektörde tekilleştirmenin bir başka yolu: \code{sort} + \code{unique} + \code{erase} deyimi (yine \bigO{n \log n}).}

\soru{Bir sayı akışında (vektör) \emph{en büyük $k$} elemanı verimli biçimde bulun. Tüm diziyi sıralamak yerine \code{std::priority\_queue} tabanlı \bigO{n \log k} bir çözüm yazın.}

\cozum
Tüm diziyi sıralamak \bigO{n \log n}'dir. Yalnızca en büyük $k$ eleman gerekiyorsa, boyutu $k$ ile sınırlı bir \emph{min-heap} (en küçüğü tepede tutan öncelik kuyruğu) kullanırız. Her elemanı ekleriz; heap boyutu $k$'yı aşınca tepedeki en küçüğü çıkarırız. Böylece heapte her zaman o ana kadarki en büyük $k$ eleman kalır. Her işlem \bigO{\log k}, toplam $n$ eleman için \bigO{n \log k}'dir; $k \ll n$ olduğunda tam sıralamadan hızlıdır ve daha az bellek kullanır (\bigO{k}).

\begin{lstlisting}
#include <queue>
#include <vector>
#include <iostream>
using namespace std;

int main() {
    vector<int> v = {4, 1, 7, 3, 9, 2, 8};
    int k = 3;
    // min-heap: en kucuk tepede; boyutu k ile sinirli tutulur
    priority_queue<int, vector<int>, greater<int>> minHeap;
    for (int x : v) {
        minHeap.push(x);                 // O(log k)
        if ((int)minHeap.size() > k)
            minHeap.pop();               // en kucugu at -> en buyuk k kalir
    }
    // Toplam: O(n log k), bellek: O(k)
    while (!minHeap.empty()) {           // en buyuk k eleman (artan sirada)
        cout << minHeap.top() << " ";    // 7 8 9
        minHeap.pop();
    }
    cout << "\n";
    return 0;
}
\end{lstlisting}

\bilgi{En büyük $k$ için \emph{min-heap}, en küçük $k$ için \emph{max-heap} kullanılır (heapten hangi ucun atılacağı ters çevrilir). \code{priority\_queue} varsayılan olarak max-heap'tir; \code{greater<int>} ile min-heap'e dönüşür.}

\soru{Bir tam sayı vektörünün elemanlarının toplamını ve en büyük elemanını, döngü yazmadan STL algoritmaları \code{std::accumulate} ve \code{std::max\_element} ile hesaplayın. Karmaşıklıklarını belirtin.}

\cozum
\code{std::accumulate} (başlık \code{<numeric>}) bir aralığın elemanlarını verilen başlangıç değeriyle toplar; tek geçişte çalıştığı için \bigO{n}'dir. Başlangıç değerinin türü sonuç türünü belirler: taşmayı önlemek için \code{0LL} (long long) vermek iyi bir alışkanlıktır. \code{std::max\_element} (başlık \code{<algorithm>}) en büyük elemana bir iteratör döndürür, o da tek geçiş \bigO{n}'dir; değere \code{*it} ile erişilir. Her iki algoritma da açık indeksli döngüden hem daha okunaklı hem hataya daha az açıktır.

\begin{lstlisting}
#include <numeric>     // accumulate
#include <algorithm>   // max_element
#include <vector>
#include <iostream>
using namespace std;

int main() {
    vector<int> v = {4, 1, 7, 3, 9, 2};
    // Toplam: O(n). Baslangic 0LL -> long long, tasmayi onler.
    long long sum = accumulate(v.begin(), v.end(), 0LL);
    // En buyuk elemana iterator: O(n)
    auto it = max_element(v.begin(), v.end());
    cout << "toplam = " << sum << "\n";   // 26
    cout << "maks   = " << *it << "\n";   // 9
    return 0;
}
\end{lstlisting}

\ipucu{\code{accumulate} dördüncü argüman olarak özel bir ikili işlem alabilir (örn. çarpım için \code{multiplies<>()}). \code{max\_element} de bir comparator lambda ile özel ölçütte maksimum bulabilir. İkisi de \bigO{n}'dir.}

% #####################################################################
\gpart{2}{Temel Veri Yapıları}
% #####################################################################

% =====================================================================
\gsec{Diziler ve Dinamik Diziler}
% =====================================================================

Dizi (array), aynı tipteki elemanların bellekte \emph{ardışık} olarak
tutulduğu en temel veri yapısıdır. Ardışıklık sayesinde herhangi bir elemana
indeksle \bigO{1}'de erişilir: eleman adresi = \code{taban + indeks $\times$
eleman\_boyutu}. Bu basit formül, dizilerin gücünün ve sınırlarının kaynağıdır.

\gsub{Statik Diziler ve Bellek Düzeni}

\begin{lstlisting}
int a[5] = {10, 20, 30, 40, 50};
// Bellek:  [10][20][30][40][50]
//  indeks:   0   1   2   3   4
// a[3] = *(a + 3) = *(taban + 3*sizeof(int))  ->  O(1) erisim

std::array<int, 5> b = {1, 2, 3, 4, 5};  // C++ tarzi, boyut tip bilgisinde
cout << b.size() << "\n";                 // 5
\end{lstlisting}

\gsub{Dinamik Dizi Nasıl Çalışır?}

Statik dizinin boyutu derleme zamanında sabittir. \code{std::vector} ise
büyüyebilen bir dizidir. İçeride bir bloğu ayırır; dolduğunda \emph{daha büyük}
(genelde iki kat) yeni bir blok ayırır, eski elemanları kopyalar ve eskiyi
serbest bırakır. Bunu kendimiz yazarak mekanizmayı görelim:

\begin{lstlisting}
class DynamicArray {
    int* data;
    int  capacity;
    int  size;

    void grow() {                     // kapasiteyi iki katina cikar
        capacity = (capacity == 0) ? 1 : capacity * 2;
        int* newData = new int[capacity];
        for (int i = 0; i < size; ++i)
            newData[i] = data[i];     // O(n) kopyalama
        delete[] data;
        data = newData;
    }
public:
    DynamicArray() : data(nullptr), capacity(0), size(0) {}
    ~DynamicArray() { delete[] data; }

    void add(int x) {                 // amortize O(1)
        if (size == capacity) grow();
        data[size++] = x;
    }
    int  get(int i) const { return data[i]; }   // O(1)
    int  getSize() const { return size; }
};
\end{lstlisting}

\notk[Neden iki kat?]{Kapasiteyi her seferinde \emph{sabit miktarda}
(örneğin +1) artırsaydık, $n$ ekleme $1+2+\dots+n = O(n^2)$ kopyalamaya
mal olurdu. İki katına çıkarınca toplam kopyalama $1+2+4+\dots+n < 2n$
olur; yani $n$ ekleme \bigO{n}, ekleme başına \emph{amortize} \bigO{1}.}

\gsub{Dizi İşlemlerinin Karmaşıklığı}

\begin{center}
\begin{tabularx}{\textwidth}{@{}>{\bfseries\color{anaMavi}}l c X@{}}
\toprule
\rowcolor{acikMavi}\color{anaMavi}\textbf{İşlem} & \color{anaMavi}\textbf{Karmaşıklık} & \color{anaMavi}\textbf{Açıklama}\\
\midrule
İndeksle erişim & \bigO{1} & Adres hesabı doğrudan.\\
\rowcolor{acikGri} Sona ekleme & \bigO{1}$^*$ & Amortize; yer varsa anında.\\
Başa/ortaya ekleme & \bigO{n} & Sonraki elemanlar kaydırılır.\\
\rowcolor{acikGri} Silme (orta) & \bigO{n} & Boşluk kaydırılarak doldurulur.\\
Arama (sırasız) & \bigO{n} & Tek tek bakılır.\\
\rowcolor{acikGri} Arama (sıralı) & \bigO{\log n} & İkili arama.\\
\bottomrule
\end{tabularx}
\end{center}

\gsub{Sık Kullanılan Dizi Teknikleri}

\gsubsub{İki İşaretçi (Two Pointers)}

Sıralı bir dizide, iki uçtan ortaya doğru ilerleyen işaretçilerle birçok
problem \bigO{n}'de çözülür. Örnek: sıralı dizide toplamı \code{target} olan
iki eleman var mı?

\begin{lstlisting}
bool hasTwoSum(const vector<int>& a, int target) {
    int left = 0, right = a.size() - 1;
    while (left < right) {
        int sum = a[left] + a[right];
        if (sum == target) return true;
        if (sum < target)  ++left;    // daha buyuk toplam gerek
        else               --right;   // daha kucuk toplam gerek
    }
    return false;
}
\end{lstlisting}

\gsubsub{Kayan Pencere (Sliding Window)}

Ardışık bir alt dizinin özelliğini (toplam, en büyük vb.) ararken, pencereyi
kaydırarak her adımda baştan hesaplamak yerine güncelleriz. Örnek: boyutu
\code{k} olan en büyük pencere toplamı.

\begin{lstlisting}
int maxWindowSum(const vector<int>& a, int k) {
    int window = 0;
    for (int i = 0; i < k; ++i) window += a[i];  // ilk pencere
    int maxSum = window;
    for (int i = k; i < (int)a.size(); ++i) {
        window += a[i] - a[i - k];   // sag ekle, sol cikar -> O(1)
        maxSum = max(maxSum, window);
    }
    return maxSum;                     // toplam O(n)
}
\end{lstlisting}

\gsubsub{Önek Toplamı (Prefix Sum)}

Bir dizide birçok kez ``$[l, r]$ aralığının toplamı nedir?'' sorusu
soruluyorsa, önce önek toplamı tablosu kurarız; sonra her sorgu \bigO{1}'dir.

\begin{lstlisting}
vector<long long> prefix(a.size() + 1, 0);
for (int i = 0; i < (int)a.size(); ++i)
    prefix[i + 1] = prefix[i] + a[i];
// [l, r] araligi toplami (dahil):
long long rangeSum(int l, int r) { return prefix[r + 1] - prefix[l]; }
\end{lstlisting}

% =====================================================================
\gsec{Bağlı Listeler}
% =====================================================================

Bağlı liste (linked list), her elemanı bir \emph{düğüm} (node) olarak tutan ve
düğümleri işaretçilerle birbirine bağlayan bir yapıdır. Diziden farkı,
elemanların bellekte ardışık \emph{olmamasıdır}; her düğüm bir sonrakinin
adresini saklar. Bu, ekleme/silmeyi ucuzlatır ama indeksle erişimi kaybettirir.

\begin{center}
\begin{tabularx}{\textwidth}{@{}>{\bfseries\color{anaMavi}}l X X@{}}
\toprule
\rowcolor{acikMavi}\color{anaMavi}\textbf{Özellik} & \color{anaMavi}\textbf{Dizi} & \color{anaMavi}\textbf{Bağlı Liste}\\
\midrule
İndeksle erişim & \bigO{1} & \bigO{n}\\
\rowcolor{acikGri} Başa ekleme & \bigO{n} & \bigO{1}\\
Bellek düzeni & Ardışık & Dağınık + işaretçi maliyeti\\
\rowcolor{acikGri} Önbellek dostu mu? & Evet & Hayır\\
\bottomrule
\end{tabularx}
\end{center}

\gsub{Tek Yönlü Bağlı Liste}

Her düğüm bir değer ve \emph{sadece} bir sonraki düğüme işaretçi tutar. Son
düğümün \code{next}'i \code{nullptr}'dır.

\begin{lstlisting}
struct Node {
    int   data;
    Node* next;
    Node(int v) : data(v), next(nullptr) {}
};

class LinkedList {
    Node* head = nullptr;
public:
    // Basa ekle -> O(1)
    void pushFront(int v) {
        Node* d = new Node(v);
        d->next = head;
        head = d;
    }
    // Sona ekle -> O(n) (son dugumu bulmak gerekir)
    void pushBack(int v) {
        Node* d = new Node(v);
        if (!head) { head = d; return; }
        Node* p = head;
        while (p->next) p = p->next;
        p->next = d;
    }
    // Bir degeri sil -> O(n)
    void remove(int v) {
        Node** p = &head;                 // isaretciye isaretci: sik idiom
        while (*p && (*p)->data != v) p = &(*p)->next;
        if (*p) { Node* toDelete = *p; *p = (*p)->next; delete toDelete; }
    }
    void print() const {
        for (Node* p = head; p; p = p->next) cout << p->data << " -> ";
        cout << "nullptr\n";
    }
    ~LinkedList() {
        while (head) { Node* s = head; head = head->next; delete s; }
    }
};
\end{lstlisting}

\gsub{Listeyi Ters Çevirme (Klasik Mülakat Sorusu)}

Üç işaretçiyle listeyi tek geçişte, \bigO{n} zaman ve \bigO{1} ek bellekte
ters çeviririz. Her adımda mevcut düğümün bağını geriye döndürürüz.

\begin{lstlisting}
Node* reverse(Node* head) {
    Node* prev = nullptr;
    Node* curr = head;
    while (curr) {
        Node* next = curr->next;         // ilerisini sakla
        curr->next = prev;               // bagi geriye cevir
        prev = curr;                     // ikisini de bir adim ilerlet
        curr = next;
    }
    return prev;     // yeni bas
}
\end{lstlisting}

\gsub{Çift Yönlü Bağlı Liste ve Döngü Tespiti}

Çift yönlü listede her düğüm hem önceki hem sonraki düğümü tutar; bu, geriye
gitmeyi ve bir düğümü \bigO{1}'de silmeyi mümkün kılar. Bir listede döngü olup
olmadığını, \textbf{Floyd'un tavşan-kaplumbağa} algoritmasıyla (biri iki, biri
tek adım atan iki işaretçi) \bigO{n} zaman, \bigO{1} bellekte buluruz.

\begin{lstlisting}
bool hasCycle(Node* head) {
    Node* slow = head;
    Node* fast = head;
    while (fast && fast->next) {
        slow = slow->next;               // 1 adim
        fast = fast->next->next;         // 2 adim
        if (slow == fast) return true;   // karsilastilar -> dongu var
    }
    return false;
}
\end{lstlisting}

\uyari{Ham işaretçilerle (\code{new}/\code{delete}) bağlı liste yazmak
eğitseldir ama gerçek kodda bellek sızıntısı ve çift-serbestleme (double-free)
hatalarına açıktır. Üretimde \code{std::list} / \code{std::forward\_list}
veya \code{std::unique\_ptr} tabanlı düğümler tercih edilir.}

% =====================================================================
\gsec{Yığın (Stack)}
% =====================================================================

Yığın, \textbf{LIFO} (Last In, First Out --- son giren ilk çıkar) ilkesiyle
çalışan bir yapıdır. Sadece bir uçtan (tepe) ekleme ve çıkarma yapılır; tıpkı
üst üste dizilmiş tabaklar gibi. Fonksiyon çağrı yığını, geri-al (undo)
işlemleri, ifade değerlendirme ve derinlik-öncelikli arama yığınla çalışır.

\gsub{Diziyle Yığın Uygulaması}

\begin{lstlisting}
template <typename T>
class Stack {
    vector<T> v;
public:
    void   push(const T& x) { v.push_back(x); }   // O(1)
    void   pop()            { v.pop_back();   }    // O(1)
    T&     top()            { return v.back(); }   // O(1)
    bool   empty() const    { return v.empty(); }
    size_t size() const     { return v.size();  }
};

// STL karsiligi:
#include <stack>
stack<int> s;
s.push(1); s.push(2);
cout << s.top();   // 2
s.pop();
\end{lstlisting}

\gsub{Uygulama: Parantez Dengeleme}

Bir ifadedeki parantezlerin \code{()\{\}[]} doğru eşleşip eşleşmediğini
yığınla kontrol ederiz: açılışı yığına at, kapanışta tepedeki eşleşmeli.

\begin{lstlisting}
bool isBalanced(const string& s) {
    stack<char> y;
    for (char c : s) {
        if (c == '(' || c == '{' || c == '[') y.push(c);
        else if (c == ')' || c == '}' || c == ']') {
            if (y.empty()) return false;
            char a = y.top(); y.pop();
            if ((c == ')' && a != '(') ||
                (c == '}' && a != '{') ||
                (c == ']' && a != '[')) return false;
        }
    }
    return y.empty();   // acik kalan olmamali
}
\end{lstlisting}

\gsub{Uygulama: Sonraki Büyük Eleman (Monoton Yığın)}

Bir dizide her eleman için ``sağındaki ilk daha büyük eleman'' sorusunu,
azalan bir \emph{monoton yığın} ile toplam \bigO{n}'de yanıtlarız.

\begin{lstlisting}
vector<int> nextGreater(const vector<int>& a) {
    int n = a.size();
    vector<int> result(n, -1);
    stack<int> y;                        // indeks tutar (azalan degerler)
    for (int i = 0; i < n; ++i) {
        while (!y.empty() && a[y.top()] < a[i]) {
            result[y.top()] = a[i];      // a[i], yigindaki daha kucukleri asar
            y.pop();
        }
        y.push(i);
    }
    return result;
}
\end{lstlisting}

% =====================================================================
\gsec{Kuyruk ve Çift Uçlu Kuyruk (Deque)}
% =====================================================================

Kuyruk, \textbf{FIFO} (First In, First Out --- ilk giren ilk çıkar) ilkesiyle
çalışır: bir uçtan (arka) eklenir, diğer uçtan (ön) çıkarılır; market sırası
gibi. Genişlik-öncelikli arama, iş zamanlama ve tampon (buffer) yapıları
kuyruk kullanır.

\gsub{Dairesel Kuyruk (Circular Buffer)}

Sabit boyutlu bir diziyi ``dairesel'' kullanarak, ön ve arka indeksleri modulo
ile döndürerek \bigO{1} ekleme/çıkarma elde ederiz. Böylece eleman kaydırmaya
gerek kalmaz.

\begin{lstlisting}
class CircularQueue {
    vector<int> v;
    int front = 0, rear = 0, count = 0, cap;
public:
    CircularQueue(int k) : v(k), cap(k) {}
    bool isFull()  const { return count == cap; }
    bool isEmpty() const { return count == 0;   }

    bool enqueue(int x) {                // arkaya ekle -> O(1)
        if (isFull()) return false;
        v[rear] = x;
        rear = (rear + 1) % cap;         // dairesel ilerleme
        ++count;
        return true;
    }
    bool dequeue(int& x) {               // onden cikar -> O(1)
        if (isEmpty()) return false;
        x = v[front];
        front = (front + 1) % cap;
        --count;
        return true;
    }
};

// STL karsiligi:
#include <queue>
queue<int> q;
q.push(1); q.push(2);
cout << q.front();   // 1
q.pop();
\end{lstlisting}

\gsub{Deque ve Kayan Pencerede Maksimum}

\code{std::deque} (double-ended queue) her iki uçtan da \bigO{1} ekleme/çıkarma
sunar. Güçlü bir uygulaması, boyutu \code{k} olan her pencerenin maksimumunu
toplam \bigO{n}'de bulan \textbf{monoton deque} tekniğidir.

\begin{lstlisting}
#include <deque>
vector<int> windowMax(const vector<int>& a, int k) {
    deque<int> dq;                       // azalan degerli indeksler
    vector<int> result;
    for (int i = 0; i < (int)a.size(); ++i) {
        if (!dq.empty() && dq.front() == i - k) dq.pop_front(); // pencereden cikan
        while (!dq.empty() && a[dq.back()] <= a[i]) dq.pop_back();
        dq.push_back(i);
        if (i >= k - 1) result.push_back(a[dq.front()]);  // on = maksimum
    }
    return result;
}
\end{lstlisting}

% =====================================================================
\gsec{Hash Tabloları}
% =====================================================================

Hash tablosu, anahtar-değer çiftlerini \emph{ortalama} \bigO{1} zamanda
saklayan ve arayan bir yapıdır. Fikir: bir \emph{hash fonksiyonu} anahtarı bir
dizi indeksine çevirir; böylece anahtarı ``hesaplayarak'' doğrudan yerine
gideriz. C++'ta \code{std::unordered\_map} ve \code{std::unordered\_set} bunu
sağlar.

\gsub{Nasıl Çalışır: Hash Fonksiyonu ve Çakışma}

İki farklı anahtar aynı indekse düşerse buna \emph{çakışma} (collision) denir.
En yaygın çözüm \textbf{zincirleme} (chaining): her kova (bucket) bir bağlı
liste tutar, çakışanlar aynı listeye eklenir. İyi bir hash fonksiyonuyla
çakışma az olur ve işlemler \bigO{1} kalır.

\begin{lstlisting}
class HashSet {
    static const int BUCKETS = 1009;     // asal sayi: dagilim daha iyi
    vector<list<int>> table;
    int hashKey(int key) const { return ((key % BUCKETS) + BUCKETS) % BUCKETS; }
public:
    HashSet() : table(BUCKETS) {}
    void add(int x) {
        auto& bucket = table[hashKey(x)];
        for (int v : bucket) if (v == x) return;   // zaten var
        bucket.push_back(x);
    }
    bool contains(int x) const {
        const auto& bucket = table[hashKey(x)];
        for (int v : bucket) if (v == x) return true;
        return false;
    }
    void remove(int x) { table[hashKey(x)].remove(x); }
};
\end{lstlisting}

\gsub{STL Hash Konteynerleri}

\begin{lstlisting}
#include <unordered_map>
#include <unordered_set>

unordered_map<string, int> counter;
counter["elma"]++;               // yoksa 0'dan baslar, sonra 1
counter["elma"]++;               // 2
cout << counter["elma"];         // 2

unordered_set<int> seen;
seen.insert(5);
if (seen.count(5)) cout << "var\n";
\end{lstlisting}

\gsub{Uygulama: İki Sayının Toplamı (Two Sum)}

Sırasız bir dizide toplamı \code{target} olan iki elemanın indekslerini,
hash haritasıyla tek geçişte \bigO{n}'de buluruz.

\begin{lstlisting}
vector<int> twoSum(const vector<int>& a, int target) {
    unordered_map<int, int> seen;        // deger -> indeks
    for (int i = 0; i < (int)a.size(); ++i) {
        int needed = target - a[i];
        auto it = seen.find(needed);
        if (it != seen.end()) return {it->second, i};
        seen[a[i]] = i;
    }
    return {};
}
\end{lstlisting}

\bilgi{\textbf{\code{map} vs \code{unordered\_map}:} \code{map} kırmızı-siyah
ağaçtır, anahtarları \emph{sıralı} tutar ve tüm işlemler \bigO{\log n}'dir.
\code{unordered\_map} hash tablosudur, sırasızdır ve \emph{ortalama}
\bigO{1}'dir ama en kötü durumda (kötü hash / saldırı) \bigO{n}'e düşebilir.
Sıraya ihtiyacın yoksa ve hız önemliyse \code{unordered\_map} seç.}

\gsec{Temel Veri Yapıları --- Alıştırmalar ve Çözümler}

\bilgi{Bu bölümde dizilerden hash tablolarına kadar temel veri yapılarını, klasik algoritmik desenler (two-pointer, sliding window, prefix sum, monoton yığın/deque) üzerinden çalışacağız. Her sorunun altında Türkçe açıklama, tam çalışan C++ çözümü ve zaman/bellek karmaşıklığı verilmiştir. Kod içinde tanımlar İngilizce, yorumlar Türkçedir.}

\gsub{Diziler}

\soru{İki Toplam --- Sıralı Dizi (Two-Pointer). Artan sırada verilmiş bir \code{nums} dizisi ve bir \code{target} değeri verildiğinde, toplamları hedefe eşit olan iki elemanın (1-tabanlı) indekslerini bulun. Tam olarak bir çözüm olduğunu varsayın.}

\cozum

Dizi sıralı olduğu için tüm ikilileri denemeye gerek yoktur. Bir işaretçiyi başa (\code{left}), diğerini sona (\code{right}) koyarız. Toplam hedeften küçükse \code{left} sağa kayar (toplamı büyütmek için), büyükse \code{right} sola kayar. Bu sayede her elemana en fazla bir kez bakarız.

\begin{lstlisting}
#include <vector>
using namespace std;

vector<int> twoSumSorted(const vector<int>& nums, int target) {
    int left = 0, right = (int)nums.size() - 1;
    while (left < right) {
        int sum = nums[left] + nums[right];
        if (sum == target)
            return {left + 1, right + 1}; // 1-tabanli indeksler
        else if (sum < target)
            ++left;   // toplami buyutmek icin soldan ilerle
        else
            --right;  // toplami kucultmek icin sagdan geri gel
    }
    return {}; // bulunamadi (varsayima gore olmaz)
}
\end{lstlisting}

Karmaşıklık: Zaman \bigO{n}, bellek \bigO{1}. Sıralı olmayan dizide bu yöntem çalışmaz; o durumda hash tablosu (bkz. Two Sum) kullanılır.

\soru{Sıralı Dizinin Karelerini Sıralı Üret. Artan sırada (negatifler dahil) verilmiş \code{nums} dizisinin her elemanının karesini içeren, yine artan sıralı bir dizi döndürün. \bigO{n} olmalı.}

\cozum

Negatif sayıların karesi büyük olabileceğinden en büyük kareler dizinin iki ucundadır. İki işaretçiyi uçlara koyup mutlak değeri büyük olanın karesini sonuç dizisinin sonundan başlayarak yerleştiririz.

\begin{lstlisting}
#include <vector>
using namespace std;

vector<int> sortedSquares(const vector<int>& nums) {
    int n = (int)nums.size();
    vector<int> result(n);
    int left = 0, right = n - 1;
    for (int pos = n - 1; pos >= 0; --pos) {
        int leftSq = nums[left] * nums[left];
        int rightSq = nums[right] * nums[right];
        if (leftSq > rightSq) {   // sol uc daha buyuk kareye sahip
            result[pos] = leftSq;
            ++left;
        } else {
            result[pos] = rightSq;
            --right;
        }
    }
    return result;
}
\end{lstlisting}

Karmaşıklık: Zaman \bigO{n}, bellek \bigO{n} (sonuç dizisi). Önce kareleyip sonra \code{sort} çağırmak \bigO{n \log n} olurdu; two-pointer bunu doğrusala indirir.

\soru{Boyutu k Olan Pencerede Maksimum Toplam (Sliding Window). \code{nums} dizisinde, ardışık \code{k} elemanlık pencereler arasındaki en büyük toplamı bulun.}

\cozum

İlk \code{k} elemanın toplamını hesaplarız. Sonra pencereyi bir sağa kaydırırken çıkan elemanı çıkarıp giren elemanı ekleriz; her pencereyi baştan toplamayız. Bu sabit maliyetli kaydırma sayesinde toplam iş doğrusaldır.

\begin{lstlisting}
#include <vector>
#include <algorithm>
#include <climits>
using namespace std;

int maxSumWindow(const vector<int>& nums, int k) {
    int n = (int)nums.size();
    if (n < k) return 0; // gecersiz durum
    int windowSum = 0;
    for (int i = 0; i < k; ++i) windowSum += nums[i]; // ilk pencere
    int best = windowSum;
    for (int i = k; i < n; ++i) {
        windowSum += nums[i] - nums[i - k]; // gireni ekle, cikani cikar
        best = max(best, windowSum);
    }
    return best;
}
\end{lstlisting}

Karmaşıklık: Zaman \bigO{n}, bellek \bigO{1}.

\soru{En Uzun Tekrarsız Alt Dizi/Alt Karakter Dizisi (Değişken Pencere). Bir \code{string} içinde tekrar eden karakter içermeyen en uzun bitişik alt dizinin uzunluğunu bulun.}

\cozum

Sağ kenarı (\code{right}) ilerletirken bir karakter ikinci kez görülürse, sol kenarı (\code{left}) o karakterin önceki konumundan sonraya çekeriz. Her karakterin son görüldüğü indeksi bir tabloda tutarız. Pencere her zaman tekrarsız kalır.

\begin{lstlisting}
#include <string>
#include <vector>
#include <algorithm>
using namespace std;

int longestUniqueSubstring(const string& s) {
    vector<int> lastSeen(256, -1); // her karakterin son indeksi
    int left = 0, best = 0;
    for (int right = 0; right < (int)s.size(); ++right) {
        unsigned char c = s[right];
        if (lastSeen[c] >= left)
            left = lastSeen[c] + 1; // tekrari disarida birak
        lastSeen[c] = right;
        best = max(best, right - left + 1);
    }
    return best;
}
\end{lstlisting}

Karmaşıklık: Zaman \bigO{n}, bellek \bigO{1} (sabit alfabe). Her kenar en fazla \code{n} kez ilerlediğinden iç içe döngü yoktur.

\soru{Prefix Sum ile Aralık Toplamı Sorguları ve Hedefe Eşit Alt Dizi Sayısı. (a) Çok sayıda \code{[i, j]} aralık toplamı sorgusunu her biri \bigO{1} olacak şekilde yanıtlayın. (b) Toplamı tam \code{target} olan bitişik alt dizilerin sayısını bulun.}

\cozum

Önek toplamı \code{prefix[i]}, ilk \code{i} elemanın toplamıdır. Bir aralığın toplamı \code{prefix[j+1] - prefix[i]} ile sabit zamanda bulunur. Alt dizi sayımı için ise şu gözlemi kullanırız: bir alt dizinin toplamı \code{target} ise, o anki önek toplamından \code{target} çıkardığımız değer daha önce kaç kez görülmüşse o kadar geçerli alt dizi biter. Görülen önek toplamlarını bir hash tablosunda sayarız.

\begin{lstlisting}
#include <vector>
#include <unordered_map>
using namespace std;

// (a) On isleme: aralik toplami sorgulari
struct RangeSum {
    vector<long long> prefix;
    RangeSum(const vector<int>& nums) {
        prefix.assign(nums.size() + 1, 0);
        for (size_t i = 0; i < nums.size(); ++i)
            prefix[i + 1] = prefix[i] + nums[i];
    }
    long long query(int i, int j) const { // kapsayici [i, j]
        return prefix[j + 1] - prefix[i];
    }
};

// (b) Toplami target olan alt dizi sayisi
int countSubarraysWithSum(const vector<int>& nums, int target) {
    unordered_map<long long, int> seen;
    seen[0] = 1;            // bos onek: toplam 0 bir kez gorulmus sayilir
    long long running = 0;
    int count = 0;
    for (int x : nums) {
        running += x;
        auto it = seen.find(running - target);
        if (it != seen.end()) count += it->second;
        ++seen[running];
    }
    return count;
}
\end{lstlisting}

Karmaşıklık: (a) Ön işleme \bigO{n}, her sorgu \bigO{1}. (b) Zaman \bigO{n}, bellek \bigO{n}.

\gsub{Bağlı Listeler}

\bilgi{Aşağıdaki çözümlerde tekil bağlı liste düğümü için standart bir tanım kullanılır:}

\begin{lstlisting}
struct ListNode {
    int val;
    ListNode* next;
    ListNode(int v) : val(v), next(nullptr) {}
};
\end{lstlisting}

\soru{Ortadaki Düğümü Bul (Slow/Fast). Bir bağlı listenin ortasındaki düğümü tek geçişte bulun. Çift uzunlukta iki ortadan ikincisini döndürün.}

\cozum

Klasik "tavşan ve kaplumbağa" tekniği: \code{slow} birer, \code{fast} ikişer adım ilerler. \code{fast} sona ulaştığında \code{slow} tam ortadadır. Böylece uzunluğu önceden saymadan tek geçişte ortayı buluruz.

\begin{lstlisting}
ListNode* middleNode(ListNode* head) {
    ListNode* slow = head;
    ListNode* fast = head;
    while (fast && fast->next) {
        slow = slow->next;        // birer adim
        fast = fast->next->next;  // ikiser adim
    }
    return slow; // cift uzunlukta ikinci ortayi verir
}
\end{lstlisting}

Karmaşıklık: Zaman \bigO{n}, bellek \bigO{1}.

\soru{Bağlı Listeyi Ters Çevir. Bir tekil bağlı listeyi yerinde (in-place) ters çevirip yeni başı döndürün.}

\cozum

Üç işaretçi ile ilerleriz: \code{prev}, \code{curr}, ve geçici \code{nextNode}. Her adımda mevcut düğümün \code{next} bağını bir önceki düğüme çeviririz. Döngü bittiğinde \code{prev} yeni baştır.

\begin{lstlisting}
ListNode* reverseList(ListNode* head) {
    ListNode* prev = nullptr;
    ListNode* curr = head;
    while (curr) {
        ListNode* nextNode = curr->next; // baglantiyi kaybetmemek icin sakla
        curr->next = prev;               // yonu ters cevir
        prev = curr;                     // prev'i ilerlet
        curr = nextNode;                 // curr'i ilerlet
    }
    return prev; // yeni bas
}
\end{lstlisting}

Karmaşıklık: Zaman \bigO{n}, bellek \bigO{1}.

\ipucu{Ters çevirme deseni birçok liste sorusunun yapı taşıdır: palindrom kontrolü, k'lı grup ters çevirme ve listenin ikinci yarısını çevirip birleştirme gibi problemlerde tekrar tekrar karşımıza çıkar.}

\soru{İki Sıralı Listeyi Birleştir. Artan sıralı iki bağlı listeyi tek bir sıralı listede birleştirin.}

\cozum

Bir "kukla" (dummy) baş düğümü kullanmak, sonuç listesinin başını özel durum olarak ele almaktan kurtarır. \code{tail} işaretçisi ile her adımda iki listenin küçük başlı olanını sonuca ekleriz.

\begin{lstlisting}
ListNode* mergeTwoLists(ListNode* a, ListNode* b) {
    ListNode dummy(0);       // kukla bas: kod sadelesir
    ListNode* tail = &dummy;
    while (a && b) {
        if (a->val <= b->val) { tail->next = a; a = a->next; }
        else                  { tail->next = b; b = b->next; }
        tail = tail->next;
    }
    tail->next = a ? a : b;  // kalan kuyrugu ekle
    return dummy.next;
}
\end{lstlisting}

Karmaşıklık: Zaman \bigO{n + m}, bellek \bigO{1} (yeni düğüm oluşturulmaz).

\soru{Döngü Tespiti (Floyd Cycle Detection). Bir bağlı listede döngü (cycle) olup olmadığını sabit bellekte belirleyin.}

\cozum

Yine slow/fast işaretçileri kullanılır. Döngü varsa hızlı işaretçi eninde sonunda yavaş işaretçiyi yakalar (aralarındaki mesafe her adımda bir azalır). Döngü yoksa \code{fast} listenin sonuna (\code{nullptr}) ulaşır.

\begin{lstlisting}
bool hasCycle(ListNode* head) {
    ListNode* slow = head;
    ListNode* fast = head;
    while (fast && fast->next) {
        slow = slow->next;
        fast = fast->next->next;
        if (slow == fast) return true; // isaretciler carpisti: dongu var
    }
    return false; // fast sona ulasti: dongu yok
}
\end{lstlisting}

Karmaşıklık: Zaman \bigO{n}, bellek \bigO{1}.

\soru{Liste Palindrom mu? Bir bağlı listenin değerlerinin palindrom olup olmadığını \bigO{1} ek bellekle kontrol edin.}

\cozum

Üç adım: (1) slow/fast ile ortayı bul, (2) ikinci yarıyı ters çevir, (3) iki yarıyı baştan karşılaştır. Böylece diziyi kopyalamadan sabit bellekle çalışırız.

\begin{lstlisting}
bool isPalindrome(ListNode* head) {
    if (!head || !head->next) return true;
    // 1) ortayi bul
    ListNode* slow = head; ListNode* fast = head;
    while (fast && fast->next) { slow = slow->next; fast = fast->next->next; }
    // 2) ikinci yariyi ters cevir
    ListNode* prev = nullptr;
    while (slow) {
        ListNode* nextNode = slow->next;
        slow->next = prev;
        prev = slow;
        slow = nextNode;
    }
    // 3) iki yariyi karsilastir
    ListNode* firstHalf = head; ListNode* secondHalf = prev;
    while (secondHalf) {
        if (firstHalf->val != secondHalf->val) return false;
        firstHalf = firstHalf->next;
        secondHalf = secondHalf->next;
    }
    return true;
}
\end{lstlisting}

Karmaşıklık: Zaman \bigO{n}, bellek \bigO{1}.

\uyari{Bu çözüm listenin ikinci yarısını kalıcı olarak değiştirir. Çağıran taraf listenin bozulmamasını bekliyorsa, karşılaştırmadan sonra ikinci yarıyı tekrar ters çevirerek eski haline döndürmelisiniz.}

\gsub{Stack (Yığın)}

\soru{Min Stack --- O(1) getMin. \code{push}, \code{pop}, \code{top} ve o anki minimumu döndüren \code{getMin} işlemlerini hepsi \bigO{1} olacak şekilde destekleyen bir yığın tasarlayın.}

\cozum

Ana yığının yanında ikinci bir "minimum yığını" tutarız. Her yeni elemanla birlikte, o ana kadarki minimumu \code{minStack} tepesinde saklarız. Böylece minimum sorgusu sadece tepeye bakmaktır.

\begin{lstlisting}
#include <stack>
#include <algorithm>
using namespace std;

class MinStack {
    stack<int> data;
    stack<int> minStack; // her seviyedeki calisan minimum
public:
    void push(int x) {
        data.push(x);
        if (minStack.empty()) minStack.push(x);
        else minStack.push(min(x, minStack.top()));
    }
    void pop() {
        data.pop();
        minStack.pop();
    }
    int top() const { return data.top(); }
    int getMin() const { return minStack.top(); }
};
\end{lstlisting}

Karmaşıklık: Tüm işlemler \bigO{1}, bellek \bigO{n}.

\soru{Geçerli Parantezler. Yalnızca \code{()[]\{\}} karakterlerinden oluşan bir dizginin doğru şekilde kapanıp kapanmadığını kontrol edin.}

\cozum

Açan parantezleri bir yığına iteriz. Kapatan bir parantez geldiğinde, yığının tepesindeki açanla eşleşmesi gerekir; eşleşmezse veya yığın boşsa dizgi geçersizdir. Sonda yığın boş olmalıdır.

\begin{lstlisting}
#include <string>
#include <stack>
using namespace std;

bool isValidParentheses(const string& s) {
    stack<char> st;
    for (char c : s) {
        if (c == '(' || c == '[' || c == '{') {
            st.push(c); // acani sakla
        } else {
            if (st.empty()) return false; // eslesecek acan yok
            char open = st.top(); st.pop();
            if ((c == ')' && open != '(') ||
                (c == ']' && open != '[') ||
                (c == '}' && open != '{'))
                return false; // yanlis tur eslesme
        }
    }
    return st.empty(); // kapanmamis acan kalmamali
}
\end{lstlisting}

Karmaşıklık: Zaman \bigO{n}, bellek \bigO{n}.

\soru{Postfix (Ters Lehçe) İfade Değerlendirme. Token'lardan oluşan bir postfix ifadeyi (örneğin \code{["2","1","+","3","*"]}) değerlendirin.}

\cozum

Postfix ifadeler yığın için biçilmiş kaftandır: sayı görünce yığına iteriz, operatör görünce tepeden iki operandı çekip işlemi yapar ve sonucu geri iteriz. İkinci çekilen operand, işlemin sol tarafıdır (çıkarma ve bölmede sıra önemlidir).

\begin{lstlisting}
#include <vector>
#include <string>
#include <stack>
using namespace std;

int evalRPN(const vector<string>& tokens) {
    stack<long long> st;
    for (const string& tok : tokens) {
        if (tok == "+" || tok == "-" || tok == "*" || tok == "/") {
            long long b = st.top(); st.pop(); // sag operand
            long long a = st.top(); st.pop(); // sol operand
            if      (tok == "+") st.push(a + b);
            else if (tok == "-") st.push(a - b);
            else if (tok == "*") st.push(a * b);
            else                 st.push(a / b); // tam sayi bolmesi
        } else {
            st.push(stoll(tok)); // sayiyi yigina it
        }
    }
    return (int)st.top();
}
\end{lstlisting}

Karmaşıklık: Zaman \bigO{n}, bellek \bigO{n}.

\soru{Günlük Sıcaklıklar (Monoton Yığın / Next Greater). Her gün için, daha sıcak bir günün gelmesine kaç gün kaldığını bulun. Yoksa 0 yazın.}

\cozum

İndeksleri, sıcaklıkları azalan bir yığında tutarız (monoton azalan yığın). Yeni bir gün, yığının tepesindeki günden daha sıcaksa, o tepe gün için cevabı bulmuş oluruz; onu çıkarıp gün farkını yazarız. Her indeks yığına en fazla bir kez girip çıktığından toplam iş doğrusaldır.

\begin{lstlisting}
#include <vector>
#include <stack>
using namespace std;

vector<int> dailyTemperatures(const vector<int>& temps) {
    int n = (int)temps.size();
    vector<int> answer(n, 0);
    stack<int> st; // cozulmemis gunlerin indeksleri (azalan sicaklik)
    for (int i = 0; i < n; ++i) {
        while (!st.empty() && temps[i] > temps[st.top()]) {
            int prev = st.top(); st.pop();
            answer[prev] = i - prev; // kac gun sonra daha sicak
        }
        st.push(i);
    }
    return answer;
}
\end{lstlisting}

Karmaşıklık: Zaman \bigO{n}, bellek \bigO{n}. Monoton yığın, "next greater/smaller element" ailesindeki tüm problemlerin ortak kalıbıdır.

\gsub{Queue ve Deque}

\soru{İki Stack ile Kuyruk (Queue) Uygula. FIFO davranışı gösteren \code{push}, \code{pop} ve \code{peek} işlemlerini yalnızca iki yığın kullanarak gerçekleştirin.}

\cozum

Bir \code{inStack} gelen elemanları alır, bir \code{outStack} çıkışları verir. \code{outStack} boşaldığında \code{inStack}'in tamamını ona aktarırız; bu aktarma sırayı ters çevirdiği için en eski eleman \code{outStack}'in tepesine gelir. Amortize edilmiş maliyet sabittir: her eleman ömrü boyunca en fazla bir kez aktarılır.

\begin{lstlisting}
#include <stack>
using namespace std;

class MyQueue {
    stack<int> inStack;  // gelenler
    stack<int> outStack; // cikacaklar (ters sirada)
    void transferIfNeeded() {
        if (outStack.empty()) {
            while (!inStack.empty()) {
                outStack.push(inStack.top()); // sirayi ters cevirerek aktar
                inStack.pop();
            }
        }
    }
public:
    void push(int x) { inStack.push(x); }
    int pop() {
        transferIfNeeded();
        int front = outStack.top(); outStack.pop();
        return front;
    }
    int peek() {
        transferIfNeeded();
        return outStack.top();
    }
    bool empty() const { return inStack.empty() && outStack.empty(); }
};
\end{lstlisting}

Karmaşıklık: \code{push} \bigO{1}; \code{pop} ve \code{peek} amortize \bigO{1} (en kötü tek çağrı \bigO{n}). Bellek \bigO{n}.

\soru{Kayan Pencere Maksimumu (Monoton Deque). \code{nums} dizisinde \code{k} boyutlu her pencerenin maksimumunu içeren diziyi \bigO{n} zamanda üretin.}

\cozum

İndeksleri, değerleri azalan bir çift uçlu kuyrukta (deque) tutarız. Yeni eleman gelince arkadan ondan küçük veya eşit tüm indeksleri atarız (onlar artık asla maksimum olamaz). Önden ise pencereden çıkmış indeksleri atarız. Böylece deque'in başı her zaman mevcut pencerenin maksimumudur.

\begin{lstlisting}
#include <vector>
#include <deque>
using namespace std;

vector<int> maxSlidingWindow(const vector<int>& nums, int k) {
    vector<int> result;
    deque<int> dq; // indeksler, degerleri azalan sirada
    for (int i = 0; i < (int)nums.size(); ++i) {
        // arkadan: yeni elemandan kucuk/esit olanlar gereksiz
        while (!dq.empty() && nums[dq.back()] <= nums[i])
            dq.pop_back();
        dq.push_back(i);
        // onden: pencere disina cikan indeksi at
        if (dq.front() <= i - k)
            dq.pop_front();
        // pencere tamamlaninca maksimumu yaz
        if (i >= k - 1)
            result.push_back(nums[dq.front()]);
    }
    return result;
}
\end{lstlisting}

Karmaşıklık: Zaman \bigO{n} (her indeks en fazla bir kez girip çıkar), bellek \bigO{k}.

\gsub{Hash Tabloları}

\soru{İki Toplam (Two Sum --- İndeksler). Sırasız bir \code{nums} dizisinde, toplamları \code{target} olan iki elemanın indekslerini döndürün.}

\cozum

Diziyi tararken her eleman için "aradığım tamamlayıcı" olan \code{target - nums[i]} değerini daha önce görüp görmediğimizi hash tablosunda kontrol ederiz. Görülen değerleri indeksleriyle birlikte tabloya ekleriz. Tek geçişte çözülür.

\begin{lstlisting}
#include <vector>
#include <unordered_map>
using namespace std;

vector<int> twoSum(const vector<int>& nums, int target) {
    unordered_map<int, int> seen; // deger -> indeks
    for (int i = 0; i < (int)nums.size(); ++i) {
        int need = target - nums[i]; // aranan tamamlayici
        auto it = seen.find(need);
        if (it != seen.end())
            return {it->second, i};
        seen[nums[i]] = i;
    }
    return {};
}
\end{lstlisting}

Karmaşıklık: Zaman ortalama \bigO{n}, bellek \bigO{n}.

\soru{Anagramları Grupla. Bir kelime listesini, harfleri aynı olan (anagram) kelimeler aynı grupta olacak şekilde gruplayın.}

\cozum

Her kelimenin harflerini sıraladığımızda, aynı anagram grubundaki tüm kelimeler aynı "imzayı" verir. Bu sıralı imzayı hash tablosunda anahtar olarak kullanıp kelimeleri gruplarız.

\begin{lstlisting}
#include <vector>
#include <string>
#include <unordered_map>
#include <algorithm>
using namespace std;

vector<vector<string>> groupAnagrams(const vector<string>& words) {
    unordered_map<string, vector<string>> groups;
    for (const string& w : words) {
        string key = w;
        sort(key.begin(), key.end()); // anagram imzasi
        groups[key].push_back(w);
    }
    vector<vector<string>> result;
    result.reserve(groups.size());
    for (auto& kv : groups)
        result.push_back(move(kv.second));
    return result;
}
\end{lstlisting}

Karmaşıklık: \code{n} kelime ve en fazla \code{L} uzunluk için zaman \bigO{n \cdot L \log L}, bellek \bigO{n \cdot L}. Alternatif olarak 26 harf sayımıyla oluşturulan imza sıralamayı ortadan kaldırır.

\soru{En Sık k Eleman. \code{nums} dizisinde en sık geçen \code{k} elemanı döndürün.}

\cozum

Önce hash tablosuyla her elemanın frekansını sayarız. Sonra en sık \code{k} elemanı seçmek için bir min-heap (öncelik kuyruğu) tutarız; boyutu \code{k}'yı aşınca en az frekanslıyı atarız. Böylece heap işlemleri \bigO{\log k} ile sınırlı kalır.

\begin{lstlisting}
#include <vector>
#include <unordered_map>
#include <queue>
using namespace std;

vector<int> topKFrequent(const vector<int>& nums, int k) {
    unordered_map<int, int> freq;
    for (int x : nums) ++freq[x]; // frekanslari say
    // min-heap: (frekans, deger); en az sik tepede
    priority_queue<pair<int,int>, vector<pair<int,int>>, greater<>> pq;
    for (auto& kv : freq) {
        pq.push({kv.second, kv.first});
        if ((int)pq.size() > k) pq.pop(); // en az sikligi cikar
    }
    vector<int> result;
    while (!pq.empty()) { result.push_back(pq.top().second); pq.pop(); }
    return result;
}
\end{lstlisting}

Karmaşıklık: Zaman \bigO{n \log k}, bellek \bigO{n}. Frekansa göre bucket sort ile \bigO{n} de yapılabilir.

\soru{İlk Tekrarsız Karakter. Bir dizgide yalnızca bir kez geçen ilk karakterin indeksini döndürün; yoksa -1 döndürün.}

\cozum

İki geçiş: birinci geçişte her karakterin sayısını hash tablosunda (veya sabit boyutlu dizide) tutarız; ikinci geçişte soldan ilk kez sayısı 1 olan karakteri buluruz. Dizgi sırasını koruduğumuz için ikinci geçiş "ilk" olanı garanti eder.

\begin{lstlisting}
#include <string>
#include <array>
using namespace std;

int firstUniqueChar(const string& s) {
    array<int, 256> count{}; // sifirla baslatilir
    for (unsigned char c : s) ++count[c]; // 1. gecis: say
    for (int i = 0; i < (int)s.size(); ++i)
        if (count[(unsigned char)s[i]] == 1) // 2. gecis: ilk tekil
            return i;
    return -1;
}
\end{lstlisting}

Karmaşıklık: Zaman \bigO{n}, bellek \bigO{1} (sabit alfabe).

\gsub{Uygulama: LRU Cache}

\notk[LRU Cache Nedir?]{LRU (Least Recently Used --- En Az Kısa Süre Önce Kullanılan) önbellek, sabit kapasiteli bir yapıdır: kapasite dolduğunda, en uzun süredir dokunulmamış (en eski erişilen) elemanı atarak yer açar. Amaç hem \code{get} hem \code{put} işlemini \bigO{1} yapmaktır. Bunun için iki veri yapısını birlikte kullanırız: bir \code{unordered\_map} (anahtardan düğüme \bigO{1} erişim için) ve bir çift yönlü bağlı liste (\code{std::list}, kullanım sırasını korumak ve herhangi bir elemanı \bigO{1} taşımak için). Liste başı = en son kullanılan; liste sonu = atılacak aday.}

\soru{LRU Cache Tam Implementasyonu. \code{capacity} kapasiteli bir \code{LRUCache} sınıfı yazın. \code{get(key)} anahtarı varsa değerini döndürsün ve onu "en son kullanılan" yapsın (yoksa -1). \code{put(key, value)} ekleme/güncelleme yapsın; kapasite aşılırsa en az kullanılanı atsın. Her iki işlem de \bigO{1} olsun.}

\cozum

\code{std::list<pair<int,int>>} kullanım sırasını (baş = en yeni) tutar; her düğüm \code{(key, value)} çiftidir. \code{unordered\_map<int, list iterator>} her anahtarı listedeki düğümüne bağlar. \code{list::splice} ile bir düğümü kopyalamadan, iteratörlerini geçersiz kılmadan listenin başına \bigO{1} taşıyabiliriz --- bu, tasarımın kilit noktasıdır.

\begin{lstlisting}
#include <list>
#include <unordered_map>
#include <utility>
using namespace std;

class LRUCache {
    int capacity;
    // liste: on = en son kullanilan, arka = en az kullanilan
    list<pair<int,int>> items;                 // (key, value)
    unordered_map<int, list<pair<int,int>>::iterator> index; // key -> dugum
public:
    explicit LRUCache(int capacity) : capacity(capacity) {}

    int get(int key) {
        auto it = index.find(key);
        if (it == index.end()) return -1; // bulunamadi
        // erisildi: bu dugumu listenin basina tasi (en son kullanilan)
        items.splice(items.begin(), items, it->second);
        return it->second->second; // value
    }

    void put(int key, int value) {
        auto it = index.find(key);
        if (it != index.end()) {
            // var olan anahtar: degeri guncelle ve basa tasi
            it->second->second = value;
            items.splice(items.begin(), items, it->second);
            return;
        }
        // kapasite doluysa en az kullanilani (listenin arkasi) at
        if ((int)items.size() == capacity) {
            int oldKey = items.back().first;
            index.erase(oldKey);
            items.pop_back();
        }
        // yeni elemani basa ekle ve indeksle
        items.emplace_front(key, value);
        index[key] = items.begin();
    }
};
\end{lstlisting}

Karmaşıklık: \code{get} ve \code{put} her ikisi de \bigO{1} (splice, emplace, erase hepsi sabit zaman). Bellek \bigO{capacity}.

\uyari{\code{splice} kullanmanın kritik nedeni, taşınan düğümün iteratörünü geçersiz kılmamasıdır. Bu yüzden \code{index} tablosundaki iteratörler taşımadan sonra da geçerli kalır; ayrıca yeniden ekleme yapmaya gerek kalmaz. \code{items.erase} + \code{push\_front} ile yapsaydınız iteratörü güncellemeyi unutmak yaygın bir hata olurdu.}

\soru{LRU Cache Kullanımı. Kapasitesi 2 olan bir \code{LRUCache} üzerinde şu işlemler sırayla uygulanıyor: \code{put(1,1)}, \code{put(2,2)}, \code{get(1)}, \code{put(3,3)}, \code{get(2)}, \code{put(4,4)}, \code{get(1)}, \code{get(3)}, \code{get(4)}. Her \code{get} çağrısının çıktısını ve \code{put(3,3)} ile \code{put(4,4)} sırasında hangi anahtarın atıldığını izleyin.}

\cozum

Adım adım (liste solu = en son kullanılan):

\begin{tabularx}{\textwidth}{@{}lX@{}}
İşlem & Durum ve Sonuç \\
\code{put(1,1)} & liste: [1] \\
\code{put(2,2)} & liste: [2, 1] \\
\code{get(1)} & 1 döner; 1 başa taşınır $\to$ [1, 2] \\
\code{put(3,3)} & kapasite dolu; en eski olan \textbf{2 atılır} $\to$ [3, 1] \\
\code{get(2)} & 2 atıldığı için \textbf{-1} döner \\
\code{put(4,4)} & kapasite dolu; en eski olan \textbf{1 atılır} $\to$ [4, 3] \\
\code{get(1)} & 1 atıldığı için \textbf{-1} döner \\
\code{get(3)} & \textbf{3} döner; 3 başa taşınır $\to$ [3, 4] \\
\code{get(4)} & \textbf{4} döner; 4 başa taşınır $\to$ [4, 3] \\
\end{tabularx}

Doğrulama kodu:

\begin{lstlisting}
#include <iostream>

int main() {
    LRUCache cache(2);
    cache.put(1, 1);
    cache.put(2, 2);
    std::cout << cache.get(1) << "\n"; // 1
    cache.put(3, 3);                   // 2 atilir
    std::cout << cache.get(2) << "\n"; // -1
    cache.put(4, 4);                   // 1 atilir
    std::cout << cache.get(1) << "\n"; // -1
    std::cout << cache.get(3) << "\n"; // 3
    std::cout << cache.get(4) << "\n"; // 4
    return 0;
}
\end{lstlisting}

Beklenen çıktı satırları: \code{1}, \code{-1}, \code{-1}, \code{3}, \code{4}. Bu örnek, LRU mantığında bir elemana \code{get} veya \code{put} ile "dokunmanın" onu atılmaktan koruduğunu, dokunulmayan elemanınsa zamanla atıldığını gösterir.

\bilgi{Karşılaştırmalı özet --- işlem başına tipik karmaşıklıklar: two-pointer / sliding window / prefix sum $\to$ \bigO{n}; bağlı liste temel işlemleri $\to$ \bigO{n} zaman, \bigO{1} bellek; yığın/kuyruk işlemleri $\to$ \bigO{1} (kuyruk için amortize); hash tablosu arama/ekleme $\to$ ortalama \bigO{1}; LRU Cache get/put $\to$ \bigO{1}. Bu desenleri tanımak, birçok karmaşık problemi tanıdık yapı taşlarına indirgemenizi sağlar.}

% #####################################################################
\gpart{3}{Ağaçlar}
% #####################################################################

% =====================================================================
\gsec{Ağaçlar ve İkili Ağaçlar}
% =====================================================================

Ağaç, düğümlerin hiyerarşik olarak bağlandığı, döngüsüz bir yapıdır. En üstteki
düğüm \emph{kök} (root), altındakiler \emph{çocuk} (child), çocuğu olmayanlar
\emph{yaprak} (leaf) düğümlerdir. Dosya sistemleri, ifade ağaçları, arama
yapıları ve daha fazlası ağaçlarla modellenir.

\begin{center}
\begin{tabularx}{\textwidth}{@{}>{\bfseries\color{anaMavi}}l X@{}}
\toprule
\rowcolor{acikMavi}\color{anaMavi}\textbf{Terim} & \color{anaMavi}\textbf{Anlamı}\\
\midrule
Kök (root) & En üstteki, ebeveyni olmayan düğüm.\\
\rowcolor{acikGri} Yaprak (leaf) & Çocuğu olmayan düğüm.\\
Derinlik (depth) & Kökten o düğüme olan kenar sayısı.\\
\rowcolor{acikGri} Yükseklik (height) & Düğümden en uzak yaprağa olan kenar sayısı.\\
Alt ağaç (subtree) & Bir düğüm ve altındaki tüm düğümler.\\
\bottomrule
\end{tabularx}
\end{center}

\gsub{İkili Ağaç ve Dolaşımlar (Traversals)}

İkili ağaçta her düğümün en fazla iki çocuğu (sol, sağ) vardır. Bir ağacı
``dolaşmak'', tüm düğümleri belirli bir sırada ziyaret etmektir. Üç temel
derinlik-öncelikli dolaşım vardır:

\begin{center}
\begin{tabularx}{\textwidth}{@{}>{\bfseries\color{anaMavi}}l l X@{}}
\toprule
\rowcolor{acikMavi}\color{anaMavi}\textbf{Dolaşım} & \color{anaMavi}\textbf{Sıra} & \color{anaMavi}\textbf{Kullanım}\\
\midrule
Pre-order & Kök, Sol, Sağ & Ağacı kopyalama, önek ifade.\\
\rowcolor{acikGri} In-order & Sol, Kök, Sağ & BST'de \emph{sıralı} çıktı verir.\\
Post-order & Sol, Sağ, Kök & Ağacı silme, sonek ifade.\\
\bottomrule
\end{tabularx}
\end{center}

\begin{lstlisting}
struct TreeNode {
    int       data;
    TreeNode* left = nullptr;
    TreeNode* right = nullptr;
    TreeNode(int v) : data(v) {}
};

void inorder(TreeNode* k) {              // Sol -> Kok -> Sag
    if (!k) return;
    inorder(k->left);
    cout << k->data << " ";
    inorder(k->right);
}

void preorder(TreeNode* k) {             // Kok -> Sol -> Sag
    if (!k) return;
    cout << k->data << " ";
    preorder(k->left);
    preorder(k->right);
}

void postorder(TreeNode* k) {            // Sol -> Sag -> Kok
    if (!k) return;
    postorder(k->left);
    postorder(k->right);
    cout << k->data << " ";
}
\end{lstlisting}

\gsub{Seviye-Sıralı Dolaşım (BFS) ve Yükseklik}

Seviye sıralı dolaşım, ağacı katman katman gezer ve bir kuyruk kullanır.
Yükseklik ise özyinelemeyle bulunur.

\begin{lstlisting}
void levelOrder(TreeNode* root) {
    if (!root) return;
    queue<TreeNode*> q;
    q.push(root);
    while (!q.empty()) {
        TreeNode* d = q.front(); q.pop();
        cout << d->data << " ";
        if (d->left) q.push(d->left);
        if (d->right) q.push(d->right);
    }
}

int height(TreeNode* k) {
    if (!k) return -1;                    // bos agac: -1 (kenar sayisi)
    return 1 + max(height(k->left), height(k->right));
}
\end{lstlisting}

% =====================================================================
\gsec{İkili Arama Ağacı (BST)}
% =====================================================================

İkili arama ağacı (Binary Search Tree), \emph{sıralı} bir ikili ağaçtır: her
düğüm için \textbf{sol alt ağacındaki tüm değerler ondan küçük, sağ alt
ağacındakiler büyüktür}. Bu değişmez (invariant) sayesinde arama, ekleme ve
silme, ağacın yüksekliği kadar sürede --- dengeli bir ağaçta \bigO{\log n}'de
--- yapılır.

\gsub{Arama ve Ekleme}

Aranan değeri kökle karşılaştırırız: küçükse sola, büyükse sağa gideriz.
Her adımda arama uzayını yarıya indiririz.

\begin{lstlisting}
struct Node {
    int   data;
    Node* left = nullptr;
    Node* right = nullptr;
    Node(int v) : data(v) {}
};

Node* search(Node* k, int x) {           // O(h)
    while (k && k->data != x)
        k = (x < k->data) ? k->left : k->right;
    return k;                             // bulunamazsa nullptr
}

Node* insert(Node* k, int x) {           // O(h)
    if (!k) return new Node(x);
    if (x < k->data)      k->left = insert(k->left, x);
    else if (x > k->data) k->right = insert(k->right, x);
    // esitse: eklemiyoruz (tekrarsiz BST)
    return k;
}
\end{lstlisting}

\gsub{Silme: Üç Durum}

Silme en incelikli işlemdir çünkü üç ayrı durum vardır. Anahtar fikir: iki
çocuklu bir düğümü silerken, onu \emph{sağ alt ağacının en küçük değeri}
(in-order ardılı) ile değiştiririz.

\begin{center}
\begin{tabularx}{\textwidth}{@{}>{\bfseries\color{anaMavi}}l X@{}}
\toprule
\rowcolor{acikMavi}\color{anaMavi}\textbf{Durum} & \color{anaMavi}\textbf{Yapılan}\\
\midrule
Yaprak (çocuk yok) & Doğrudan sil.\\
\rowcolor{acikGri} Tek çocuk & Düğümü çocuğuyla değiştir.\\
İki çocuk & In-order ardılı ile değiştir, sonra ardılı sil.\\
\bottomrule
\end{tabularx}
\end{center}

\begin{lstlisting}
Node* findMin(Node* k) { while (k->left) k = k->left; return k; }

Node* remove(Node* k, int x) {
    if (!k) return nullptr;
    if (x < k->data)      k->left = remove(k->left, x);
    else if (x > k->data) k->right = remove(k->right, x);
    else {                                // dugum bulundu
        if (!k->left) { Node* s = k->right; delete k; return s; }
        if (!k->right) { Node* s = k->left; delete k; return s; }
        Node* successor = findMin(k->right);   // iki cocuklu durum
        k->data = successor->data;        // ardili yukari tasi
        k->right = remove(k->right, successor->data);
    }
    return k;
}
\end{lstlisting}

\uyari{BST'nin \bigO{\log n} garantisi \emph{dengeli} olmasına bağlıdır.
Sıralı veriyi (1, 2, 3, 4, \dots) sırayla eklersen ağaç dev bir bağlı listeye
dönüşür ve işlemler \bigO{n}'e çıkar. Bu problemi bir sonraki bölümdeki
kendi-kendini dengeleyen ağaçlar çözer.}

% =====================================================================
\gsec{Dengeli Ağaçlar: AVL ve Kırmızı-Siyah}
% =====================================================================

Dengeli arama ağaçları, her eklemeden/silmeden sonra yapıyı yeniden düzenleyerek
yüksekliği \bigO{\log n}'de tutar; böylece BST'nin en kötü durum sorununu ortadan
kaldırır. Dengeleme, \emph{rotasyon} (rotation) adı verilen \bigO{1}'lik yerel
işlemlerle yapılır.

\gsub{Rotasyonlar: Dengelemenin Temeli}

Rotasyon, BST sırasını bozmadan düğümleri yeniden bağlayarak alt ağaç
yüksekliklerini değiştirir. Sol ve sağ rotasyon birbirinin simetriğidir.

\begin{lstlisting}
struct AVLNode {
    int      data, height = 1;
    AVLNode* left = nullptr;
    AVLNode* right = nullptr;
    AVLNode(int v) : data(v) {}
};

int getHeight(AVLNode* d) { return d ? d->height : 0; }
int balance(AVLNode* d) { return d ? getHeight(d->left) - getHeight(d->right) : 0; }
void update(AVLNode* d) { d->height = 1 + max(getHeight(d->left), getHeight(d->right)); }

AVLNode* rotateRight(AVLNode* y) {       // sol agir dugumu dengele
    AVLNode* x = y->left;
    AVLNode* T = x->right;
    x->right = y;  y->left = T;           // rotasyon
    update(y); update(x);
    return x;                             // yeni kok
}

AVLNode* rotateLeft(AVLNode* x) {        // sag agir dugumu dengele
    AVLNode* y = x->right;
    AVLNode* T = y->left;
    y->left = x;  x->right = T;
    update(x); update(y);
    return y;
}
\end{lstlisting}

\gsub{AVL Ağacına Ekleme}

AVL ağacında her düğüm için \emph{denge faktörü} (sol yükseklik $-$ sağ
yükseklik) $\{-1, 0, +1\}$ aralığında tutulur. Ekleme sonrası bir düğüm bu
aralığı aşarsa, dört durumdan (LL, RR, LR, RL) uygun rotasyon(lar) uygulanır.

\begin{lstlisting}
AVLNode* insert(AVLNode* d, int x) {
    if (!d) return new AVLNode(x);
    if (x < d->data)      d->left = insert(d->left, x);
    else if (x > d->data) d->right = insert(d->right, x);
    else return d;                        // tekrar yok
    update(d);

    int b = balance(d);
    if (b > 1 && x < d->left->data) return rotateRight(d);        // LL
    if (b < -1 && x > d->right->data) return rotateLeft(d);       // RR
    if (b > 1 && x > d->left->data) {                             // LR
        d->left = rotateLeft(d->left); return rotateRight(d);
    }
    if (b < -1 && x < d->right->data) {                          // RL
        d->right = rotateRight(d->right); return rotateLeft(d);
    }
    return d;
}
\end{lstlisting}

\gsub{AVL vs Kırmızı-Siyah Ağaç}

\begin{center}
\begin{tabularx}{\textwidth}{@{}>{\bfseries\color{anaMavi}}l X X@{}}
\toprule
\rowcolor{acikMavi}\color{anaMavi}\textbf{Özellik} & \color{anaMavi}\textbf{AVL} & \color{anaMavi}\textbf{Kırmızı-Siyah}\\
\midrule
Denge & Katı (daha kısa) & Gevşek\\
\rowcolor{acikGri} Arama & Daha hızlı & Biraz yavaş\\
Ekleme/silme & Daha çok rotasyon & Daha az rotasyon\\
\rowcolor{acikGri} Kullanım & Okuma ağırlıklı & \code{std::map}, \code{std::set}\\
\bottomrule
\end{tabularx}
\end{center}

\bilgi{C++'ın \code{std::map} ve \code{std::set} yapıları içeride
kırmızı-siyah ağaç kullanır. Yani pratikte kendi dengeli ağacını yazman
nadiren gerekir; ama nasıl çalıştığını bilmek, \bigO{\log n} garantisinin
nereden geldiğini anlamanı sağlar.}

% =====================================================================
\gsec{Heap ve Öncelik Kuyruğu}
% =====================================================================

Heap (yığıt/öbek), \textbf{tam ikili ağaç} (complete binary tree) biçiminde
bir yapıdır ve bir \emph{öncelik} değişmezini korur: \textbf{max-heap}'te her
ebeveyn çocuklarından büyük veya eşittir (min-heap'te küçük veya eşit). Bu
sayede en büyük (veya en küçük) elemana \bigO{1}'de erişilir; ekleme ve
çıkarma \bigO{\log n}'dir. Öncelik kuyruğu (priority queue) heap ile uygulanır.

\gsub{Diziyle Heap Gösterimi}

Tam ikili ağaç olduğundan heap'i bir dizide, işaretçi olmadan tutabiliriz.
$i$ indeksli düğüm için: sol çocuk $2i+1$, sağ çocuk $2i+2$, ebeveyn
$(i-1)/2$'dir.

\begin{lstlisting}
class MaxHeap {
    vector<int> h;
    void siftUp(int i) {                  // ekleme sonrasi: O(log n)
        while (i > 0 && h[(i - 1) / 2] < h[i]) {
            swap(h[i], h[(i - 1) / 2]);
            i = (i - 1) / 2;
        }
    }
    void siftDown(int i) {                // cikarma sonrasi: O(log n)
        int n = h.size();
        while (true) {
            int largest = i, left = 2*i + 1, right = 2*i + 2;
            if (left < n && h[left] > h[largest]) largest = left;
            if (right < n && h[right] > h[largest]) largest = right;
            if (largest == i) break;
            swap(h[i], h[largest]);
            i = largest;
        }
    }
public:
    void push(int x) { h.push_back(x); siftUp(h.size() - 1); }
    int  top() const { return h[0]; }             // O(1)
    void pop() {                                  // en buyugu cikar
        h[0] = h.back(); h.pop_back();
        if (!h.empty()) siftDown(0);
    }
    bool empty() const { return h.empty(); }
};
\end{lstlisting}

\gsub{STL Öncelik Kuyruğu}

\begin{lstlisting}
#include <queue>
// Varsayilan: MAX-heap (en buyuk tepede)
priority_queue<int> pq;
pq.push(3); pq.push(1); pq.push(4);
cout << pq.top();   // 4
pq.pop();

// MIN-heap icin:
priority_queue<int, vector<int>, greater<int>> minpq;
minpq.push(3); minpq.push(1); minpq.push(4);
cout << minpq.top();   // 1
\end{lstlisting}

\gsub{Uygulama: En Küçük k Eleman (ve Heapsort)}

Bir akıştaki en küçük \code{k} elemanı, boyutu \code{k} olan bir \emph{max-heap}
ile \bigO{n\log k}'de buluruz: heap doluysa ve yeni eleman tepeden küçükse
tepeyi at, yenisini ekle.

\begin{lstlisting}
vector<int> kSmallest(const vector<int>& a, int k) {
    priority_queue<int> maxHeap;         // boyutu en fazla k
    for (int x : a) {
        maxHeap.push(x);
        if ((int)maxHeap.size() > k) maxHeap.pop();  // en buyugu at
    }
    vector<int> result;
    while (!maxHeap.empty()) { result.push_back(maxHeap.top()); maxHeap.pop(); }
    return result;
}
\end{lstlisting}

\notk[Heapsort]{Tüm elemanları bir heap'e atıp tek tek çıkarırsak sıralı elde
ederiz: bu \bigO{n\log n}'lik \emph{heapsort} algoritmasıdır. STL'de
\code{std::make\_heap}, \code{std::push\_heap}, \code{std::pop\_heap} ile aynı
mekanizmaya doğrudan erişilir.}

% =====================================================================
\gsec{Trie (Önek Ağacı)}
% =====================================================================

Trie (önek ağacı / prefix tree), string'leri karakter karakter saklayan bir
ağaçtır. Kökten bir düğüme giden yol, bir öneki temsil eder. Bir kelimenin
aranması, eklenmesi veya bir önekle başlayan kelimelerin bulunması, kelime
uzunluğu $L$ olmak üzere \bigO{L}'de yapılır --- sözlük boyutundan bağımsız!
Otomatik tamamlama, yazım denetimi ve IP yönlendirme trie kullanır.

\begin{lstlisting}
struct TrieNode {
    TrieNode* children[26] = {nullptr};  // 'a'-'z'
    bool isEndOfWord = false;
};

class Trie {
    TrieNode* root = new TrieNode();
public:
    void insert(const string& s) {       // O(L)
        TrieNode* d = root;
        for (char c : s) {
            int i = c - 'a';
            if (!d->children[i]) d->children[i] = new TrieNode();
            d = d->children[i];
        }
        d->isEndOfWord = true;
    }
    bool search(const string& s) const { // tam kelime var mi?
        TrieNode* d = root;
        for (char c : s) {
            int i = c - 'a';
            if (!d->children[i]) return false;
            d = d->children[i];
        }
        return d->isEndOfWord;
    }
    bool startsWith(const string& prefix) const {  // bu onekle baslayan var mi?
        TrieNode* d = root;
        for (char c : prefix) {
            int i = c - 'a';
            if (!d->children[i]) return false;
            d = d->children[i];
        }
        return true;                      // isEndOfWord'e bakmiyoruz
    }
};
\end{lstlisting}

\ipucu{Trie, bellek karşılığında hız satın alır: her düğüm 26 işaretçi
tutabilir. Bellek kritikse çocukları \code{unordered\_map<char, TrieNode*>}
ile seyrek tutmak veya sıkıştırılmış trie (radix tree) kullanmak daha
verimlidir.}

% =====================================================================
\gsec{Segment Ağacı ve Fenwick Ağacı}
% =====================================================================

Bir dizide hem \emph{aralık sorguları} (örneğin ``$[l,r]$ toplamı / minimumu'')
hem de \emph{güncellemeler} sık yapılıyorsa, düz dizi yetersizdir: sorgu
\bigO{n} sürer. Segment ağacı ve Fenwick ağacı, her ikisini de
\bigO{\log n}'de yapar.

\gsub{Segment Ağacı}

Segment ağacı, diziyi aralıklara bölen bir ikili ağaçtır: her düğüm bir
aralığın özetini (toplam, min, max\dots) tutar. Kök tüm diziyi, yapraklar tek
elemanları temsil eder.

\begin{lstlisting}
class SegmentTree {
    int n;
    vector<long long> tree;              // 1-tabanli, 4n boyut
    void build(const vector<int>& a, int node, int l, int r) {
        if (l == r) { tree[node] = a[l]; return; }
        int mid = (l + r) / 2;
        build(a, 2*node,   l, mid);
        build(a, 2*node+1, mid+1, r);
        tree[node] = tree[2*node] + tree[2*node+1];   // toplam birlestir
    }
    long long query(int node, int l, int r, int ql, int qr) {
        if (qr < l || r < ql) return 0;             // hic kesismiyor
        if (ql <= l && r <= qr) return tree[node];   // tamamen iceride
        int mid = (l + r) / 2;
        return query(2*node, l, mid, ql, qr)
             + query(2*node+1, mid+1, r, ql, qr);
    }
    void update(int node, int l, int r, int i, int value) {
        if (l == r) { tree[node] = value; return; }
        int mid = (l + r) / 2;
        if (i <= mid) update(2*node, l, mid, i, value);
        else          update(2*node+1, mid+1, r, i, value);
        tree[node] = tree[2*node] + tree[2*node+1];
    }
public:
    SegmentTree(const vector<int>& a) : n(a.size()), tree(4*n) {
        build(a, 1, 0, n-1);
    }
    long long query(int l, int r)   { return query(1, 0, n-1, l, r); }
    void update(int i, int value)   { update(1, 0, n-1, i, value); }
};
\end{lstlisting}

\gsub{Fenwick Ağacı (Binary Indexed Tree)}

Fenwick ağacı, yalnızca \emph{önek toplamı} tipi sorgular için segment ağacına
göre çok daha kısa ve az bellekli bir alternatiftir. Sihri, bir indeksin son
kurulu bitini (\code{i \& -i}) kullanarak ağaçta zıplamaktır.

\begin{lstlisting}
class Fenwick {
    int n;
    vector<long long> tree;              // 1-tabanli
public:
    Fenwick(int n) : n(n), tree(n + 1, 0) {}

    void add(int i, long long value) {   // a[i] += value, O(log n)
        for (++i; i <= n; i += i & -i)   // son biti ekleyerek ilerle
            tree[i] += value;
    }
    long long prefixSum(int i) {         // a[0..i] toplami, O(log n)
        long long s = 0;
        for (++i; i > 0; i -= i & -i)    // son biti cikararak geri git
            s += tree[i];
        return s;
    }
    long long rangeSum(int l, int r) {
        return prefixSum(r) - (l > 0 ? prefixSum(l - 1) : 0);
    }
};
\end{lstlisting}

\begin{center}
\begin{tabularx}{\textwidth}{@{}>{\bfseries\color{anaMavi}}l X X@{}}
\toprule
\rowcolor{acikMavi}\color{anaMavi}\textbf{Kıstas} & \color{anaMavi}\textbf{Segment Ağacı} & \color{anaMavi}\textbf{Fenwick}\\
\midrule
Esneklik & Min/max/gcd/\dots her şey & Genelde sadece toplam\\
\rowcolor{acikGri} Kod uzunluğu & Uzun & Çok kısa\\
Bellek & $\sim 4n$ & $n$\\
\rowcolor{acikGri} Aralık güncelleme & Lazy propagation ile & Sınırlı\\
\bottomrule
\end{tabularx}
\end{center}

\gsec{Ağaçlar --- Alıştırmalar ve Çözümler}

Bu bölümde ağaç veri yapılarının temel ve ileri düzey problemlerini çözeceğiz. İkili ağaçlardan başlayıp arama ağaçları, dengeli ağaçlar, heap yapıları, Trie ve aralık sorgu ağaçlarına (Segment/Fenwick) kadar geniş bir yelpazeyi ele alacağız. Her çözümde zaman ve alan karmaşıklığını açıkça belirteceğiz.

\bilgi{Bölüm boyunca ikili ağaç düğümü için şu standart yapıyı kullanacağız:}

\begin{lstlisting}
struct TreeNode {
    int val;
    TreeNode *left, *right;
    TreeNode(int x) : val(x), left(nullptr), right(nullptr) {}
};
\end{lstlisting}

\gsub{İkili Ağaç ve Dolaşım}

\soru{Bir ikili ağacın yüksekliğini (height) hesaplayan bir fonksiyon yazın. Yükseklik, kökten en derindeki yaprağa giden yoldaki kenar sayısıdır. Boş ağacın yüksekliği $-1$, tek düğümlü ağacın yüksekliği $0$'dır.}

\cozum

Yükseklik, özyinelemeli olarak tanımlanır: bir düğümün yüksekliği, sol ve sağ alt ağaçlarının yüksekliklerinin maksimumundan bir fazladır. Boş ağaç için taban durumu $-1$ döndürürüz; böylece yaprak düğümü otomatik olarak $\max(-1,-1)+1 = 0$ değerini alır.

\begin{lstlisting}
int height(TreeNode* root) {
    if (root == nullptr) return -1;         // taban durumu: boş ağaç
    int leftH  = height(root->left);        // sol alt ağacın yüksekliği
    int rightH = height(root->right);       // sağ alt ağacın yüksekliği
    return 1 + std::max(leftH, rightH);     // en derin taraf + kenar
}
\end{lstlisting}

\notk[Karmaşıklık]{Her düğüm bir kez ziyaret edildiği için zaman \bigO{n}, özyineleme yığını en kötü durumda (dejenere/çarpık ağaç) \bigO{n}, dengeli ağaçta \bigO{\log n} alan kullanır.}

\ipucu{Yükseklik (height) ve derinlik (depth) sık karıştırılır. Derinlik kökten \emph{aşağı} sayılır (kökün derinliği $0$), yükseklik yapraktan \emph{yukarı} sayılır (yaprağın yüksekliği $0$). Bir düğümün derinliğini bulmak için kökten aşağı inerken sayaç arttırılır.}

\soru{Bir ikili ağacın çapını (diameter) bulun. Çap, ağaçtaki herhangi iki düğüm arasındaki en uzun yolun kenar sayısıdır. Bu yolun kökten geçmesi gerekmez.}

\cozum

Herhangi bir düğümden geçen en uzun yol, o düğümün sol alt ağacının yüksekliği ile sağ alt ağacının yüksekliğinin toplamıdır. Çap, tüm düğümler üzerindeki bu toplamın maksimumudur. Yüksekliği hesaplarken aynı gezinme sırasında çapı da güncelleyerek \bigO{n} çözüm elde ederiz; her düğümde ayrı ayrı yükseklik hesaplayan naif yaklaşım \bigO{n^2} olurdu.

\begin{lstlisting}
int diameterHelper(TreeNode* node, int& best) {
    if (node == nullptr) return -1;             // boş ağaç yüksekliği -1
    int leftH  = diameterHelper(node->left,  best);
    int rightH = diameterHelper(node->right, best);
    // bu düğümden geçen yolun kenar sayısı: (leftH+1) + (rightH+1)
    best = std::max(best, leftH + rightH + 2);
    return 1 + std::max(leftH, rightH);         // yüksekliği yukarı taşı
}

int diameter(TreeNode* root) {
    int best = 0;
    diameterHelper(root, best);
    return best;
}
\end{lstlisting}

\notk[Karmaşıklık]{Tek geçişte hem yükseklik hem çap hesaplandığı için zaman \bigO{n}, alan \bigO{h} (ağaç yüksekliği kadar yığın).}

\soru{Yığın (stack) kullanarak özyinelemesiz (iteratif) inorder dolaşımı yazın ve düğüm değerlerini bir vektörde döndürün.}

\cozum

İnorder dolaşımda sıra: sol alt ağaç, kök, sağ alt ağaç. İteratif çözümde bir yığın kullanırız: mevcut düğümden başlayarak sürekli sola inip yol boyunca düğümleri yığına koyarız. Sola gidilemediğinde yığından bir düğüm çıkarıp değerini işleriz ve o düğümün sağ çocuğuna geçeriz. Bu, özyinelemenin yaptığı işi elle simüle eder.

\begin{lstlisting}
std::vector<int> inorderIterative(TreeNode* root) {
    std::vector<int> result;
    std::stack<TreeNode*> st;
    TreeNode* curr = root;
    while (curr != nullptr || !st.empty()) {
        while (curr != nullptr) {   // olabildiğince sola in
            st.push(curr);
            curr = curr->left;
        }
        curr = st.top(); st.pop();  // en soldaki işlenmemiş düğüm
        result.push_back(curr->val);// kökü ziyaret et
        curr = curr->right;         // sağ alt ağaca geç
    }
    return result;
}
\end{lstlisting}

\notk[Karmaşıklık]{Her düğüm yığına bir kez girip bir kez çıktığı için zaman \bigO{n}, yığın en kötü durumda \bigO{h} alan (çarpık ağaçta \bigO{n}).}

\soru{Bir ikili ağacın seviye sıralı (level-order / BFS) dolaşımını, her seviyeyi ayrı bir alt vektör olacak şekilde döndürün.}

\cozum

Seviye sıralı dolaşım için bir kuyruk (queue) kullanırız. Her iterasyonda kuyruğun o anki boyutu, işlenecek seviyedeki düğüm sayısını verir. Bu sayı kadar düğümü çıkarıp değerlerini o seviyenin vektörüne ekler, çocuklarını kuyruğa koyarız. Böylece seviyeleri birbirinden ayırmış oluruz.

\begin{lstlisting}
std::vector<std::vector<int>> levelOrder(TreeNode* root) {
    std::vector<std::vector<int>> levels;
    if (root == nullptr) return levels;
    std::queue<TreeNode*> q;
    q.push(root);
    while (!q.empty()) {
        int sz = q.size();              // bu seviyedeki düğüm sayısı
        std::vector<int> level;
        for (int i = 0; i < sz; ++i) {
            TreeNode* node = q.front(); q.pop();
            level.push_back(node->val);
            if (node->left)  q.push(node->left);
            if (node->right) q.push(node->right);
        }
        levels.push_back(level);
    }
    return levels;
}
\end{lstlisting}

\notk[Karmaşıklık]{Her düğüm bir kez kuyruğa girer/çıkar: zaman \bigO{n}. Kuyruk en geniş seviye kadar düğüm tutabilir, en kötü durumda \bigO{n} alan.}

\soru{İki ikili ağacın yapısal olarak aynı ve düğüm değerlerinin özdeş olup olmadığını kontrol eden bir fonksiyon yazın.}

\cozum

İki ağaç aynıysa: ya ikisi de boştur, ya da ikisinin de kökleri aynı değere sahiptir ve sol alt ağaçları birbiriyle, sağ alt ağaçları birbiriyle aynıdır. Bu tanımı doğrudan özyinelemeye çeviririz.

\begin{lstlisting}
bool isSameTree(TreeNode* p, TreeNode* q) {
    if (p == nullptr && q == nullptr) return true;   // ikisi de boş
    if (p == nullptr || q == nullptr) return false;  // biri boş, diğeri değil
    if (p->val != q->val) return false;              // değerler farklı
    return isSameTree(p->left, q->left)              // sol alt ağaçlar
        && isSameTree(p->right, q->right);           // sağ alt ağaçlar
}
\end{lstlisting}

\notk[Karmaşıklık]{En fazla $\min(n_1, n_2)$ düğüm çifti karşılaştırılır: zaman \bigO{n}, özyineleme yığını \bigO{h} alan.}

\soru{Bir ikili ağacın ayna görüntüsünü alın (invert): her düğümün sol ve sağ çocuklarını yer değiştirin.}

\cozum

Ağacı tersine çevirmek için her düğümde sol ve sağ çocukları takas eder, sonra alt ağaçlar için özyinelemeli olarak aynı işlemi yaparız. Sıra önemli değildir; önce takas edip sonra inebilir ya da önce inip sonra takas edebiliriz.

\begin{lstlisting}
TreeNode* invertTree(TreeNode* root) {
    if (root == nullptr) return nullptr;
    std::swap(root->left, root->right);  // çocukları yer değiştir
    invertTree(root->left);              // yeni sol (eski sağ) alt ağacı çevir
    invertTree(root->right);             // yeni sağ (eski sol) alt ağacı çevir
    return root;
}
\end{lstlisting}

\notk[Karmaşıklık]{Her düğüm bir kez ziyaret edilir: zaman \bigO{n}, alan \bigO{h}.}

\gsub{İkili Arama Ağacı (BST)}

\bilgi{İkili arama ağacında (BST) her düğüm için: sol alt ağaçtaki tüm değerler düğümün değerinden \emph{küçük}, sağ alt ağaçtaki tüm değerler düğümün değerinden \emph{büyüktür}. Bu özellik alt ağaçların \emph{tamamı} için geçerli olmalıdır, sadece doğrudan çocuklar için değil.}

\soru{Verilen bir ikili ağacın geçerli bir BST olup olmadığını doğrulayın.}

\cozum

Yaygın bir hata, yalnızca her düğümü doğrudan çocuklarıyla karşılaştırmaktır; bu yeterli değildir çünkü BST kuralı tüm alt ağaç için geçerli olmalıdır. Doğru yöntem, her düğüme izin verilen bir $(\text{alt sınır}, \text{üst sınır})$ aralığı taşımaktır. Sola inerken üst sınır düğümün değeri olur, sağa inerken alt sınır düğümün değeri olur.

\begin{lstlisting}
bool validate(TreeNode* node, long lo, long hi) {
    if (node == nullptr) return true;
    // düğüm değeri açık aralık (lo, hi) içinde olmalı
    if (node->val <= lo || node->val >= hi) return false;
    return validate(node->left, lo, node->val)     // üst sınır daralır
        && validate(node->right, node->val, hi);   // alt sınır daralır
}

bool isValidBST(TreeNode* root) {
    return validate(root, LONG_MIN, LONG_MAX);
}
\end{lstlisting}

\uyari{Sınırlar için \code{long} kullandık çünkü düğüm değeri \code{INT\_MIN} veya \code{INT\_MAX} olabilir; \code{int} sınırlarla başlarsak taşma ve yanlış karşılaştırma olur.}

\notk[Karmaşıklık]{Her düğüm bir kez ziyaret edilir: zaman \bigO{n}, alan \bigO{h}. Alternatif olarak inorder dolaşımın kesin artan olup olmadığı da kontrol edilebilir.}

\soru{Bir BST'de $k$. en küçük elemanı bulun ($1 \le k \le n$).}

\cozum

BST'nin inorder dolaşımı değerleri artan sırada verir. Dolayısıyla inorder sırasında $k$. ziyaret edilen düğüm cevaptır. Tüm dolaşımı bitirmeden, $k$. düğüme ulaştığımızda durabiliriz. İteratif inorder ile bunu temiz biçimde yaparız.

\begin{lstlisting}
int kthSmallest(TreeNode* root, int k) {
    std::stack<TreeNode*> st;
    TreeNode* curr = root;
    while (curr != nullptr || !st.empty()) {
        while (curr != nullptr) {      // sola in
            st.push(curr);
            curr = curr->left;
        }
        curr = st.top(); st.pop();
        if (--k == 0) return curr->val;// k. ziyaret edilen düğüm
        curr = curr->right;
    }
    return -1; // k geçersizse (buraya normalde ulaşılmaz)
}
\end{lstlisting}

\notk[Karmaşıklık]{En fazla $k$ düğüm işlenerek durulur: zaman \bigO{h + k}, alan \bigO{h}. Sık sorgu gerekiyorsa her düğümde alt ağaç boyutu tutularak \bigO{h} sorgu elde edilebilir.}

\soru{Bir BST'de iki düğümün en küçük ortak atasını (Lowest Common Ancestor, LCA) bulun.}

\cozum

BST özelliği sayesinde LCA'yı kökten inerek buluruz. Kökten başlarız: eğer her iki değer de kökün değerinden büyükse LCA sağ alt ağaçtadır; her ikisi de küçükse sol alt ağaçtadır. Aksi halde (değerler kökün iki yanına düşüyor ya da biri köke eşitse) mevcut düğüm ayrışmanın (split) olduğu ilk noktadır, yani LCA'dır.

\begin{lstlisting}
TreeNode* lowestCommonAncestor(TreeNode* root, TreeNode* p, TreeNode* q) {
    TreeNode* curr = root;
    while (curr != nullptr) {
        if (p->val < curr->val && q->val < curr->val)
            curr = curr->left;        // her ikisi de solda
        else if (p->val > curr->val && q->val > curr->val)
            curr = curr->right;       // her ikisi de sağda
        else
            return curr;              // ayrışma noktası = LCA
    }
    return nullptr;
}
\end{lstlisting}

\notk[Karmaşıklık]{Kökten aşağı tek bir yol izlenir: zaman \bigO{h}, ekstra alan \bigO{1}.}

\gsub{Dengeli Ağaçlar}

\bilgi{Bir ikili ağaç yükseklik-dengeli (height-balanced) ise, \emph{her} düğümde sol ve sağ alt ağaçların yükseklik farkı en fazla $1$'dir. AVL ağaçları bu değişmezi ekleme/silme sırasında dönüşlerle (rotation) korur ve \bigO{\log n} işlem garantisi verir.}

\soru{Bir ikili ağacın yükseklik-dengeli olup olmadığını \bigO{n} zamanda kontrol edin.}

\cozum

Naif yaklaşım her düğümde yükseklik hesaplayıp fark kontrolü yapar ve \bigO{n^2} olur. Bunun yerine \emph{aşağıdan yukarı} (bottom-up) tek geçiş kullanırız: her düğüm için hem yüksekliği hesaplarız hem de alt ağacın dengeli olup olmadığını yukarı taşırız. Bir alt ağaç dengesizse özel bir işaret ($-2$) döndürerek tüm hesabı erkenden dengesiz olarak sonlandırırız.

\begin{lstlisting}
// dengeliyse yüksekliği, değilse -2 döndürür
int checkHeight(TreeNode* node) {
    if (node == nullptr) return -1;              // boş ağaç yüksekliği -1
    int lh = checkHeight(node->left);
    if (lh == -2) return -2;                     // sol alt ağaç dengesiz
    int rh = checkHeight(node->right);
    if (rh == -2) return -2;                     // sağ alt ağaç dengesiz
    if (std::abs(lh - rh) > 1) return -2;        // bu düğümde denge bozuk
    return 1 + std::max(lh, rh);                 // yüksekliği yukarı taşı
}

bool isBalanced(TreeNode* root) {
    return checkHeight(root) != -2;
}
\end{lstlisting}

\ipucu{Erken çıkış (\code{-2} işareti) sayesinde bir dengesizlik bulunduğunda ağacın kalanı taranmaz; yine de en kötü durum karmaşıklığı \bigO{n}'dir.}

\notk[Karmaşıklık]{Her düğüm en fazla bir kez ziyaret edilir: zaman \bigO{n}, özyineleme yığını \bigO{h} alan.}

\gsub{Heap ve Öncelik Kuyruğu}

\bilgi{Öncelik kuyruğu (priority queue) tipik olarak bir ikili heap ile gerçeklenir. \code{std::priority\_queue} varsayılan olarak \emph{maks-heap}'tir (en büyük eleman tepede). Min-heap için \code{std::priority\_queue<int, std::vector<int>, std::greater<int>>} kullanılır. Ekleme ve çıkarma \bigO{\log n}, tepeye bakma \bigO{1}'dir.}

\soru{Bir tam sayı akışında (stream) $k$. en büyük elemanı sürekli olarak takip eden bir yapı tasarlayın. Her yeni eleman eklendiğinde güncel $k$. en büyük değeri döndürün.}

\cozum

Boyutu en fazla $k$ olan bir \emph{min-heap} tutarız. Heap dolu değilse elemanı ekleriz. Heap doluysa ve yeni eleman heap'in tepesinden (şu ana kadarki en küçük "ilk $k$") büyükse, tepeyi çıkarıp yeni elemanı ekleriz. Böylece heap her zaman en büyük $k$ elemanı tutar ve tepesi tam olarak $k$. en büyük elemandır.

\begin{lstlisting}
class KthLargest {
    int k;
    std::priority_queue<int, std::vector<int>, std::greater<int>> minHeap;
public:
    KthLargest(int k, std::vector<int>& nums) : k(k) {
        for (int x : nums) add(x);
    }
    int add(int val) {
        minHeap.push(val);
        if ((int)minHeap.size() > k)
            minHeap.pop();          // en küçüğü at, ilk k büyüğü tut
        return minHeap.top();       // tepe = k. en büyük
    }
};
\end{lstlisting}

\notk[Karmaşıklık]{Her \code{add} çağrısı \bigO{\log k}, alan \bigO{k}. $n$ elemanlık akış için toplam \bigO{n \log k}.}

\soru{Bir tam sayı akışının medyanını her ekleme sonrası \bigO{\log n} zamanda döndüren \code{MedianFinder} sınıfını iki heap ile tasarlayın.}

\cozum

İki heap kullanırız: \code{lo} bir maks-heap (küçük yarının en büyüğü tepede), \code{hi} bir min-heap (büyük yarının en küçüğü tepede). Değişmezi koruruz: \code{lo}'nun boyutu \code{hi}'ye eşit ya da bir fazladır ve \code{lo}'daki her eleman \code{hi}'deki her elemandan küçük veya eşittir. Eleman sayısı tekse medyan \code{lo}'nun tepesidir; çiftse iki tepenin ortalamasıdır.

\begin{lstlisting}
class MedianFinder {
    std::priority_queue<int> lo;                                   // maks-heap
    std::priority_queue<int, std::vector<int>, std::greater<int>> hi; // min-heap
public:
    void addNum(int num) {
        lo.push(num);            // önce küçük yarıya ekle
        hi.push(lo.top());       // lo'nun en büyüğünü hi'ye taşı
        lo.pop();                // (böylece sıralama korunur)
        if (hi.size() > lo.size()) { // boyut dengesini koru
            lo.push(hi.top());
            hi.pop();
        }
    }
    double findMedian() {
        if (lo.size() > hi.size()) return lo.top();      // tek sayıda eleman
        return (lo.top() + hi.top()) / 2.0;              // çift sayıda eleman
    }
};
\end{lstlisting}

\notk[Karmaşıklık]{\code{addNum} \bigO{\log n}, \code{findMedian} \bigO{1}, alan \bigO{n}.}

\soru{$k$ adet artan sıralı bağlı listeyi tek bir sıralı bağlı listede birleştirin. \code{ListNode} yapısını kullanın.}

\cozum

Her listenin başındaki en küçük düğümleri bir min-heap'te tutarız. Heap'ten en küçük düğümü çıkarıp sonuç listesine ekler, o düğümün listedeki bir sonrakini heap'e koyarız. Heap boşalana kadar devam ederiz. Bu, her adımda tüm aktif listelerin başları arasından en küçüğü verimli seçmemizi sağlar.

\begin{lstlisting}
struct ListNode {
    int val;
    ListNode* next;
    ListNode(int x) : val(x), next(nullptr) {}
};

ListNode* mergeKLists(std::vector<ListNode*>& lists) {
    // heap: (değer, düğüm) çiftlerini en küçük değere göre sıralar
    auto cmp = [](ListNode* a, ListNode* b) { return a->val > b->val; };
    std::priority_queue<ListNode*, std::vector<ListNode*>, decltype(cmp)> pq(cmp);
    for (ListNode* node : lists)
        if (node) pq.push(node);      // her listenin başını ekle

    ListNode dummy(0);                // yapay baş düğüm
    ListNode* tail = &dummy;
    while (!pq.empty()) {
        ListNode* smallest = pq.top(); pq.pop();
        tail->next = smallest;        // en küçüğü sonuca ekle
        tail = tail->next;
        if (smallest->next)           // o listeden bir sonrakini ekle
            pq.push(smallest->next);
    }
    return dummy.next;
}
\end{lstlisting}

\notk[Karmaşıklık]{Toplam $N$ düğüm varsa her düğüm heap'e bir kez girip çıkar; heap boyutu en fazla $k$: zaman \bigO{N \log k}, alan \bigO{k}.}

\gsub{Trie}

\bilgi{Trie (önek ağacı), her düğümün bir karakteri temsil ettiği ve kökten bir düğüme giden yolun bir öneki (prefix) kodladığı bir ağaçtır. Kelime ekleme, arama ve önek sorgusu, kelime uzunluğu $L$ olmak üzere \bigO{L} zamanda yapılır --- alfabe boyutundan bağımsızdır.}

\soru{Küçük harflerden oluşan kelimeler için \code{insert}, \code{search} ve \code{startsWith} işlemlerini destekleyen bir Trie yapısı gerçekleyin.}

\cozum

Her Trie düğümü, 26 küçük harf için çocuk işaretçileri ve o düğümde biten bir kelime olup olmadığını belirten bir \code{isEnd} bayrağı tutar. Ekleme sırasında karakter karakter ilerler, yoksa yeni düğüm oluştururuz. \code{search} kelimenin tamamının var olduğunu \emph{ve} son düğümde \code{isEnd} olduğunu ister; \code{startsWith} ise yalnızca yolun var olmasını yeterli görür.

\begin{lstlisting}
class Trie {
    struct Node {
        Node* children[26] = {nullptr};
        bool isEnd = false;
    };
    Node* root;

    Node* findNode(const std::string& s) {
        Node* curr = root;
        for (char c : s) {
            int idx = c - 'a';
            if (!curr->children[idx]) return nullptr; // yol yok
            curr = curr->children[idx];
        }
        return curr;
    }
public:
    Trie() : root(new Node()) {}

    void insert(const std::string& word) {
        Node* curr = root;
        for (char c : word) {
            int idx = c - 'a';
            if (!curr->children[idx])
                curr->children[idx] = new Node(); // eksik düğümü oluştur
            curr = curr->children[idx];
        }
        curr->isEnd = true;                       // kelime sonunu işaretle
    }
    bool search(const std::string& word) {
        Node* node = findNode(word);
        return node != nullptr && node->isEnd;    // tam kelime olmalı
    }
    bool startsWith(const std::string& prefix) {
        return findNode(prefix) != nullptr;       // önek yeterli
    }
};
\end{lstlisting}

\notk[Karmaşıklık]{Her işlem, kelime/önek uzunluğu $L$ için \bigO{L}. Alan, toplam karakter sayısıyla orantılıdır: en kötü durumda \bigO{\text{toplam karakter} \times 26}.}

\soru{\code{addWord} ve \code{search} destekleyen bir \code{WordDictionary} tasarlayın. Aramada \code{.} karakteri herhangi bir tek harfle eşleşebilsin (joker).}

\cozum

Yapı yine bir Trie'dir; fark aramadadır. Normal karakterde tek çocuğa ineriz. \code{.} jokerinde ise mevcut düğümün \emph{tüm} çocukları üzerinden özyinelemeli olarak dener, herhangi biri kalan deseni eşliyorsa başarılı sayarız. Bu, jokerli aramayı bir geri izleme (backtracking) problemine dönüştürür.

\begin{lstlisting}
class WordDictionary {
    struct Node {
        Node* children[26] = {nullptr};
        bool isEnd = false;
    };
    Node* root;

    bool dfs(const std::string& word, int i, Node* node) {
        if (node == nullptr) return false;
        if (i == (int)word.size()) return node->isEnd; // desen bitti
        char c = word[i];
        if (c == '.') {                    // joker: tüm çocukları dene
            for (int j = 0; j < 26; ++j)
                if (dfs(word, i + 1, node->children[j])) return true;
            return false;
        }
        return dfs(word, i + 1, node->children[c - 'a']); // normal karakter
    }
public:
    WordDictionary() : root(new Node()) {}

    void addWord(const std::string& word) {
        Node* curr = root;
        for (char c : word) {
            int idx = c - 'a';
            if (!curr->children[idx]) curr->children[idx] = new Node();
            curr = curr->children[idx];
        }
        curr->isEnd = true;
    }
    bool search(const std::string& word) {
        return dfs(word, 0, root);
    }
};
\end{lstlisting}

\notk[Karmaşıklık]{Ekleme \bigO{L}. Arama, jokersiz durumda \bigO{L}; her joker en kötü durumda 26 dala ayrılabildiği için jokerli durumda üst sınır \bigO{26^{L}} olsa da pratikte çok daha hızlıdır.}

\soru{Bir kelime dizisinin en uzun ortak önekini (longest common prefix) bulun. Ortak önek yoksa boş string döndürün.}

\cozum

Bu problem için tam Trie kurmak şart değildir; dikey tarama (vertical scanning) daha basittir. İlk kelimenin karakterlerini birer birer alır, aynı konumdaki karakterin \emph{tüm} kelimelerde aynı olup olmadığını kontrol ederiz. İlk uyuşmazlıkta veya bir kelime bittiğinde o ana kadar biriken önek cevaptır. (Kavramsal olarak bu, Trie'de dallanma olmayan tek zinciri izlemeye denktir.)

\begin{lstlisting}
std::string longestCommonPrefix(std::vector<std::string>& strs) {
    if (strs.empty()) return "";
    for (int i = 0; i < (int)strs[0].size(); ++i) {
        char c = strs[0][i];                 // referans karakter
        for (int j = 1; j < (int)strs.size(); ++j) {
            // kelime bittiyse ya da karakter uyuşmuyorsa önek burada biter
            if (i >= (int)strs[j].size() || strs[j][i] != c)
                return strs[0].substr(0, i);
        }
    }
    return strs[0];  // ilk kelimenin tamamı ortak önek
}
\end{lstlisting}

\notk[Karmaşıklık]{Tüm karakterlerin toplam sayısı $S$ olmak üzere en kötü durumda zaman \bigO{S}, ekstra alan \bigO{1}.}

\gsub{Segment/Fenwick Ağacı}

\bilgi{Fenwick ağacı (Binary Indexed Tree, BIT) ve segment ağacı, bir dizi üzerinde \emph{nokta güncelleme} ve \emph{aralık sorgusunu} her ikisini de \bigO{\log n} zamanda yapan yapılardır. Fenwick daha az kod ve bellekle önek toplamı gibi tersinir işlemler için idealdir; segment ağacı minimum, maksimum, gcd gibi daha genel birleştirmeleri destekler.}

\soru{Fenwick ağacı (BIT) kullanarak nokta güncelleme ve önek toplamı (dolayısıyla aralık toplamı) sorgusunu destekleyen tam bir yapı yazın. Diziyi $1$-tabanlı indeksleyin.}

\cozum

Fenwick ağacında her indeks, en düşük anlamlı bitine (\code{i \& -i}) bağlı bir aralığın toplamını tutar. \code{update} işleminde indeksi \code{i += i \& -i} ile büyüterek etkilenen tüm bloklara değeri ekleriz. \code{prefixSum} işleminde \code{i -= i \& -i} ile küçülerek $[1, i]$ aralığının toplamını parçalardan toplarız. Aralık toplamı $[l, r]$, \code{prefixSum(r) - prefixSum(l-1)} ile bulunur.

\begin{lstlisting}
class Fenwick {
    std::vector<long> tree;    // 1-tabanlı; tree[0] kullanılmaz
    int n;
public:
    Fenwick(int size) : tree(size + 1, 0), n(size) {}

    // i konumundaki değere delta ekle
    void update(int i, long delta) {
        for (; i <= n; i += i & (-i))   // etkilenen blokları gez
            tree[i] += delta;
    }
    // [1, i] önek toplamı
    long prefixSum(int i) {
        long sum = 0;
        for (; i > 0; i -= i & (-i))    // parça toplamlarını topla
            sum += tree[i];
        return sum;
    }
    // [l, r] aralık toplamı
    long rangeSum(int l, int r) {
        return prefixSum(r) - prefixSum(l - 1);
    }
};
\end{lstlisting}

\uyari{Fenwick ağacı $1$-tabanlıdır; \code{i \& -i} ifadesi $i=0$'da sonsuz döngüye yol açacağından indeksler mutlaka $1$'den başlamalıdır. Dış diziniz $0$-tabanlıysa çağrılarda indekse $1$ ekleyin.}

\notk[Karmaşıklık]{\code{update} ve \code{prefixSum} her ikisi de \bigO{\log n}, alan \bigO{n}. İlk kurulum $n$ güncelleme ile \bigO{n \log n}.}

\soru{Bir segment ağacı ile aralık minimum sorgusu (Range Minimum Query) ve nokta güncellemeyi destekleyen bir yapı yazın.}

\cozum

Segment ağacını bir dizi üzerinde kurarız: her düğüm bir aralığın minimumunu tutar, çocukları o aralığın iki yarısını temsil eder. \code{build} özyinelemeli olarak yaprakları dizideki değerlerle doldurur, iç düğümleri çocuklarının minimumu yapar. \code{query} sorgu aralığıyla düğüm aralığının kesişimine göre üç durumu ele alır: tamamen dışında (etkisiz eleman $+\infty$), tamamen içinde (düğüm değeri), kısmen örtüşme (iki çocuğu birleştir). \code{update} ilgili yaprağı değiştirip kök yolundaki düğümleri günceller.

\begin{lstlisting}
class SegmentTree {
    std::vector<int> tree;     // 4*n boyutunda
    int n;

    void build(const std::vector<int>& a, int node, int lo, int hi) {
        if (lo == hi) { tree[node] = a[lo]; return; }  // yaprak
        int mid = (lo + hi) / 2;
        build(a, 2*node,   lo,    mid);
        build(a, 2*node+1, mid+1, hi);
        tree[node] = std::min(tree[2*node], tree[2*node+1]);
    }
    int query(int node, int lo, int hi, int l, int r) {
        if (r < lo || hi < l) return INT_MAX;          // tamamen dışında
        if (l <= lo && hi <= r) return tree[node];     // tamamen içinde
        int mid = (lo + hi) / 2;                        // kısmi örtüşme
        return std::min(query(2*node,   lo,    mid, l, r),
                        query(2*node+1, mid+1, hi,  l, r));
    }
    void update(int node, int lo, int hi, int pos, int val) {
        if (lo == hi) { tree[node] = val; return; }    // yaprağı değiştir
        int mid = (lo + hi) / 2;
        if (pos <= mid) update(2*node,   lo,    mid, pos, val);
        else            update(2*node+1, mid+1, hi,  pos, val);
        tree[node] = std::min(tree[2*node], tree[2*node+1]); // yukarı güncelle
    }
public:
    SegmentTree(const std::vector<int>& a) : tree(4 * a.size()), n(a.size()) {
        build(a, 1, 0, n - 1);
    }
    int queryMin(int l, int r) { return query(1, 0, n - 1, l, r); }
    void pointUpdate(int pos, int val) { update(1, 0, n - 1, pos, val); }
};
\end{lstlisting}

\notk[Karmaşıklık]{Kurulum \bigO{n}, her sorgu ve her güncelleme \bigO{\log n}, alan \bigO{n} (\code{4n} düğüm ayrılır).}

\gsub{Ağaç Serileştirme}

\bilgi{Serileştirme (serialization), bir ağaç yapısını diske yazılabilecek ya da ağ üzerinden gönderilebilecek düz bir string'e çevirmektir. Deserileştirme, bu string'den özgün ağacı birebir geri kurar. İkili ağaçlar için önsıralı (preorder) dolaşıma boş çocuk işaretçilerini de eklemek en yaygın ve sağlam yöntemdir.}

\soru{Bir ikili ağacı bir string'e çeviren \code{serialize} ve bu string'den ağacı geri kuran \code{deserialize} fonksiyonlarını yazın. Değerler arası ilişkiyi koruyun; boş ağaç da desteklensin.}

\cozum

Önsıralı (kök, sol, sağ) dolaşım kullanırız ve boş çocukları özel bir işaretle (\code{\#}) kodlarız. Boş işaretçilerin kaydedilmesi, ağacın yapısını tek başına preorder'ın belirleyebilmesini sağlar --- bu yüzden ayrıca inorder'a gerek kalmaz. Değerleri boşlukla ayırırız. Geri kurarken token'ları aynı önsıralı düzende okuruz: bir \code{\#} görürsek \code{nullptr} döndürür, aksi halde düğümü oluşturup önce sol sonra sağ alt ağacı özyinelemeli okuruz. Okuma sırası yazma sırasıyla birebir aynı olduğundan ağaç eksiksiz geri kurulur.

\begin{lstlisting}
class Codec {
    void serializeHelper(TreeNode* node, std::ostringstream& out) {
        if (node == nullptr) { out << "# "; return; } // boş işaretçi
        out << node->val << ' ';                       // kök
        serializeHelper(node->left, out);              // sol
        serializeHelper(node->right, out);             // sağ
    }
    TreeNode* deserializeHelper(std::istringstream& in) {
        std::string token;
        in >> token;
        if (token == "#") return nullptr;              // boş alt ağaç
        TreeNode* node = new TreeNode(std::stoi(token));
        node->left  = deserializeHelper(in);           // sol (aynı sıra)
        node->right = deserializeHelper(in);           // sağ
        return node;
    }
public:
    std::string serialize(TreeNode* root) {
        std::ostringstream out;
        serializeHelper(root, out);
        return out.str();
    }
    TreeNode* deserialize(const std::string& data) {
        std::istringstream in(data);
        return deserializeHelper(in);
    }
};
\end{lstlisting}

\ipucu{Preorder + boş işaretçi kombinasyonu ağacı tek başına benzersiz biçimde belirler. Buna karşılık boş işaretçi olmadan tek bir dolaşım (yalnızca preorder veya yalnızca inorder) ağacı belirlemeye yetmez; bu yüzden \code{\#} işaretleri kritiktir.}

\notk[Karmaşıklık]{Hem \code{serialize} hem \code{deserialize} her düğümü (ve boş işaretçileri) bir kez işler: zaman \bigO{n}, üretilen string ve özyineleme yığını \bigO{n} alan kullanır.}

% #####################################################################
\gpart{4}{Graflar}
% #####################################################################

% =====================================================================
\gsec{Graf Gösterimi}
% =====================================================================

Graf, \emph{düğümlerin} (vertex) \emph{kenarlarla} (edge) bağlandığı bir
yapıdır. Ağacın aksine graf döngü içerebilir ve bir düğümden diğerine birden
çok yol olabilir. Sosyal ağlar, haritalar, bağımlılık çizgeleri ve daha
fazlası graflarla modellenir. Kenarlar \emph{yönlü} (directed) veya
\emph{yönsüz} (undirected), \emph{ağırlıklı} (weighted) veya ağırlıksız olabilir.

\gsub{Komşuluk Listesi vs Komşuluk Matrisi}

Bir grafı iki temel biçimde saklarız. Seçim, grafın \emph{yoğunluğuna} bağlıdır:

\begin{center}
\begin{tabularx}{\textwidth}{@{}>{\bfseries\color{anaMavi}}l c c X@{}}
\toprule
\rowcolor{acikMavi}\color{anaMavi}\textbf{Gösterim} & \color{anaMavi}\textbf{Bellek} & \color{anaMavi}\textbf{Kenar var mı?} & \color{anaMavi}\textbf{Ne zaman?}\\
\midrule
Komşuluk listesi & \bigO{V+E} & \bigO{derece} & Seyrek graf (çoğu durum)\\
\rowcolor{acikGri} Komşuluk matrisi & \bigO{V^2} & \bigO{1} & Yoğun graf, hızlı kenar sorgusu\\
\bottomrule
\end{tabularx}
\end{center}

\begin{lstlisting}
int V = 5;                                // dugum sayisi
// Komsuluk listesi (agirliksiz)
vector<vector<int>> adj(V);
adj[0].push_back(1);   adj[1].push_back(0);   // yonsuz kenar 0-1

// Komsuluk listesi (agirlikli): {komsu, agirlik}
vector<vector<pair<int,int>>> g(V);
g[0].push_back({1, 4});  g[1].push_back({0, 4});  // 0-1 arasi agirlik 4

// Komsuluk matrisi
vector<vector<int>> mat(V, vector<int>(V, 0));
mat[0][1] = mat[1][0] = 1;
\end{lstlisting}

% =====================================================================
\gsec{Graf Dolaşımı: BFS ve DFS}
% =====================================================================

Bir grafı ``dolaşmak'', bir başlangıç düğümünden erişilebilen tüm düğümleri
ziyaret etmektir. İki temel strateji vardır ve pek çok graf algoritmasının
temelini oluştururlar.

\gsub{Genişlik-Öncelikli Arama (BFS)}

BFS, başlangıçtan başlayarak düğümleri \emph{katman katman} (mesafeye göre)
gezer; bir kuyruk kullanır. Ağırlıksız graflarda \textbf{en kısa yolu}
(en az kenarla) bulur.

\begin{lstlisting}
vector<int> bfs(const vector<vector<int>>& adj, int start) {
    int V = adj.size();
    vector<int>  dist(V, -1);             // -1 = ziyaret edilmedi
    queue<int>   q;
    dist[start] = 0;
    q.push(start);
    while (!q.empty()) {
        int u = q.front(); q.pop();
        for (int v : adj[u]) {
            if (dist[v] == -1) {           // ilk kez goruluyor
                dist[v] = dist[u] + 1;     // bir katman uzakta
                q.push(v);
            }
        }
    }
    return dist;                           // O(V + E)
}
\end{lstlisting}

\gsub{Derinlik-Öncelikli Arama (DFS)}

DFS, bir yolu \emph{sonuna kadar} takip eder, tıkanınca geri döner (backtrack).
Özyineleme (veya açık bir yığın) ile yapılır. Döngü tespiti, bağlı bileşenler
ve topolojik sıralama DFS üzerine kuruludur.

\begin{lstlisting}
void dfs(const vector<vector<int>>& adj, int u, vector<bool>& visited) {
    visited[u] = true;
    cout << u << " ";
    for (int v : adj[u])
        if (!visited[v])
            dfs(adj, v, visited);          // O(V + E)
}

// Bagli bilesen sayisi (yonsuz graf)
int countComponents(const vector<vector<int>>& adj) {
    int V = adj.size(), count = 0;
    vector<bool> visited(V, false);
    for (int i = 0; i < V; ++i)
        if (!visited[i]) { dfs(adj, i, visited); ++count; }
    return count;
}
\end{lstlisting}

\uyari{Derin graflarda özyinelemeli DFS \emph{yığın taşmasına} (stack overflow)
yol açabilir ($V \sim 10^5$ ve zincir biçimli graf). Böyle durumlarda DFS'i
açık bir \code{std::stack} ile yinelemeli (iterative) yaz.}

% =====================================================================
\gsec{En Kısa Yol Algoritmaları}
% =====================================================================

Ağırlıklı graflarda ``bir düğümden diğerine en düşük toplam ağırlıklı yol
nedir?'' sorusu temeldir. Farklı kısıtlar için farklı algoritmalar kullanılır.

\begin{center}
\begin{tabularx}{\textwidth}{@{}>{\bfseries\color{anaMavi}}l c c X@{}}
\toprule
\rowcolor{acikMavi}\color{anaMavi}\textbf{Algoritma} & \color{anaMavi}\textbf{Karmaşıklık} & \color{anaMavi}\textbf{Negatif kenar?} & \color{anaMavi}\textbf{Kapsam}\\
\midrule
BFS & \bigO{V+E} & --- (ağırlıksız) & Tek kaynak\\
\rowcolor{acikGri} Dijkstra & \bigO{(V+E)\log V} & Hayır & Tek kaynak\\
Bellman-Ford & \bigO{VE} & Evet & Tek kaynak\\
\rowcolor{acikGri} Floyd-Warshall & \bigO{V^3} & Evet & Tüm çiftler\\
\bottomrule
\end{tabularx}
\end{center}

\gsub{Dijkstra Algoritması}

Dijkstra, \emph{negatif olmayan} ağırlıklı graflarda tek bir kaynaktan tüm
düğümlere en kısa yolu bulur. Açgözlü çalışır: her adımda henüz kesinleşmemiş
en yakın düğümü seçer. Bir \emph{min-heap} (öncelik kuyruğu) ile verimlidir.

\begin{lstlisting}
vector<long long> dijkstra(const vector<vector<pair<int,int>>>& g, int source) {
    int V = g.size();
    const long long INF = 1e18;
    vector<long long> dist(V, INF);
    priority_queue<pair<long long,int>,
                   vector<pair<long long,int>>,
                   greater<>> pq;          // min-heap: {mesafe, dugum}
    dist[source] = 0;
    pq.push({0, source});
    while (!pq.empty()) {
        auto [d, u] = pq.top(); pq.pop();
        if (d > dist[u]) continue;         // eski (bayat) kayit, atla
        for (auto [v, w] : g[u]) {
            if (dist[u] + w < dist[v]) {       // gevsetme (relaxation)
                dist[v] = dist[u] + w;
                pq.push({dist[v], v});
            }
        }
    }
    return dist;
}
\end{lstlisting}

\gsub{Bellman-Ford: Negatif Kenarlar}

Bellman-Ford negatif kenarlarla çalışır ve \textbf{negatif döngü} olup
olmadığını da tespit eder. Fikir: tüm kenarları $V-1$ kez gevşet; hâlâ bir
gevşetme mümkünse negatif döngü vardır.

\begin{lstlisting}
struct Edge { int u, v, w; };

bool bellmanFord(int V, const vector<Edge>& edges, int source,
                 vector<long long>& dist) {
    const long long INF = 1e18;
    dist.assign(V, INF);
    dist[source] = 0;
    for (int i = 0; i < V - 1; ++i)        // V-1 tur
        for (auto& e : edges)
            if (dist[e.u] != INF && dist[e.u] + e.w < dist[e.v])
                dist[e.v] = dist[e.u] + e.w;
    // V. tur: hala gevseme varsa negatif dongu
    for (auto& e : edges)
        if (dist[e.u] != INF && dist[e.u] + e.w < dist[e.v])
            return false;                  // negatif dongu var
    return true;
}
\end{lstlisting}

\gsub{Floyd-Warshall: Tüm Çiftler}

Floyd-Warshall, her düğüm çifti arasındaki en kısa yolu tek seferde bulur.
Üç iç içe döngüyle her düğümü ``ara nokta'' olarak dener; kısacık ama güçlü
bir dinamik programlama uygulamasıdır.

\begin{lstlisting}
void floydWarshall(vector<vector<long long>>& d) {   // d: baslangic matris
    int V = d.size();
    for (int k = 0; k < V; ++k)            // ara nokta
        for (int i = 0; i < V; ++i)
            for (int j = 0; j < V; ++j)
                if (d[i][k] + d[k][j] < d[i][j])
                    d[i][j] = d[i][k] + d[k][j];
}
\end{lstlisting}

% =====================================================================
\gsec{Minimum Kapsayan Ağaç (MST)}
% =====================================================================

Ağırlıklı, bağlı ve yönsüz bir grafta, tüm düğümleri en düşük toplam ağırlıkla
birbirine bağlayan (döngüsüz) kenar alt kümesine \emph{minimum kapsayan ağaç}
denir. Ağ tasarımı, kümeleme ve yaklaşık algoritmalar MST kullanır. İki klasik
algoritma vardır; her ikisi de açgözlüdür.

\gsub{Kruskal Algoritması (Union-Find ile)}

Kruskal, kenarları ağırlığa göre sıralar ve döngü oluşturmayan en ucuz
kenarları sırayla ekler. Döngü kontrolü için \emph{Birleşim-Bulma}
(Disjoint Set Union) yapısını kullanır.

\begin{lstlisting}
struct DSU {
    vector<int> parent, rank_;
    DSU(int n) : parent(n), rank_(n, 0) {
        for (int i = 0; i < n; ++i) parent[i] = i;
    }
    int find(int x) {                      // yol sikistirma
        return parent[x] == x ? x : parent[x] = find(parent[x]);
    }
    bool unite(int a, int b) {             // rank ile birlestirme
        a = find(a); b = find(b);
        if (a == b) return false;          // zaten ayni kume -> dongu
        if (rank_[a] < rank_[b]) swap(a, b);
        parent[b] = a;
        if (rank_[a] == rank_[b]) ++rank_[a];
        return true;
    }
};

long long kruskal(int V, vector<tuple<int,int,int>> edges) {
    sort(edges.begin(), edges.end());         // agirliga gore (ilk eleman)
    DSU dsu(V);
    long long sum = 0;
    for (auto [w, u, v] : edges)
        if (dsu.unite(u, v)) sum += w;        // dongu olusturmuyorsa ekle
    return sum;                               // O(E log E)
}
\end{lstlisting}

\gsub{Prim Algoritması}

Prim, tek bir düğümden başlar ve ağacı, her adımda ağaca bağlanan en ucuz
kenarı ekleyerek büyütür. Dijkstra'ya çok benzer; bir min-heap kullanır.

\begin{lstlisting}
long long prim(const vector<vector<pair<int,int>>>& g) {
    int V = g.size();
    vector<bool> inTree(V, false);
    priority_queue<pair<int,int>, vector<pair<int,int>>, greater<>> pq;
    pq.push({0, 0});                       // {agirlik, dugum}
    long long sum = 0;
    int added = 0;
    while (!pq.empty() && added < V) {
        auto [w, u] = pq.top(); pq.pop();
        if (inTree[u]) continue;           // zaten agacta
        inTree[u] = true;
        sum += w;
        ++added;
        for (auto [v, weight] : g[u])
            if (!inTree[v]) pq.push({weight, v});
    }
    return sum;                            // O(E log V)
}
\end{lstlisting}

% =====================================================================
\gsec{Topolojik Sıralama}
% =====================================================================

Yönlü ve döngüsüz bir grafta (DAG), her kenar $u \to v$ için $u$'nun $v$'den
önce geldiği bir düğüm sıralamasına \emph{topolojik sıralama} denir. Görev
zamanlama, derleme bağımlılıkları ve ders ön-koşulları bunu kullanır. İki
yöntem vardır.

\gsub{Kahn Algoritması (BFS Tabanlı)}

Giriş derecesi (in-degree) 0 olan düğümlerden başlarız; her birini sıraya
koyup çıkardıkça komşularının derecesini azaltırız. Bu, aynı zamanda döngü
tespiti de yapar: tüm düğümler işlenemezse graf döngülüdür.

\begin{lstlisting}
vector<int> topoSort(const vector<vector<int>>& adj) {
    int V = adj.size();
    vector<int> inDegree(V, 0), result;
    for (int u = 0; u < V; ++u)
        for (int v : adj[u]) ++inDegree[v];

    queue<int> q;
    for (int u = 0; u < V; ++u)
        if (inDegree[u] == 0) q.push(u);

    while (!q.empty()) {
        int u = q.front(); q.pop();
        result.push_back(u);
        for (int v : adj[u])
            if (--inDegree[v] == 0) q.push(v);
    }
    if ((int)result.size() != V) return {}; // dongu var -> gecersiz
    return result;                          // O(V + E)
}
\end{lstlisting}

\bilgi{Topolojik sıralama \emph{yalnızca} yönlü döngüsüz graflar (DAG) için
tanımlıdır. Sonuç genelde tek değildir: birden fazla geçerli sıralama olabilir.
Kahn algoritmasının döndürdüğü boş vektör, grafta bir döngü olduğunu ---
yani sıralamanın imkânsız olduğunu --- gösterir.}

\gsec{Graflar --- Alıştırmalar ve Çözümler}

\gsub{BFS/DFS Uygulamaları}

Bu bölümde graf gezinmesinin (traversal) klasik uygulamalarını inceliyoruz. Hem \code{grid} (ızgara) üzerinde hem de komşuluk listesi (adjacency list) ile temsil edilen graflar üzerinde çalışacağız. Grid'ler örtük (implicit) graflardır: her hücre bir düğüm, her hücrenin 4 (bazen 8) komşusu ise kenardır.

\soru{Bir \code{grid} verildiğinde (\code{'1'} kara, \code{'0'} deniz), yatay ve dikey olarak birbirine bağlı kara hücrelerinin oluşturduğu ada sayısını bulun. (Number of Islands)}

\cozum

Her kara hücresini bir düğüm gibi düşünürüz. Henüz ziyaret edilmemiş bir kara hücresi bulduğumuzda yeni bir ada başlatır ve DFS/BFS ile o adaya ait tüm bağlı kara hücrelerini gezip işaretleriz. Böylece her ada tam bir kez sayılır. Ziyaret işaretini ayrı bir matris tutmak yerine, gezdiğimiz hücreyi \code{'0'} yaparak yerinde (in-place) yapıyoruz; bu bellek tasarrufu sağlar.

\bilgi{Bağlı bileşen (connected component) sayma mantığı burada gizlidir: her DFS çağrısı bir bileşeni tamamen tüketir, dıştaki döngü kaç kez yeni DFS başlattıysa o kadar bileşen (ada) vardır.}

\begin{lstlisting}
class Solution {
public:
    int numIslands(vector<vector<char>>& grid) {
        int rows = grid.size();
        if (rows == 0) return 0;
        int cols = grid[0].size();
        int count = 0;

        for (int r = 0; r < rows; ++r) {
            for (int c = 0; c < cols; ++c) {
                if (grid[r][c] == '1') {
                    ++count;          // yeni ada bulundu
                    dfs(grid, r, c);  // bu adaya ait her seyi bat
                }
            }
        }
        return count;
    }

private:
    void dfs(vector<vector<char>>& grid, int r, int c) {
        int rows = grid.size(), cols = grid[0].size();
        // sinir disi veya deniz ise dur
        if (r < 0 || r >= rows || c < 0 || c >= cols || grid[r][c] != '1')
            return;

        grid[r][c] = '0';       // ziyaret edildi olarak isaretle (batir)
        dfs(grid, r + 1, c);    // asagi
        dfs(grid, r - 1, c);    // yukari
        dfs(grid, r, c + 1);    // saga
        dfs(grid, r, c - 1);    // sola
    }
};
\end{lstlisting}

Karmaşıklık: her hücre en fazla bir kez ziyaret edilir, dolayısıyla zaman \bigO{R \cdot C}, bellek en kötü durumda (tüm grid kara) özyineleme yığını için \bigO{R \cdot C}.

\uyari{Çok büyük ızgaralarda özyinelemeli DFS yığın taşmasına (stack overflow) yol açabilir. Bu durumda açık bir \code{stack} veya \code{queue} ile iteratif (döngüsel) BFS/DFS tercih edin.}

\soru{Bir grafın kenarları verilmiş; \code{n} düğüm ve \code{edges} listesi var. Grafın kaç bağlı bileşeni (connected component) olduğunu iteratif BFS ile bulun.}

\cozum

Komşuluk listesini kurar, tüm düğümleri dolaşırız. Ziyaret edilmemiş her düğüm yeni bir bileşendir; o düğümden başlayan BFS o bileşenin tamamını ziyaret eder. Bir önceki soruyla aynı fikir, yalnızca temsil ızgara değil komşuluk listesidir.

\begin{lstlisting}
int countComponents(int n, vector<vector<int>>& edges) {
    vector<vector<int>> adj(n);
    for (auto& e : edges) {
        adj[e[0]].push_back(e[1]);
        adj[e[1]].push_back(e[0]);   // yonsuz graf: cift yonlu ekle
    }

    vector<bool> visited(n, false);
    int components = 0;

    for (int start = 0; start < n; ++start) {
        if (visited[start]) continue;
        ++components;                 // yeni bilesen
        queue<int> q;
        q.push(start);
        visited[start] = true;
        while (!q.empty()) {
            int node = q.front(); q.pop();
            for (int next : adj[node]) {
                if (!visited[next]) {
                    visited[next] = true;
                    q.push(next);
                }
            }
        }
    }
    return components;
}
\end{lstlisting}

Karmaşıklık: her düğüm ve kenar bir kez işlenir, zaman \bigO{V + E}, bellek \bigO{V + E} (komşuluk listesi + kuyruk + ziyaret dizisi).

\soru{Yönsüz bir grafın iki parçalı (bipartite) olup olmadığını belirleyin: düğümler iki renge boyanabilir mi, öyle ki her kenar farklı renkli iki düğümü bağlasın?}

\cozum

BFS ile renklendirme yaparız. Başlangıç düğümüne \code{0} rengini veririz; her komşuya karşıt rengi (\code{1 - color}) atarız. Eğer bir komşu zaten boyanmışsa ve rengi mevcut düğümle aynıysa, iki parçalı değildir (aynı renkte bir kenar bulundu). Graf bağlantısız olabileceğinden her bileşen için ayrı başlatırız.

\notk[Sezgi]{İki parçalılık, grafın tek uzunlukta (odd-length) döngü içermemesiyle eşdeğerdir. Renklendirme çelişkisi tam olarak böyle bir tek döngüyü yakalar.}

\begin{lstlisting}
bool isBipartite(vector<vector<int>>& adj) {
    int n = adj.size();
    vector<int> color(n, -1);   // -1: henuz boyanmadi

    for (int start = 0; start < n; ++start) {
        if (color[start] != -1) continue;
        color[start] = 0;
        queue<int> q;
        q.push(start);
        while (!q.empty()) {
            int node = q.front(); q.pop();
            for (int next : adj[node]) {
                if (color[next] == -1) {
                    color[next] = 1 - color[node];  // karsit renk
                    q.push(next);
                } else if (color[next] == color[node]) {
                    return false;   // ayni renkli kenar: iki parcali degil
                }
            }
        }
    }
    return true;
}
\end{lstlisting}

Karmaşıklık: zaman \bigO{V + E}, bellek \bigO{V}.

\soru{Bir \code{grid}'de \code{0} geçilebilir, \code{1} engel. Sol üst köşeden (\code{0,0}) sağ alt köşeye 4 yönlü hareketle en kısa yol uzunluğunu (kaç hücre) bulun; yoksa \code{-1} döndürün.}

\cozum

Kenarların ağırlığı eşit (her adım maliyet 1) olduğundan en kısa yol için BFS yeterlidir. BFS düğümleri uzaklık sırasına göre keşfeder; hedefe ilk ulaştığımız anda bulduğumuz uzaklık en kısadır. \code{dist} matrisi hem ziyaret işareti hem de mesafe tutar.

\ipucu{Ağırlıksız (veya tüm ağırlıkları eşit) graflarda en kısa yol her zaman BFS ile bulunur; Dijkstra'ya gerek yoktur ve BFS daha hızlıdır.}

\begin{lstlisting}
int shortestPath(vector<vector<int>>& grid) {
    int n = grid.size();
    if (n == 0 || grid[0][0] == 1 || grid[n-1][n-1] == 1) return -1;

    vector<vector<int>> dist(n, vector<int>(n, -1));
    queue<pair<int,int>> q;
    q.push({0, 0});
    dist[0][0] = 1;   // baslangic hucresi de sayilir

    int dr[] = {1, -1, 0, 0};
    int dc[] = {0, 0, 1, -1};

    while (!q.empty()) {
        auto [r, c] = q.front(); q.pop();
        if (r == n - 1 && c == n - 1) return dist[r][c];
        for (int k = 0; k < 4; ++k) {
            int nr = r + dr[k], nc = c + dc[k];
            if (nr >= 0 && nr < n && nc >= 0 && nc < n
                && grid[nr][nc] == 0 && dist[nr][nc] == -1) {
                dist[nr][nc] = dist[r][c] + 1;
                q.push({nr, nc});
            }
        }
    }
    return -1;   // hedefe ulasilamadi
}
\end{lstlisting}

Karmaşıklık: zaman \bigO{R \cdot C}, bellek \bigO{R \cdot C}.

\gsub{Döngü Tespiti ve Topolojik Sıralama}

Yönlü asiklik graflar (DAG --- Directed Acyclic Graph) bağımlılık modellemenin temelidir: dersler, derleme sırası, iş zamanlaması. Topolojik sıralama, her kenar \code{u -> v} için \code{u}'nun \code{v}'den önce geldiği bir düğüm dizilimidir. Böyle bir sıralama ancak ve ancak grafta yönlü döngü yoksa vardır.

\soru{Yönlü bir grafta döngü olup olmadığını DFS ile üç-renk yöntemiyle tespit edin.}

\cozum

Her düğümü üç durumdan biriyle işaretleriz: \code{0} beyaz (hiç ziyaret edilmedi), \code{1} gri (işleniyor, özyineleme yığınında), \code{2} siyah (tamamen bitti). DFS sırasında gri bir düğüme geri kenar (back edge) bulursak döngü vardır. Yalnızca siyah bir düğüme rastlamak döngü değildir; o dal zaten güvenle tamamlanmıştır.

\uyari{Yönlü graflarda döngü tespiti için tek bir \code{visited} bayrağı yetmez; düğümün hâlâ işlenmekte olduğunu (gri) ayrıca ayırt etmelisiniz. Yönsüz graflarda ise ebeveyni (parent) hariç tutmak gerekir.}

\begin{lstlisting}
class CycleDetector {
public:
    bool hasCycle(int n, vector<vector<int>>& adj) {
        vector<int> state(n, 0);   // 0 beyaz, 1 gri, 2 siyah
        for (int i = 0; i < n; ++i)
            if (state[i] == 0 && dfs(i, adj, state))
                return true;
        return false;
    }

private:
    bool dfs(int node, vector<vector<int>>& adj, vector<int>& state) {
        state[node] = 1;   // gri: yigina girdi
        for (int next : adj[node]) {
            if (state[next] == 1) return true;          // geri kenar: dongu
            if (state[next] == 0 && dfs(next, adj, state)) return true;
        }
        state[node] = 2;   // siyah: tamamlandi
        return false;
    }
};
\end{lstlisting}

Karmaşıklık: zaman \bigO{V + E}, bellek \bigO{V} (durum dizisi + özyineleme yığını).

\soru{\code{numCourses} ders ve \code{prerequisites} önkoşul çiftleri (\code{[a, b]}: \code{a} için önce \code{b} alınmalı) verildi. Tüm dersleri tamamlamak mümkün müdür? (Course Schedule)}

\cozum

Bu, ders bağımlılık grafında yönlü döngü olup olmadığı sorusudur: döngü varsa (\code{A}, \code{B}'yi \code{B} de \code{A}'yı gerektirir gibi) tamamlamak imkânsızdır. Kahn algoritmasını (BFS tabanlı topolojik sıralama) kullanıyoruz: giriş derecesi (in-degree) sıfır olan düğümlerden başlarız, her işlediğimizde komşuların giriş derecesini azaltırız. Sonunda işlenen düğüm sayısı toplamdan azsa graf bir döngü içerir.

\bilgi{Kahn algoritmasının güzel yanı, aynı anda hem döngü tespiti yapması hem de döngü yoksa geçerli bir topolojik sıra üretmesidir.}

\begin{lstlisting}
bool canFinish(int numCourses, vector<vector<int>>& prerequisites) {
    vector<vector<int>> adj(numCourses);
    vector<int> indegree(numCourses, 0);

    for (auto& p : prerequisites) {
        // b -> a: once b, sonra a
        adj[p[1]].push_back(p[0]);
        ++indegree[p[0]];
    }

    queue<int> q;
    for (int i = 0; i < numCourses; ++i)
        if (indegree[i] == 0) q.push(i);   // onkosulsuz dersler

    int processed = 0;
    while (!q.empty()) {
        int node = q.front(); q.pop();
        ++processed;
        for (int next : adj[node]) {
            if (--indegree[next] == 0)   // son onkosul da karsilandi
                q.push(next);
        }
    }
    return processed == numCourses;   // hepsi islendiyse dongu yok
}
\end{lstlisting}

Karmaşıklık: zaman \bigO{V + E}, bellek \bigO{V + E}.

\soru{Bir DAG için geçerli bir topolojik sıralama üretin; graf döngü içeriyorsa boş bir liste döndürün.}

\cozum

Yine Kahn algoritmasını kullanırız, ancak bu kez işlenen düğümleri sırayla bir sonuç vektörüne ekleriz. Kuyruktan çıkardığımız sıra, geçerli bir topolojik dizilim verir. Sonunda listenin boyutu düğüm sayısından küçükse döngü vardır ve boş liste döndürürüz.

\begin{lstlisting}
vector<int> topologicalSort(int n, vector<vector<int>>& adj) {
    vector<int> indegree(n, 0);
    for (int u = 0; u < n; ++u)
        for (int v : adj[u])
            ++indegree[v];

    queue<int> q;
    for (int i = 0; i < n; ++i)
        if (indegree[i] == 0) q.push(i);

    vector<int> order;
    while (!q.empty()) {
        int node = q.front(); q.pop();
        order.push_back(node);        // sonuca ekle
        for (int next : adj[node])
            if (--indegree[next] == 0)
                q.push(next);
    }

    if ((int)order.size() != n)
        return {};   // dongu var: gecerli siralama yok
    return order;
}
\end{lstlisting}

Karmaşıklık: zaman \bigO{V + E}, bellek \bigO{V}.

\notk[Alternatif]{Topolojik sıra DFS ile de üretilebilir: her düğümün DFS'i bittiğinde onu bir yığına itersiniz; sonunda yığını ters çevirmek topolojik sırayı verir. Kahn yöntemi ise döngüyü daha doğrudan raporlar.}

\gsub{En Kısa Yol}

Ağırlıklı graflarda en kısa yol için BFS yetmez; kenar maliyetleri farklı olduğunda Dijkstra algoritması gerekir. Ağırlıklar negatif olmadığı sürece Dijkstra en kısa yolları doğru bulur.

\soru{Negatif olmayan ağırlıklı yönlü bir grafta, \code{source} düğümünden diğer tüm düğümlere en kısa mesafeleri Dijkstra ile bulun.}

\cozum

Dijkstra, açgözlü (greedy) bir yaklaşımla her adımda henüz kesinleşmemiş en küçük mesafeli düğümü seçer ve komşularını gevşetir (relax). Bu en küçük seçimi verimli yapmak için minimum öbeği (\code{priority\_queue}) kullanırız. Bir düğümü ilk kez kuyruktan minimum mesafeyle çıkardığımızda o mesafe kesinleşmiştir. Eski (bayat) girişleri, saklanan mesafe güncel \code{dist} değerinden büyükse atlayarak eleriz.

\bilgi{\code{priority\_queue} varsayılan olarak maksimum öbektir. Minimum öbek için \code{greater<>} karşılaştırıcısı kullanılır: \code{priority\_queue<T, vector<T>, greater<T>>}.}

\begin{lstlisting}
vector<long long> dijkstra(int n, vector<vector<pair<int,int>>>& adj, int source) {
    const long long INF = LLONG_MAX;
    vector<long long> dist(n, INF);
    dist[source] = 0;

    // (mesafe, dugum) ciftleri; en kucuk mesafe tepede
    using State = pair<long long, int>;
    priority_queue<State, vector<State>, greater<State>> pq;
    pq.push({0, source});

    while (!pq.empty()) {
        auto [d, u] = pq.top(); pq.pop();
        if (d > dist[u]) continue;   // bayat giris: atla

        for (auto& [v, w] : adj[u]) {
            long long nd = d + w;
            if (nd < dist[v]) {      // gevsetme (relaxation)
                dist[v] = nd;
                pq.push({nd, v});
            }
        }
    }
    return dist;
}
\end{lstlisting}

Karmaşıklık: öbek tabanlı Dijkstra ile zaman \bigO{(V + E) \log V}, bellek \bigO{V + E}. Negatif kenar varsa Dijkstra yanlış sonuç verir; o durumda Bellman-Ford (\bigO{V \cdot E}) gerekir.

\soru{\code{beginWord}, \code{endWord} ve bir \code{wordList} verildi. Her adımda tek bir harf değiştirerek (ve her ara kelime listede olmalı) \code{beginWord}'den \code{endWord}'e ulaşan en kısa dönüşüm dizisinin uzunluğunu bulun. (Word Ladder)}

\cozum

Kelimeleri düğüm, tek harf farkla ayrışan kelime çiftlerini kenar kabul eden örtük bir graf kurarız. Tüm kenarların maliyeti 1 olduğundan en kısa yol BFS ile bulunur. Her kelimenin komşularını bulmak için her konumdaki harfi \code{'a'}--\code{'z'} ile değiştirir, oluşan kelime listede (\code{set}) varsa komşu sayarız. Ziyaret edilen kelimeleri kümeden silerek tekrar gezmeyi önleriz.

\begin{lstlisting}
int ladderLength(string beginWord, string endWord, vector<string>& wordList) {
    unordered_set<string> dict(wordList.begin(), wordList.end());
    if (dict.find(endWord) == dict.end()) return 0;

    queue<string> q;
    q.push(beginWord);
    int steps = 1;   // baslangic kelimesi de sayilir

    while (!q.empty()) {
        int levelSize = q.size();
        for (int i = 0; i < levelSize; ++i) {   // bir seviye (uzaklik) bat
            string word = q.front(); q.pop();
            if (word == endWord) return steps;

            for (int pos = 0; pos < (int)word.size(); ++pos) {
                char original = word[pos];
                for (char c = 'a'; c <= 'z'; ++c) {
                    word[pos] = c;
                    if (dict.find(word) != dict.end()) {
                        dict.erase(word);   // ziyaret edildi: tekrar bakma
                        q.push(word);
                    }
                }
                word[pos] = original;   // kelimeyi geri yukle
            }
        }
        ++steps;
    }
    return 0;   // ulasilamadi
}
\end{lstlisting}

Karmaşıklık: \code{N} kelime ve \code{L} kelime uzunluğu için zaman \bigO{N \cdot L^2} (her kelime için \code{L} konum \times 26 deneme, her deneme \code{L} maliyetli karma), bellek \bigO{N \cdot L}.

\soru{Kenar ağırlıkları yalnızca \code{0} veya \code{1} olan bir grafta \code{source}'tan en kısa yolu, \code{priority\_queue} yerine \code{deque} ile \bigO{V+E} zamanda bulun. (0-1 BFS) Yaklaşımı açıklayın.}

\cozum

Tüm ağırlıklar 0 veya 1 olduğunda Dijkstra'nın öbeğine gerek kalmaz. Çift uçlu kuyruk (\code{deque}) kullanırız: bir kenarı gevşetirken ağırlık \code{0} ise düğümü kuyruğun \emph{önüne}, \code{1} ise \emph{arkasına} ekleriz. Böylece kuyruk her zaman mesafeye göre (en fazla iki farklı değer) sıralı kalır ve öbeğin \bigO{\log V} çarpanından kurtuluruz.

\ipucu{0-1 BFS mantığı, ağırlıkların \code{0..k} aralığında olduğu durumlara "dial's algorithm" ile genelleşir. Ancak temel fikir aynıdır: küçük ağırlığı öne, büyüğü arkaya koy.}

\begin{lstlisting}
vector<int> zeroOneBFS(int n, vector<vector<pair<int,int>>>& adj, int source) {
    const int INF = INT_MAX;
    vector<int> dist(n, INF);
    dist[source] = 0;

    deque<int> dq;
    dq.push_front(source);

    while (!dq.empty()) {
        int u = dq.front(); dq.pop_front();
        for (auto& [v, w] : adj[u]) {   // w yalnizca 0 veya 1
            if (dist[u] + w < dist[v]) {
                dist[v] = dist[u] + w;
                if (w == 0) dq.push_front(v);   // 0 kenar: one
                else        dq.push_back(v);    // 1 kenar: arkaya
            }
        }
    }
    return dist;
}
\end{lstlisting}

Karmaşıklık: zaman \bigO{V + E}, bellek \bigO{V}. Klasik Dijkstra'ya göre \code{log} çarpanı yoktur.

\gsub{Birleşim-Bulma (DSU)}

Birleşim-Bulma (Union-Find / Disjoint Set Union --- DSU) veri yapısı, ayrık kümeler üzerinde iki işlemi neredeyse sabit zamanda yapar: \code{find} (bir elemanın hangi kümede olduğunu bul) ve \code{union} (iki kümeyi birleştir). İki optimizasyon kritiktir: yol sıkıştırma (path compression) ve rütbeye/boyuta göre birleştirme (union by rank/size). İkisi birlikte işlem başına ters-Ackermann \bigO{\alpha(n)} verir; pratikte sabit sayılır.

\notk[Temel Yapı]{Aşağıdaki üç problemde de aynı DSU sınıfını kullanacağız. Önce sınıfı bir kez tanımlayıp sonra uygulamalarına bakalım.}

\begin{lstlisting}
class DSU {
public:
    DSU(int n) : parent(n), rank_(n, 0), count(n) {
        for (int i = 0; i < n; ++i) parent[i] = i;   // her eleman kendi kokunu gosterir
    }

    int find(int x) {
        if (parent[x] != x)
            parent[x] = find(parent[x]);   // yol sikistirma
        return parent[x];
    }

    bool unite(int a, int b) {
        int ra = find(a), rb = find(b);
        if (ra == rb) return false;        // zaten ayni kumede
        if (rank_[ra] < rank_[rb]) swap(ra, rb);
        parent[rb] = ra;                   // rutbeye gore birlestir
        if (rank_[ra] == rank_[rb]) ++rank_[ra];
        --count;                           // kume sayisi bir azaldi
        return true;
    }

    int componentCount() const { return count; }

private:
    vector<int> parent;
    vector<int> rank_;
    int count;   // ayrik kume sayisi
};
\end{lstlisting}

\soru{\code{n} düğüm ve \code{edges} verildi. Bağlı bileşen sayısını DSU ile bulun.}

\cozum

Başlangıçta her düğüm kendi kümesindedir, yani \code{n} bileşen vardır. Her kenar için iki uç düğümü \code{unite} ederiz; farklı kümelerdeyseler birleşir ve sayaç azalır. Tüm kenarlar işlendiğinde \code{componentCount} bağlı bileşen sayısını verir.

\begin{lstlisting}
int countComponentsDSU(int n, vector<vector<int>>& edges) {
    DSU dsu(n);
    for (auto& e : edges)
        dsu.unite(e[0], e[1]);
    return dsu.componentCount();
}
\end{lstlisting}

Karmaşıklık: \code{E} kenar için zaman \bigO{E \cdot \alpha(n)}, bellek \bigO{n}.

\soru{Yönsüz bir grafta döngü olup olmadığını DSU ile tespit edin.}

\cozum

Kenarları tek tek işleriz. Bir kenarın iki ucu \code{unite} edilmeden önce zaten aynı kümedeyse (\code{find} kökleri eşitse), bu kenar zaten bağlı iki düğümü yeniden birleştiriyor demektir; yani bir döngü oluşur. \code{unite} bu durumda \code{false} döndürdüğü için tespiti tek satırda yaparız.

\uyari{Bu yöntem yönsüz graflar içindir. Yönlü graflarda döngü tespiti için DSU uygun değildir; oradaki üç-renk DFS veya Kahn yöntemini kullanın.}

\begin{lstlisting}
bool hasCycleUndirected(int n, vector<vector<int>>& edges) {
    DSU dsu(n);
    for (auto& e : edges) {
        if (!dsu.unite(e[0], e[1]))
            return true;   // uclar zaten baglı: dongu
    }
    return false;
}
\end{lstlisting}

Karmaşıklık: zaman \bigO{E \cdot \alpha(n)}, bellek \bigO{n}.

\soru{\code{accounts} listesi verildi; her eleman \code{[isim, email1, email2, ...]} biçimindedir. Ortak en az bir e-postayı paylaşan hesaplar aynı kişiye aittir (isimler aynıdır). Hesapları birleştirip her kişinin sıralı e-posta listesini üretin. (Accounts Merge)}

\cozum

Her hesabı bir düğüm sayarız. Aynı e-postayı içeren hesapları DSU ile birleştiririz: her e-postayı ilk gördüğümüz hesaba bağlarız (\code{emailToAccount} eşlemesi); aynı e-posta başka bir hesapta tekrar görünürse iki hesabı \code{unite} ederiz. Sonra her e-postayı ait olduğu hesabın kökü altında toplar, her kök için e-postaları sıralayıp isimle birlikte çıktı veririz.

\bilgi{DSU'nun gücü tam burada görülür: "geçişli birleşme" (A ile B ortak posta, B ile C ortak posta $\Rightarrow$ A, B, C aynı kişi) doğal olarak ele alınır; DSU geçişliliği (transitivity) ücretsiz sağlar.}

\begin{lstlisting}
vector<vector<string>> accountsMerge(vector<vector<string>>& accounts) {
    int n = accounts.size();
    DSU dsu(n);
    unordered_map<string, int> emailToAccount;   // email -> ilk hesap indeksi

    for (int i = 0; i < n; ++i) {
        for (int j = 1; j < (int)accounts[i].size(); ++j) {
            const string& email = accounts[i][j];
            auto it = emailToAccount.find(email);
            if (it == emailToAccount.end())
                emailToAccount[email] = i;
            else
                dsu.unite(i, it->second);   // ayni email: hesaplari birlestir
        }
    }

    // her kok icin e-postalari topla (siralilik icin set kullan)
    unordered_map<int, set<string>> merged;
    for (auto& [email, idx] : emailToAccount)
        merged[dsu.find(idx)].insert(email);

    vector<vector<string>> result;
    for (auto& [root, emails] : merged) {
        vector<string> entry;
        entry.push_back(accounts[root][0]);   // isim
        entry.insert(entry.end(), emails.begin(), emails.end());
        result.push_back(entry);
    }
    return result;
}
\end{lstlisting}

Karmaşıklık: toplam e-posta sayısı \code{K} için zaman \bigO{K \log K} (sıralama baskın), bellek \bigO{K}.

\gsub{Güçlü Bağlı Bileşenler (SCC)}

Yönlü bir grafta güçlü bağlı bileşen (Strongly Connected Component --- SCC), içindeki her düğüm çiftinin karşılıklı olarak birbirine ulaşabildiği maksimal düğüm kümesidir. SCC'leri tek bir düğüme büzersek, geriye her zaman bir DAG (yönlü asiklik graf) kalır; bu "yoğunlaştırma" (condensation) birçok problemi basitleştirir.

Kosaraju algoritması iki DFS geçişiyle SCC'leri bulur:

\begin{tabularx}{\textwidth}{@{}lX@{}}
Adım 1 & Orijinal grafta DFS yap; her düğümü \emph{bitiş sırasına} göre bir yığına it. \\
Adım 2 & Grafın transpozunu (tüm kenarları ters çevir) oluştur. \\
Adım 3 & Yığından düğümleri çıkararak transpoz grafta DFS yap; her yeni DFS bir SCC verir. \\
\end{tabularx}

\notk[Neden Çalışır?]{İlk DFS'in bitiş sıraları, SCC'leri yoğunlaştırma DAG'ında topolojik sıraya dizer. Kenarları ters çevirdiğimizde, en geç biten SCC'den başlayan DFS o SCC'nin dışına çıkamaz --- çünkü ters grafta o SCC'ye giden kenarlar artık başka SCC'lere gitmez. Böylece her DFS tam bir SCC toplar.}

\soru{Yönlü bir graf verildiğinde tüm güçlü bağlı bileşenleri Kosaraju algoritmasıyla bulun; her SCC'yi ayrı bir düğüm listesi olarak döndürün.}

\cozum

Üç adımı doğrudan uygularız. İlk DFS bitiş sırasına göre bir \code{order} yığını doldurur. Transpoz grafı \code{adjT} olarak kurarız. Son olarak \code{order}'ı tersten (en geç biten önce) gezerek transpoz grafta DFS yaparız; her başlatılan DFS bir SCC toplar.

\begin{lstlisting}
class Kosaraju {
public:
    vector<vector<int>> findSCCs(int n, vector<vector<int>>& adj) {
        vector<bool> visited(n, false);
        vector<int> order;   // bitis sirasi yigini

        // 1. adim: orijinal grafta DFS, bitis sirasini biriktir
        for (int i = 0; i < n; ++i)
            if (!visited[i])
                dfs1(i, adj, visited, order);

        // 2. adim: transpoz grafi kur
        vector<vector<int>> adjT(n);
        for (int u = 0; u < n; ++u)
            for (int v : adj[u])
                adjT[v].push_back(u);   // kenari ters cevir

        // 3. adim: bitis sirasinin tersinden transpozda DFS
        fill(visited.begin(), visited.end(), false);
        vector<vector<int>> sccs;
        for (int i = (int)order.size() - 1; i >= 0; --i) {
            int node = order[i];
            if (!visited[node]) {
                vector<int> component;
                dfs2(node, adjT, visited, component);
                sccs.push_back(component);
            }
        }
        return sccs;
    }

private:
    void dfs1(int u, vector<vector<int>>& adj,
              vector<bool>& visited, vector<int>& order) {
        visited[u] = true;
        for (int v : adj[u])
            if (!visited[v])
                dfs1(v, adj, visited, order);
        order.push_back(u);   // dugum bitti: yigina it
    }

    void dfs2(int u, vector<vector<int>>& adjT,
              vector<bool>& visited, vector<int>& component) {
        visited[u] = true;
        component.push_back(u);
        for (int v : adjT[u])
            if (!visited[v])
                dfs2(v, adjT, visited, component);
    }
};
\end{lstlisting}

Karmaşıklık: iki DFS ve bir transpoz kurulumu, her biri \bigO{V + E}; toplam zaman \bigO{V + E}, bellek \bigO{V + E} (transpoz graf + yığın).

\bilgi{Kosaraju'ya alternatif Tarjan algoritmasıdır; o tek DFS geçişiyle (low-link değerleri ile) SCC bulur ve transpoz gerektirmez. Kosaraju kavramsal olarak daha anlaşılır, Tarjan pratikte biraz daha hızlıdır. İkisi de \bigO{V + E}'dir.}

% #####################################################################
\gpart{5}{Algoritma Teknikleri}
% #####################################################################

% =====================================================================
\gsec{Sıralama Algoritmaları}
% =====================================================================

Sıralama, verilerin belirli bir düzene sokulmasıdır ve pek çok başka
algoritmanın (ikili arama, birçok açgözlü yöntem) ön koşuludur. Farklı sıralama
algoritmaları farklı ödünleşmeler sunar. Anahtar kavramlar: \emph{kararlılık}
(stable --- eşit elemanların göreli sırası korunur mu?) ve \emph{yerinde}
(in-place --- ek bellek gerekir mi?).

\begin{center}
\begin{tabularx}{\textwidth}{@{}>{\bfseries\color{anaMavi}}l c c c c@{}}
\toprule
\rowcolor{acikMavi}\color{anaMavi}\textbf{Algoritma} & \color{anaMavi}\textbf{Ort.} & \color{anaMavi}\textbf{En kötü} & \color{anaMavi}\textbf{Bellek} & \color{anaMavi}\textbf{Kararlı?}\\
\midrule
Bubble & \bigO{n^2} & \bigO{n^2} & \bigO{1} & Evet\\
\rowcolor{acikGri} Selection & \bigO{n^2} & \bigO{n^2} & \bigO{1} & Hayır\\
Insertion & \bigO{n^2} & \bigO{n^2} & \bigO{1} & Evet\\
\rowcolor{acikGri} Merge & \bigO{n\log n} & \bigO{n\log n} & \bigO{n} & Evet\\
Quick & \bigO{n\log n} & \bigO{n^2} & \bigO{\log n} & Hayır\\
\rowcolor{acikGri} Heap & \bigO{n\log n} & \bigO{n\log n} & \bigO{1} & Hayır\\
Counting & \bigO{n+k} & \bigO{n+k} & \bigO{k} & Evet\\
\bottomrule
\end{tabularx}
\end{center}

\gsub{Basit Sıralamalar: Insertion Sort}

Ekleme sıralaması, kart oyununda elini sıralaman gibi çalışır: her elemanı
alıp solundaki sıralı kısımda doğru yerine yerleştirir. Küçük veya neredeyse
sıralı diziler için pratikte hızlıdır.

\begin{lstlisting}
void insertionSort(vector<int>& a) {
    for (int i = 1; i < (int)a.size(); ++i) {
        int key = a[i], j = i - 1;
        while (j >= 0 && a[j] > key) {     // buyukleri saga kaydir
            a[j + 1] = a[j];
            --j;
        }
        a[j + 1] = key;                    // dogru yere koy
    }
}
\end{lstlisting}

\gsub{Merge Sort: Böl ve Birleştir}

Birleştirme sıralaması bir \emph{böl-yönet} algoritmasıdır: diziyi ikiye böler,
her yarıyı özyinelemeli sıralar, sonra iki sıralı yarıyı birleştirir.
\bigO{n\log n} garantisi verir ve kararlıdır.

\begin{lstlisting}
void merge(vector<int>& a, int left, int mid, int right) {
    vector<int> temp;
    int i = left, j = mid + 1;
    while (i <= mid && j <= right)         // iki yariyi sirali birlestir
        temp.push_back(a[i] <= a[j] ? a[i++] : a[j++]);
    while (i <= mid) temp.push_back(a[i++]);
    while (j <= right)  temp.push_back(a[j++]);
    for (int k = 0; k < (int)temp.size(); ++k)
        a[left + k] = temp[k];
}

void mergeSort(vector<int>& a, int left, int right) {
    if (left >= right) return;
    int mid = left + (right - left) / 2;
    mergeSort(a, left, mid);               // sol yariyi sirala
    mergeSort(a, mid + 1, right);          // sag yariyi sirala
    merge(a, left, mid, right);            // birlestir
}
\end{lstlisting}

\gsub{Quick Sort: Bölümleme (Partition)}

Hızlı sıralama bir \emph{pivot} seçer ve diziyi ``pivottan küçükler'' ve
``büyükler'' diye ikiye böler; sonra her parçayı özyinelemeli sıralar.
Ortalama \bigO{n\log n}'dir ama kötü pivot seçiminde \bigO{n^2}'ye düşebilir.

\begin{lstlisting}
int partition(vector<int>& a, int left, int right) {
    int pivot = a[right];                  // son elemani pivot sec
    int i = left - 1;
    for (int j = left; j < right; ++j)
        if (a[j] < pivot) swap(a[++i], a[j]);
    swap(a[i + 1], a[right]);              // pivotu ortasina koy
    return i + 1;
}

void quickSort(vector<int>& a, int left, int right) {
    if (left >= right) return;
    int p = partition(a, left, right);
    quickSort(a, left, p - 1);             // sol parca
    quickSort(a, p + 1, right);            // sag parca
}
\end{lstlisting}

\gsub{Counting Sort: Karşılaştırmasız Sıralama}

Elemanlar küçük bir tam sayı aralığındaysa, karşılaştırma yapmadan sayarak
\bigO{n+k}'de sıralarız. Bu, karşılaştırma sıralamalarının \bigO{n\log n} alt
sınırını aşar çünkü karşılaştırma kullanmaz.

\begin{lstlisting}
void countingSort(vector<int>& a, int maxValue) {
    vector<int> count(maxValue + 1, 0);
    for (int x : a) ++count[x];            // her degeri say
    int idx = 0;
    for (int v = 0; v <= maxValue; ++v)    // sirayla geri yaz
        while (count[v]-- > 0) a[idx++] = v;
}
\end{lstlisting}

\ipucu{Gerçek kodda kendi sıralamanı yazma: \code{std::sort} (introsort ---
quicksort+heapsort+insertion karışımı) \bigO{n\log n} garantisi verir ve son
derece optimize edilmiştir. Kararlılık gerekiyorsa \code{std::stable\_sort}
kullan. Kendi karşılaştırıcını lambda ile verebilirsin:
\code{sort(v.begin(), v.end(), [](auto\& a, auto\& b)\{ return a > b; \});}}

% =====================================================================
\gsec{Arama ve İkili Arama}
% =====================================================================

\gsub{Doğrusal ve İkili Arama}

Sırasız veride tek yol doğrusal aramadır (\bigO{n}). Ama veri \emph{sıralıysa},
ikili arama her adımda arama uzayını yarıya indirerek \bigO{\log n}'de bulur.

\begin{lstlisting}
int binarySearch(const vector<int>& a, int target) {
    int left = 0, right = a.size() - 1;
    while (left <= right) {
        int mid = left + (right - left) / 2;  // tasmayi onleyen orta
        if (a[mid] == target)     return mid;
        else if (a[mid] < target) left = mid + 1;
        else                      right = mid - 1;
    }
    return -1;                             // bulunamadi
}
\end{lstlisting}

\gsub{Cevap Üzerinde İkili Arama (Binary Search on Answer)}

İkili aramanın güçlü bir genellemesi: cevabın \emph{monoton} bir yes/no
koşulunu sağladığı problemlerde, cevabın kendisini arayabiliriz. Örnek:
``$k$ gün içinde bitirmek için gereken minimum günlük kapasite nedir?''

\begin{lstlisting}
// feasible(x): x kapasite ile is k gunde bitiyor mu? (monoton: buyudukce true)
int smallestValid(int lo, int hi, function<bool(int)> feasible) {
    while (lo < hi) {
        int mid = lo + (hi - lo) / 2;
        if (feasible(mid)) hi = mid;       // orta yeterli -> daha kucuk dene
        else               lo = mid + 1;   // yetersiz -> daha buyuk gerek
    }
    return lo;                             // ilk 'true' olan deger
}
\end{lstlisting}

% =====================================================================
\gsec{Böl ve Yönet (Divide and Conquer)}
% =====================================================================

Böl-yönet stratejisi bir problemi (1) daha küçük alt problemlere \emph{böler},
(2) her birini özyinelemeli \emph{çözer}, (3) sonuçları \emph{birleştirir}.
Merge sort ve quick sort bunun örnekleridir. Çalışma süresi genelde
\emph{Master teoremi} ile analiz edilir.

\bilgi{\textbf{Master Teoremi:} $T(n) = a\,T(n/b) + f(n)$ biçimindeki
yinelemeler için: alt problem sayısı $a$, küçülme oranı $b$, birleştirme
maliyeti $f(n)$. Merge sort'ta $a=2, b=2, f(n)=O(n)$ olduğundan
$T(n) = O(n\log n)$ çıkar.}

\gsub{Uygulama: Hızlı Üs Alma (Binary Exponentiation)}

$a^n$'i \bigO{n} yerine \bigO{\log n}'de hesaplarız: $a^n = (a^{n/2})^2$
özdeşliğini kullanarak üssü her adımda yarıya indiririz.

\begin{lstlisting}
long long power(long long base, long long exp, long long mod) {
    long long result = 1;
    base %= mod;
    while (exp > 0) {
        if (exp & 1) result = result * base % mod;  // us tekse carp
        base = base * base % mod;                   // tabani karesini al
        exp >>= 1;                                  // ussu yariya indir
    }
    return result;                                  // O(log us)
}
\end{lstlisting}

% =====================================================================
\gsec{Açgözlü Algoritmalar (Greedy)}
% =====================================================================

Açgözlü algoritma, her adımda o an \emph{en iyi görünen} yerel seçimi yapar ve
geri dönmez. Her problem için işe yaramaz --- ama işe yaradığında (bir
``açgözlü seçim özelliği'' varsa) hem basit hem hızlıdır. Doğruluğu genelde
kanıt gerektirir.

\gsub{Uygulama: Etkinlik Seçimi (Activity Selection)}

Bitiş zamanlarına göre sıralı etkinliklerden, çakışmadan en fazla kaçını
seçebiliriz? Açgözlü kural: \emph{her zaman en erken biteni seç}.

\begin{lstlisting}
int maxActivities(vector<pair<int,int>> activities) {   // {bitis, baslangic}
    sort(activities.begin(), activities.end());   // bitis zamanina gore
    int count = 0, lastEnd = -1;
    for (auto [end, start] : activities)
        if (start >= lastEnd) {               // cakismiyorsa sec
            ++count;
            lastEnd = end;
        }
    return count;                             // O(n log n)
}
\end{lstlisting}

\gsub{Açgözlü Ne Zaman Yanılır?}

Açgözlü yaklaşım her zaman en iyiyi (optimum) bulmaz. Klasik örnek: bozuk para
problemi. \{1, 3, 4\} birimleriyle 6 tutarını açgözlü (en büyük parayı önce al)
$4+1+1 = 3$ parayla verir, ama optimum $3+3 = 2$ paradır! Böyle durumlarda
\emph{dinamik programlama} gerekir --- ki bir sonraki bölümün konusu budur.

\uyari{Bir açgözlü algoritma kurmadan önce, açgözlü seçimin gerçekten optimuma
götürdüğünü \emph{kanıtla} (değiş-tokuş argümanı / exchange argument) veya
küçük karşı-örneklerle test et. ``Mantıklı görünüyor'' yeterli değildir;
yukarıdaki bozuk para örneği tam da bu tuzağı gösterir.}

% =====================================================================
\gsec{Geri İzleme (Backtracking)}
% =====================================================================

Geri izleme, bir çözümü adım adım inşa eden ve bir adımın çözüme
götürmeyeceği anlaşıldığında \emph{geri alıp} (backtrack) başka bir yol deneyen
bir yöntemdir. Tüm olası kombinasyonları sistemli biçimde --- ama umutsuz
dalları budayarak --- tarar. Permütasyonlar, alt kümeler, bulmacalar
(sudoku, N-vezir) bununla çözülür.

\gsub{Uygulama: Tüm Permütasyonlar}

\begin{lstlisting}
void permutations(vector<int>& a, int k, vector<vector<int>>& result) {
    if (k == (int)a.size()) { result.push_back(a); return; }
    for (int i = k; i < (int)a.size(); ++i) {
        swap(a[k], a[i]);                  // sec
        permutations(a, k + 1, result);    // derinlere in
        swap(a[k], a[i]);                  // geri al (backtrack)
    }
}
\end{lstlisting}

\gsub{Uygulama: N-Vezir Problemi}

$N \times N$ tahtaya, birbirini tehdit etmeyecek şekilde $N$ vezir yerleştir.
Her satıra bir vezir koyar, güvenli sütun ararız; çıkmaza girersek geri döneriz.

\begin{lstlisting}
bool isSafe(vector<int>& cols, int row, int s) {
    for (int i = 0; i < row; ++i)
        if (cols[i] == s ||                        // ayni sutun
            abs(cols[i] - s) == abs(i - row))      // ayni capraz
            return false;
    return true;
}

int nQueens(int n, int row, vector<int>& cols) {
    if (row == n) return 1;                // bir cozum tamamlandi
    int count = 0;
    for (int s = 0; s < n; ++s)
        if (isSafe(cols, row, s)) {
            cols[row] = s;                 // vezir koy
            count += nQueens(n, row + 1, cols);
            // ustune yazacagimiz icin ayrica geri almaya gerek yok
        }
    return count;                          // toplam cozum sayisi
}
\end{lstlisting}

\bilgi{Geri izlemenin gücü \emph{budamadır} (pruning): olanaksız dalları erken
elemek arama uzayını dramatik biçimde küçültür. N-vezirde \code{isSafe}
kontrolü, milyonlarca imkânsız yerleşimi hiç denemeden eler. İyi budama,
üstel bir aramayı pratikte çalışabilir hale getirir.}

% =====================================================================
\gsec{Dinamik Programlama (Ayrıntılı)}
% =====================================================================

Dinamik programlama (DP), bir problemi \emph{örtüşen alt problemlere} bölüp,
her alt problemi \textbf{bir kez} çözerek ve sonucunu saklayarak (tekrar
hesaplamadan) çözen bir tekniktir. Üstel çözümleri polinom zamana indirmesiyle
algoritmanın en güçlü araçlarından biridir. Bu bölüm, DP'yi sıfırdan, bol
örnekle ve iki temel stiliyle (yukarıdan-aşağı ve aşağıdan-yukarı) ele alır.

\gsub{DP Ne Zaman Uygulanır? İki Koşul}

Bir problem DP ile çözülebilir, ancak ve ancak şu iki özelliğe sahipse:

\begin{center}
\begin{tabularx}{\textwidth}{@{}>{\bfseries\color{anaMavi}}l X@{}}
\toprule
\rowcolor{acikMavi}\color{anaMavi}\textbf{Özellik} & \color{anaMavi}\textbf{Anlamı}\\
\midrule
Örtüşen alt problemler & Aynı alt problemler tekrar tekrar çözülür (özyineleme ağacında tekrar).\\
\rowcolor{acikGri} Optimal alt yapı & Problemin optimum çözümü, alt problemlerin optimum çözümlerinden kurulur.\\
\bottomrule
\end{tabularx}
\end{center}

Örneğin Fibonacci: $F(n) = F(n-1) + F(n-2)$. Saf özyineleme $F(5)$'i
hesaplarken $F(3)$'ü \emph{iki kez}, $F(2)$'yi \emph{üç kez} hesaplar ---
alt problemler örtüşür. Toplam çağrı sayısı üsteldir (\bigO{2^n}).

\begin{lstlisting}
// SAF ozyineleme: O(2^n) -- KOTU!  fib(50) evrenin yasindan uzun surer.
long long fibSlow(int n) {
    if (n <= 1) return n;
    return fibSlow(n - 1) + fibSlow(n - 2);   // ayni degerler tekrar tekrar
}
\end{lstlisting}

\gsub{İki DP Stili: Memoization vs Tabulation}

Aynı DP'yi iki şekilde yazabiliriz. İkisi de her alt problemi bir kez çözer;
fark, \emph{yönde} ve uygulamadadır.

\begin{center}
\begin{tabularx}{\textwidth}{@{}>{\bfseries\color{anaMavi}}l X X@{}}
\toprule
\rowcolor{acikMavi}\color{anaMavi}\textbf{} & \color{anaMavi}\textbf{Memoization (yukarıdan-aşağı)} & \color{anaMavi}\textbf{Tabulation (aşağıdan-yukarı)}\\
\midrule
Yaklaşım & Özyineleme + önbellek & Döngüyle tablo doldurma\\
\rowcolor{acikGri} Yön & Büyük problemden küçüğe & Küçük alt problemden büyüğe\\
Sadece gerekli alt problemler & Evet (tembel) & Hayır (hepsi)\\
\rowcolor{acikGri} Yığın taşması riski & Var (derin özyineleme) & Yok\\
Genelde tercih & Yazması kolay & Biraz daha hızlı, bellek optimize edilebilir\\
\bottomrule
\end{tabularx}
\end{center}

\begin{lstlisting}
// (1) MEMOIZATION -- ozyineleme + onbellek. O(n)
vector<long long> memo;
long long fibMemo(int n) {
    if (n <= 1) return n;
    if (memo[n] != -1) return memo[n];         // daha once hesaplandi mi?
    return memo[n] = fibMemo(n-1) + fibMemo(n-2);
}
// Kullanim: memo.assign(n+1, -1); fibMemo(n);

// (2) TABULATION -- tabloyu kucukten buyuge doldur. O(n)
long long fibTab(int n) {
    if (n <= 1) return n;
    vector<long long> dp(n + 1);
    dp[0] = 0; dp[1] = 1;
    for (int i = 2; i <= n; ++i)
        dp[i] = dp[i-1] + dp[i-2];
    return dp[n];
}

// (3) BELLEK OPTIMIZASYONU -- sadece son iki degeri tut. O(1) bellek
long long fibFast(int n) {
    if (n <= 1) return n;
    long long prev = 0, curr = 1;
    for (int i = 2; i <= n; ++i) {
        long long next = prev + curr;
        prev = curr;
        curr = next;
    }
    return curr;
}
\end{lstlisting}

\bilgi{\textbf{DP'ye yaklaşımın 5 adımı:} (1) \emph{Durumu} tanımla ---
bir alt problemi hangi parametreler benzersiz belirler? (2) \emph{Geçişi}
(recurrence) yaz --- bir durum, hangi daha küçük durumlardan hesaplanır?
(3) \emph{Taban durumları} belirle. (4) \emph{Sırayı} belirle (memo ya da
tablo). (5) Gerekiyorsa \emph{belleği optimize et}. Bu bölümdeki her örneği bu
adımlarla okuyacağız.}

\gsub{1B DP: Merdiven Çıkma ve Bozuk Para}

\gsubsub{Merdiven Çıkma}

$n$ basamaklı bir merdiveni bir seferde 1 veya 2 basamak atlayarak kaç farklı
şekilde çıkabilirsin? \textbf{Durum:} \code{dp[i]} = $i$. basamağa kaç yolla
gelinir. \textbf{Geçiş:} \code{dp[i] = dp[i-1] + dp[i-2]} (son adım 1 veya 2).

\begin{lstlisting}
int climbStairs(int n) {
    vector<int> dp(n + 1, 0);
    dp[0] = 1;                             // yerinde durmanin 1 yolu
    dp[1] = 1;
    for (int i = 2; i <= n; ++i)
        dp[i] = dp[i-1] + dp[i-2];
    return dp[n];
}
\end{lstlisting}

\gsubsub{Bozuk Para (Coin Change) --- Minimum Para}

Verilen \code{coins} birimleriyle \code{amount}'ı oluşturmak için gereken
\emph{en az} para sayısı. (Açgözlünün yanıldığı, DP'nin doğru çözdüğü problem.)
\textbf{Durum:} \code{dp[t]} = $t$ tutarı için minimum para. \textbf{Geçiş:}
her para \code{p} için \code{dp[t] = min(dp[t], dp[t-p] + 1)}.

\begin{lstlisting}
int coinChange(const vector<int>& coins, int amount) {
    const int INF = 1e9;
    vector<int> dp(amount + 1, INF);
    dp[0] = 0;                             // 0 tutar icin 0 para
    for (int t = 1; t <= amount; ++t)
        for (int p : coins)
            if (p <= t && dp[t - p] + 1 < dp[t])
                dp[t] = dp[t - p] + 1;
    return dp[amount] == INF ? -1 : dp[amount];  // -1: mumkun degil
}
\end{lstlisting}

\gsubsub{Bozuk Para --- Kaç Farklı Yol?}

Aynı paralarla \code{amount}'ı kaç \emph{farklı} şekilde oluşturabiliriz?
Burada dış döngünün \emph{paralar} olması önemlidir --- yoksa aynı kombinasyon
farklı sırayla tekrar sayılır.

\begin{lstlisting}
long long coinCombinations(const vector<int>& coins, int amount) {
    vector<long long> dp(amount + 1, 0);
    dp[0] = 1;
    for (int p : coins)                    // DIS dongu: paralar
        for (int t = p; t <= amount; ++t)  // IC dongu: tutar
            dp[t] += dp[t - p];
    return dp[amount];
}
\end{lstlisting}

\gsub{Sırt Çantası (0/1 Knapsack)}

Klasik DP problemi: her biri bir \code{weight} ve \code{value}'e sahip
nesneler var; kapasitesi \code{W} olan çantaya, toplam ağırlık \code{W}'yi
aşmadan koyabileceğin \emph{maksimum değer} nedir? ``0/1'' her nesnenin ya
alınıp ya alınmadığını (parçalanamaz) belirtir.

\textbf{Durum:} \code{dp[i][w]} = ilk $i$ nesneyle ve $w$ kapasiteyle elde
edilen maks değer. \textbf{Geçiş:} nesne $i$'yi \emph{alma}
(\code{dp[i-1][w]}) veya \emph{al} (\code{dp[i-1][w-weight] + value}) ---
büyük olanı seç.

\begin{lstlisting}
int knapsack(const vector<int>& weight, const vector<int>& value, int W) {
    int n = weight.size();
    vector<vector<int>> dp(n + 1, vector<int>(W + 1, 0));
    for (int i = 1; i <= n; ++i)
        for (int w = 0; w <= W; ++w) {
            dp[i][w] = dp[i-1][w];         // nesne i'yi ALMA
            if (weight[i-1] <= w)          // sigiyorsa AL secenegini de dene
                dp[i][w] = max(dp[i][w],
                               dp[i-1][w - weight[i-1]] + value[i-1]);
        }
    return dp[n][W];                       // O(n * W)
}
\end{lstlisting}

\notk[Bellek optimizasyonu]{\code{dp[i]} yalnızca \code{dp[i-1]}'e bağlı
olduğundan 2B tabloyu 1B'ye indirebiliriz. Kritik nokta: iç döngüyü
\emph{tersten} (\code{w = W} $\to$ \code{weight}) çevirmek, her nesnenin bir
kez kullanılmasını garantiler.}
\begin{lstlisting}
int knapsack1B(const vector<int>& weight, const vector<int>& value, int W) {
    vector<int> dp(W + 1, 0);
    for (int i = 0; i < (int)weight.size(); ++i)
        for (int w = W; w >= weight[i]; --w)    // TERSTEN!
            dp[w] = max(dp[w], dp[w - weight[i]] + value[i]);
    return dp[W];                          // O(n*W) zaman, O(W) bellek
}
\end{lstlisting}

\gsub{En Uzun Ortak Alt Dizi (LCS)}

İki string'in \emph{ortak} (ardışık olması gerekmeyen ama sıra korunan) en
uzun alt dizisinin uzunluğu. Diff araçları, DNA karşılaştırma ve sürüm kontrolü
bunu kullanır. \textbf{Durum:} \code{dp[i][j]} = \code{a}'nın ilk $i$ karakteri
ile \code{b}'nin ilk $j$ karakteri için LCS uzunluğu.

\begin{lstlisting}
int lcs(const string& a, const string& b) {
    int n = a.size(), m = b.size();
    vector<vector<int>> dp(n + 1, vector<int>(m + 1, 0));
    for (int i = 1; i <= n; ++i)
        for (int j = 1; j <= m; ++j) {
            if (a[i-1] == b[j-1])
                dp[i][j] = dp[i-1][j-1] + 1;      // karakter eslesti
            else
                dp[i][j] = max(dp[i-1][j], dp[i][j-1]);  // birini atla
        }
    return dp[n][m];                       // O(n * m)
}
\end{lstlisting}

\gsub{En Uzun Artan Alt Dizi (LIS)}

Bir dizideki, elemanları kesin artan en uzun alt dizinin uzunluğu.
İki çözüm sunuyoruz: sezgisel \bigO{n^2} DP ve zekice \bigO{n\log n}.

\begin{lstlisting}
// (1) DP -- O(n^2). dp[i] = i'de biten LIS uzunlugu
int lisSquared(const vector<int>& a) {
    int n = a.size();
    vector<int> dp(n, 1);
    for (int i = 0; i < n; ++i)
        for (int j = 0; j < i; ++j)
            if (a[j] < a[i])
                dp[i] = max(dp[i], dp[j] + 1);
    return *max_element(dp.begin(), dp.end());
}

// (2) Sabir dizisi + ikili arama -- O(n log n)
int lisLog(const vector<int>& a) {
    vector<int> tails;                     // tails[k]: uzunlugu k+1 olan
    for (int x : a) {                      // artan altdizinin en kucuk sonu
        auto it = lower_bound(tails.begin(), tails.end(), x);
        if (it == tails.end()) tails.push_back(x);   // uzat
        else *it = x;                                // daha kucukle degistir
    }
    return tails.size();
}
\end{lstlisting}

\gsub{Düzenleme Uzaklığı (Edit Distance / Levenshtein)}

Bir string'i diğerine dönüştürmek için gereken en az işlem (ekle, sil,
değiştir). Yazım denetimi ve DNA hizalama kullanır. \textbf{Durum:}
\code{dp[i][j]} = \code{a[0..i)}'yi \code{b[0..j)}'ye çevirmenin min maliyeti.

\begin{lstlisting}
int editDistance(const string& a, const string& b) {
    int n = a.size(), m = b.size();
    vector<vector<int>> dp(n + 1, vector<int>(m + 1));
    for (int i = 0; i <= n; ++i) dp[i][0] = i;   // hepsini sil
    for (int j = 0; j <= m; ++j) dp[0][j] = j;   // hepsini ekle
    for (int i = 1; i <= n; ++i)
        for (int j = 1; j <= m; ++j) {
            if (a[i-1] == b[j-1])
                dp[i][j] = dp[i-1][j-1];         // ayni -> islem yok
            else
                dp[i][j] = 1 + min({ dp[i-1][j],     // sil
                                     dp[i][j-1],     // ekle
                                     dp[i-1][j-1] }); // degistir
        }
    return dp[n][m];                       // O(n * m)
}
\end{lstlisting}

\gsub{Grid Üzerinde DP: Minimum Yol Toplamı}

Sol üstten sağ alta, yalnızca sağa/aşağı giderek en düşük maliyetli yol.
\textbf{Durum:} \code{dp[i][j]} = $(i,j)$ hücresine ulaşmanın min maliyeti;
sadece üstten veya soldan gelinir.

\begin{lstlisting}
int minPathSum(const vector<vector<int>>& grid) {
    int n = grid.size(), m = grid[0].size();
    vector<vector<int>> dp(n, vector<int>(m));
    dp[0][0] = grid[0][0];
    for (int j = 1; j < m; ++j) dp[0][j] = dp[0][j-1] + grid[0][j];
    for (int i = 1; i < n; ++i) dp[i][0] = dp[i-1][0] + grid[i][0];
    for (int i = 1; i < n; ++i)
        for (int j = 1; j < m; ++j)
            dp[i][j] = grid[i][j] + min(dp[i-1][j], dp[i][j-1]);
    return dp[n-1][m-1];                   // O(n * m)
}
\end{lstlisting}

\gsub{Bitmask DP: Gezgin Satıcı (TSP)}

Az sayıda düğüm ($n \le 20$) olduğunda, ziyaret edilen düğüm kümesini bir
\emph{bitmask} (bit maskesi) olarak kodlayıp DP yapabiliriz. TSP: tüm şehirleri
tam bir kez gezip başa dönen en kısa tur. \textbf{Durum:}
\code{dp[mask][i]} = \code{mask} kümesindeki şehirler ziyaret edilmiş ve
şu an \code{i}'deyken min maliyet.

\begin{lstlisting}
int tsp(const vector<vector<int>>& d) {    // d: mesafe matrisi
    int n = d.size();
    const int INF = 1e9;
    vector<vector<int>> dp(1 << n, vector<int>(n, INF));
    dp[1][0] = 0;                          // sadece 0. sehir ziyaret, oradayiz
    for (int mask = 1; mask < (1 << n); ++mask)
        for (int u = 0; u < n; ++u) {
            if (dp[mask][u] == INF) continue;
            if (!(mask & (1 << u))) continue;
            for (int v = 0; v < n; ++v)    // sonraki sehir
                if (!(mask & (1 << v)))
                    dp[mask | (1 << v)][v] =
                        min(dp[mask | (1 << v)][v], dp[mask][u] + d[u][v]);
        }
    int result = INF;
    for (int u = 1; u < n; ++u)            // basa (0) donus
        result = min(result, dp[(1 << n) - 1][u] + d[u][0]);
    return result;                         // O(2^n * n^2)
}
\end{lstlisting}

\gsub{Ağaç Üzerinde DP}

DP yalnızca dizilerde değil ağaçlarda da çalışır: her düğümün cevabı,
çocuklarının cevaplarından hesaplanır (post-order). Örnek: bir ağaçta
komşu olmayan düğümlerden en büyük bağımsız küme (maksimum ağırlık).

\begin{lstlisting}
// dp[u][0] = u alinmadan alt agacin en iyisi
// dp[u][1] = u alinarak  alt agacin en iyisi
void treeDP(int u, int parent, const vector<vector<int>>& adj,
            const vector<int>& weight, vector<array<int,2>>& dp) {
    dp[u][0] = 0;
    dp[u][1] = weight[u];
    for (int v : adj[u]) {
        if (v == parent) continue;
        treeDP(v, u, adj, weight, dp);
        dp[u][0] += max(dp[v][0], dp[v][1]); // u yoksa cocuk serbest
        dp[u][1] += dp[v][0];                // u varsa cocuk alinamaz
    }
}
// Cevap: max(dp[root][0], dp[root][1])
\end{lstlisting}

\gsub{Özetle: DP Problemlerini Tanımak}

\begin{center}
\begin{tabularx}{\textwidth}{@{}>{\bfseries\color{anaMavi}}l X@{}}
\toprule
\rowcolor{acikMavi}\color{anaMavi}\textbf{Problem tipi} & \color{anaMavi}\textbf{Tipik durum tanımı}\\
\midrule
Dizide seçim (knapsack) & \code{dp[i][capacity]}\\
\rowcolor{acikGri} İki string (LCS, edit) & \code{dp[i][j]} (iki önek)\\
Alt dizi (LIS) & \code{dp[i]} = \code{i}'de biten en iyi\\
\rowcolor{acikGri} Grid & \code{dp[i][j]} = hücreye ulaşma\\
Aralık (matrix chain) & \code{dp[i][j]} = \code{[i,j]} aralığı\\
\rowcolor{acikGri} Alt küme (TSP) & \code{dp[mask][i]} (bitmask)\\
Ağaç & \code{dp[node][state]} (post-order)\\
\bottomrule
\end{tabularx}
\end{center}

\ipucu{DP öğrenmenin en iyi yolu bol problem çözmektir. ``Bu problemi daha
küçük bir versiyonundan çözebilir miyim?'' diye sor. Cevap evetse ve alt
problemler örtüşüyorsa, DP büyük ihtimalle uygulanır. Önce yavaş özyinelemeyi
yaz, sonra \code{memo} dizisi ekleyerek hızlandır --- bu yol neredeyse her
zaman işe yarar.}

% =====================================================================
\gsec{String Algoritmaları}
% =====================================================================

Metin işleme, örüntü arama ve biyoinformatik string algoritmalarına dayanır.
En temel problem \emph{örüntü eşleme}: bir metin içinde bir deseni bulmak.

\gsub{KMP: Doğrusal Zamanda Örüntü Arama}

Naif arama \bigO{nm}'dir. Knuth-Morris-Pratt (KMP), desenin kendi
tekrarlarından yararlanarak geri gitmeden \bigO{n+m}'de arar. Anahtar,
\emph{başarısızlık fonksiyonu} (LPS --- en uzun uygun önek=sonek) tablosudur.

\begin{lstlisting}
vector<int> buildLps(const string& pattern) {
    int m = pattern.size();
    vector<int> lps(m, 0);
    int len = 0;
    for (int i = 1; i < m; ) {
        if (pattern[i] == pattern[len]) lps[i++] = ++len;
        else if (len) len = lps[len - 1];   // geri sic
        else lps[i++] = 0;
    }
    return lps;
}

vector<int> kmpSearch(const string& text, const string& pattern) {
    vector<int> lps = buildLps(pattern), result;
    int i = 0, j = 0;                      // i: metin, j: desen
    while (i < (int)text.size()) {
        if (text[i] == pattern[j]) { ++i; ++j; }
        if (j == (int)pattern.size()) {    // tam eslesme bulundu
            result.push_back(i - j);
            j = lps[j - 1];
        } else if (i < (int)text.size() && text[i] != pattern[j]) {
            if (j) j = lps[j - 1];         // deseni kaydir
            else ++i;
        }
    }
    return result;                         // eslesmelerin baslangic indeksleri
}
\end{lstlisting}

\gsub{Uygulama: Palindrom Kontrolü ve Sayma}

\begin{lstlisting}
bool isPalindrome(const string& s) {
    int left = 0, right = s.size() - 1;
    while (left < right)
        if (s[left++] != s[right--]) return false;
    return true;                           // O(n)
}
\end{lstlisting}

% =====================================================================
\gsec{Bit Manipülasyonu}
% =====================================================================

Sayıları ikili tabanda, bit düzeyinde işlemek bazı problemleri son derece hızlı
ve zarif çözer. Alt küme üretme, bayrak (flag) tutma ve düşük seviye
optimizasyonlar bit işlemlerine dayanır.

\begin{center}
\begin{tabularx}{\textwidth}{@{}>{\ttfamily\bfseries\color{anaMavi}}l >{\ttfamily}l X@{}}
\toprule
\rowcolor{acikMavi}\color{anaMavi}\textbf{İşlem} & \color{anaMavi}\textbf{Kod} & \color{anaMavi}\textbf{Anlamı}\\
\midrule
Bit oku & (x >> i) \& 1 & $i$. bit 1 mi?\\
\rowcolor{acikGri} Bit kur & x | (1 << i) & $i$. biti 1 yap\\
Bit temizle & x \& \~{}(1 << i) & $i$. biti 0 yap\\
\rowcolor{acikGri} Bit çevir & x \^{} (1 << i) & $i$. biti tersle\\
Son biti al & x \& -x & En düşük kurulu bit\\
\rowcolor{acikGri} 2'nin kuvveti mi? & x \&\& !(x \& (x-1)) & Tek bit kuruluysa\\
\bottomrule
\end{tabularx}
\end{center}

\begin{lstlisting}
// Kurulu bit sayisi (population count)
int countBits(unsigned x) {
    int count = 0;
    while (x) { x &= (x - 1); ++count; }   // son kurulu biti sil
    return count;
}
// Derleyici built-in'i (cok daha hizli):
int count = __builtin_popcount(x);

// Bir kumenin TUM alt kumelerini gez (bitmask ile)
void subsets(const vector<int>& a) {
    int n = a.size();
    for (int mask = 0; mask < (1 << n); ++mask) {
        for (int i = 0; i < n; ++i)
            if (mask & (1 << i)) cout << a[i] << " ";   // i. eleman kumede
        cout << "\n";
    }                                      // toplam 2^n alt kume
}
\end{lstlisting}

\bilgi{\textbf{XOR'un sihri:} \code{a \^{} a = 0} ve \code{a \^{} 0 = a}
olduğundan, bir dizide her eleman iki kez, biri tek kez geçiyorsa, hepsini
XOR'layarak tek olanı \bigO{n} zaman ve \bigO{1} bellekte buluruz:
\code{int tek = 0; for (int x : a) tek \^{}= x;}. Bu, hash tablosuna gerek
kalmadan çalışan klasik bir numaradır.}

% =====================================================================
\gsec{Algoritma Teknikleri --- Alıştırmalar ve Çözümler}

\gsub{Sıralama}

\soru{Bir tam sayı dizisindeki \emph{inversiyon} sayısını bulun. İnversiyon, $i < j$ olmasına rağmen $A[i] > A[j]$ olan $(i, j)$ ikilisidir. Kaba kuvvet \bigO{n^2} yerine merge sort tabanlı \bigO{n \log n} bir çözüm yazın.}

\cozum
İnversiyon sayısı, bir dizinin ne kadar ``karışık'' olduğunun ölçüsüdür; sıralı bir dizide 0, tersten sıralı bir dizide ise $\binom{n}{2}$ inversiyon vardır. Merge sort'un birleştirme (merge) adımını değiştirerek bu sayıyı bedavaya sayabiliriz: sol yarı ile sağ yarıyı birleştirirken, sol taraftan bir eleman yerine sağ taraftan bir eleman seçtiğimiz her seferinde, sağdaki bu eleman soldaki \emph{kalan tüm elemanlardan} küçüktür. Yani soldaki kalan eleman sayısı kadar inversiyon bir anda sayılır.

Neden çalışır? Merge sort böl-yönet mantığıyla inversiyonları üç gruba ayırır: (1) tamamen sol yarıda olanlar, (2) tamamen sağ yarıda olanlar, (3) biri solda biri sağda olan ``çapraz'' inversiyonlar. İlk iki grup özyinelemeli çağrılarda sayılır; çapraz olanlar ise birleştirme sırasında sayılır. Toplam maliyet merge sort ile aynıdır: \bigO{n \log n} zaman, \bigO{n} ek bellek.

\begin{lstlisting}
#include <vector>
using namespace std;

// [left, right) yarı acik araligini sirala ve inversiyon sayisini dondur
long long mergeCount(vector<int>& a, vector<int>& buffer, int left, int right) {
    if (right - left <= 1) return 0;        // 0 veya 1 eleman: inversiyon yok
    int mid = left + (right - left) / 2;
    long long count = 0;
    count += mergeCount(a, buffer, left, mid);   // sol yaridaki inversiyonlar
    count += mergeCount(a, buffer, mid, right);   // sag yaridaki inversiyonlar

    // Birlestirme (merge) adimi: capraz inversiyonlari say
    int i = left, j = mid, k = left;
    while (i < mid && j < right) {
        if (a[i] <= a[j]) {
            buffer[k++] = a[i++];               // sol elemani sec, inversiyon yok
        } else {
            buffer[k++] = a[j++];               // sag eleman kucuk
            count += (mid - i);                 // soldaki kalan tum elemanlarla inversiyon
        }
    }
    while (i < mid)   buffer[k++] = a[i++];      // sol kalanlar
    while (j < right) buffer[k++] = a[j++];      // sag kalanlar
    for (int t = left; t < right; ++t) a[t] = buffer[t]; // geri kopyala
    return count;
}

long long countInversions(vector<int> a) {
    vector<int> buffer(a.size());
    return mergeCount(a, buffer, 0, (int)a.size());
}
\end{lstlisting}

\uyari{İnversiyon sayısı $n$ orta büyüklükte bile $\binom{n}{2}$ mertebesine ulaşabildiği için sonucu mutlaka \code{long long} tipinde tutun; \code{int} taşması sık yapılan bir hatadır.}

\soru{Sıralanmamış bir dizide $k$. en büyük elemanı bulun. Diziyi tamamen sıralamadan (\bigO{n \log n}), ortalama \bigO{n} zamanda çalışan QuickSelect algoritmasını yazın.}

\cozum
QuickSelect, quicksort'un ``sadece bir tarafa in'' versiyonudur. Bir pivot seçip diziyi bölümlere ayırdıktan (partition) sonra, pivotun nihai konumu (index) bize onun \emph{kaçıncı sıradaki} eleman olduğunu söyler. Aradığımız konumla karşılaştırıp sadece pivotun bir tarafına özyineleniriz; diğer tarafı tamamen atarız. Quicksort iki tarafa da indiği için \bigO{n \log n}, QuickSelect tek tarafa indiği için ortalama $n + n/2 + n/4 + \dots \approx 2n = \bigO{n}$ olur.

``$k$. en büyük'' istendiği için, $k$. en büyük eleman artan sıralamada $(n - k)$. indekstir. Pivotu rastgele seçmek en kötü durum olan \bigO{n^2} olasılığını pratikte yok eder (kötü niyetli girdiye karşı koruma).

\begin{lstlisting}
#include <vector>
#include <cstdlib>
using namespace std;

// [left, right] araliginda partition; pivotun son indeksini dondur (Lomuto)
int partition(vector<int>& a, int left, int right) {
    int pivotIndex = left + rand() % (right - left + 1);
    swap(a[pivotIndex], a[right]);          // pivotu sona tasi
    int pivot = a[right];
    int store = left;
    for (int i = left; i < right; ++i)
        if (a[i] < pivot) swap(a[i], a[store++]);
    swap(a[store], a[right]);               // pivotu dogru yerine koy
    return store;
}

int kthLargest(vector<int> a, int k) {
    int n = (int)a.size();
    int target = n - k;                     // k. en buyuk = (n-k). en kucuk (0-tabanli)
    int left = 0, right = n - 1;
    while (left <= right) {
        int p = partition(a, left, right);
        if (p == target) return a[p];        // tam yerinde bulduk
        else if (p < target) left = p + 1;   // sag tarafa in
        else right = p - 1;                  // sol tarafa in
    }
    return -1;                               // buraya gelinmez (gecerli k icin)
}
\end{lstlisting}

\bilgi{Deterministik en kötü durum \bigO{n} garantisi isteniyorsa ``median of medians'' pivot seçimi kullanılır, ancak sabit çarpanı büyük olduğu için pratikte rastgele pivotlu QuickSelect tercih edilir. Alternatif olarak boyut $k$'lık bir min-heap ile \bigO{n \log k} elde edilebilir.}

\soru{Yalnızca 0, 1 ve 2 değerlerinden oluşan (kırmızı, beyaz, mavi) bir diziyi tek geçişte ve sabit ek bellekle sıralayın (Dutch National Flag / Hollanda bayrağı problemi).}

\cozum
Sayma sıralaması (counting sort) da işi görür ama iki geçiş gerektirir. Dijkstra'nın ``Hollanda bayrağı'' algoritması tek geçişte, üç işaretçi (three pointers) ile çözer. Fikir, diziyi dört bölgeye ayırmaktır: kesin 0'lar, kesin 1'ler, işlenmemiş bölge, kesin 2'ler. \code{low} işaretçisi 0'ların sınırını, \code{high} 2'lerin sınırını, \code{mid} ise tarama başını temsil eder. Değişmez (invariant): \code{[0, low)} hep 0, \code{[low, mid)} hep 1, \code{(high, n)} hep 2.

\code{mid} eleman 0 ise \code{low} ile takas edip ikisini de ilerletiriz. 2 ise \code{high} ile takas edip sadece \code{high}'ı geri çekeriz (takas edilen eleman henüz incelenmediği için \code{mid} ilerlemez!). 1 ise sadece \code{mid} ilerler. Karmaşıklık \bigO{n} zaman, \bigO{1} bellek.

\begin{lstlisting}
#include <vector>
using namespace std;

void sortColors(vector<int>& nums) {
    int low = 0, mid = 0, high = (int)nums.size() - 1;
    while (mid <= high) {
        if (nums[mid] == 0) {
            swap(nums[low++], nums[mid++]);  // 0'i one at
        } else if (nums[mid] == 2) {
            swap(nums[mid], nums[high--]);   // 2'yi sona at, mid ILERLEMEZ
        } else {
            mid++;                           // 1 zaten yerinde
        }
    }
}
\end{lstlisting}

\ipucu{Neden 2 ile takas ederken \code{mid++} yapmıyoruz? Çünkü \code{high}'dan gelen elemanın ne olduğunu bilmiyoruz; onu bir sonraki turda yeniden incelemeliyiz. Oysa 0 ile takas ettiğimiz \code{low} konumu zaten taranmış (orada ya 1 ya da 0 vardı) olduğu için güvenle ilerleyebiliriz.}

\gsub{Arama}

\soru{Artan sıralı, ancak bilinmeyen bir noktadan döndürülmüş (rotated) ve tekrar içermeyen bir dizide verilen hedefi \bigO{\log n} zamanda arayın. Örn. $[4,5,6,7,0,1,2]$ içinde 0'ı bulun.}

\cozum
Döndürülmüş sıralı dizi, iki sıralı parçadan oluşur. İkili aramanın (binary search) her adımında, \code{mid}'e baktığımızda yarılardan \emph{en az biri} kesinlikle sıralıdır. Hangi yarının sıralı olduğunu \code{nums[left] <= nums[mid]} kontrolüyle anlarız. Sıralı yarının uç değerleri arasında hedef var mı diye bakarız: varsa o yarıya, yoksa diğer yarıya ineriz. Böylece klasik ikili aramanın \bigO{\log n} garantisini korumuş oluruz.

Kritik nokta, karşılaştırmalarda \code{<=} mı \code{<} mu kullanılacağıdır; \code{nums[left] <= nums[mid]} eşitliği tek elemanlı yarıyı (left == mid) doğru sınıflandırır.

\begin{lstlisting}
#include <vector>
using namespace std;

int searchRotated(vector<int>& nums, int target) {
    int left = 0, right = (int)nums.size() - 1;
    while (left <= right) {
        int mid = left + (right - left) / 2;
        if (nums[mid] == target) return mid;
        if (nums[left] <= nums[mid]) {          // sol yari sirali
            if (nums[left] <= target && target < nums[mid])
                right = mid - 1;                 // hedef sol yarida
            else
                left = mid + 1;
        } else {                                 // sag yari sirali
            if (nums[mid] < target && target <= nums[right])
                left = mid + 1;                  // hedef sag yarida
            else
                right = mid - 1;
        }
    }
    return -1;
}
\end{lstlisting}

\uyari{Dizi tekrar eden elemanlar içeriyorsa (örn. $[1,1,1,0,1]$) \code{nums[left] == nums[mid]} durumunda hangi yarının sıralı olduğunu anlayamayız; bu durumda \code{left++} ile bir adım daralmak gerekir ve en kötü durum \bigO{n}'e çıkar.}

\soru{Cevap üzerinde ikili arama: Bir bant (konveyör) üzerinde belli sırayla taşınacak paketlerin ağırlıkları \code{weights} dizisinde verilmiş. Paketleri sırayı bozmadan en fazla \code{days} günde taşımak istiyoruz. Günlük minimum gemi kapasitesini bulun.}

\cozum
Bu problemi doğrudan çözmek zor, ama ``$X$ kapasiteyle bu işi \code{days} günde bitirebilir miyiz?'' sorusunu \bigO{n}'de yanıtlayabiliriz (açgözlüce her güne sığdığı kadar paket koy). Kritik gözlem: kapasite arttıkça gerekli gün sayısı \emph{monoton azalır}. Yani bir eşik kapasitenin altında ``hayır'', üstünde ``evet'' cevabı vardır. Bu monotonluk, cevabın kendisi üzerinde ikili arama yapmamıza olanak tanır.

Arama aralığının alt sınırı \code{max(weights)}'tir (en ağır paket tek başına bir güne sığmalı), üst sınırı \code{sum(weights)}'tir (her şeyi tek günde taşı). İkili aramayla en küçük geçerli kapasiteyi buluruz. Karmaşıklık: \bigO{n \log(\text{sum})}.

\begin{lstlisting}
#include <vector>
#include <numeric>
#include <algorithm>
using namespace std;

// capacity kapasiteyle kac gun gerekir? days'ten kucuk/esitse yeterli
bool canShip(vector<int>& weights, int capacity, int days) {
    int neededDays = 1, load = 0;
    for (int w : weights) {
        if (load + w > capacity) {   // bu paket bu gune sigmaz
            neededDays++;             // yeni gune gec
            load = 0;
        }
        load += w;
    }
    return neededDays <= days;
}

int shipWithinDays(vector<int>& weights, int days) {
    int low = *max_element(weights.begin(), weights.end());
    int high = accumulate(weights.begin(), weights.end(), 0);
    while (low < high) {              // en kucuk gecerli kapasiteyi bul
        int mid = low + (high - low) / 2;
        if (canShip(weights, mid, days))
            high = mid;               // mid yeterli, daha kucugunu dene
        else
            low = mid + 1;            // mid yetersiz, buyut
    }
    return low;
}
\end{lstlisting}

\notk[Genel Kalıp]{``Cevap üzerinde ikili arama'' deseninin anahtarı bir \emph{karar fonksiyonu} \code{feasible(x)} bulmaktır: bu fonksiyon bir eşikten sonra hep \code{true} (veya hep \code{false}) olmalı. Bu monotonluk sağlandığında, en küçük/en büyük geçerli değeri logaritmik sürede buluruz. ``Minimum en büyük değer'' ya da ``maksimum en küçük değer'' isteyen problemler genelde bu kalıba oturur.}

\soru{Unimodal (tek tepeli) bir dizide -- önce kesin artıp sonra kesin azalan -- tepe (peak) elemanı ternary search (üçlü arama) ile bulun.}

\cozum
İkili arama iki parçaya böler; ternary search aralığı iki iç noktayla üç parçaya böler ve her adımda dörtte birini eler. Unimodal bir fonksiyonda \code{m1 < m2} ise tepe kesinlikle \code{m1}'in solunda değildir (çünkü fonksiyon hâlâ artıyor), dolayısıyla \code{left = m1 + 1}. Aksi halde tepe \code{m2}'nin sağında değildir, \code{right = m2 - 1}. Her turda aralık $2/3$ oranında daraldığı için karmaşıklık \bigO{\log_{3/2} n} = \bigO{\log n}'dir.

Bu teknik özellikle sürekli (real) alanlarda konveks/konkav fonksiyonların ekstremumunu bulmakta çok kullanışlıdır; orada döngüyü sabit sayıda (örn. 200) iterasyonla ya da bir \code{eps} eşiğiyle sınırlarız.

\begin{lstlisting}
#include <vector>
using namespace std;

int findPeak(vector<int>& a) {
    int left = 0, right = (int)a.size() - 1;
    while (left < right) {
        int m1 = left + (right - left) / 3;
        int m2 = right - (right - left) / 3;
        if (a[m1] < a[m2])
            left = m1 + 1;      // tepe m1'in sagi tarafinda
        else
            right = m2 - 1;     // tepe m2'nin sol tarafinda
    }
    return left;                // left == right: tepe
}
\end{lstlisting}

\bilgi{Ayrık (integer) dizilerde tepe bulmak için ikili arama ($a[mid]$ ile $a[mid+1]$ karşılaştırması) daha yaygın ve sabit çarpanı düşüktür. Ternary search'ün asıl gücü, konveks/unimodal \emph{sürekli} fonksiyonların minimum/maksimumunu bulurken ortaya çıkar.}

\gsub{Böl ve Yönet}

\soru{Maksimum alt dizi toplamını (maximum subarray sum) böl-yönet ile \bigO{n \log n} çözün, ardından Kadane algoritmasının \bigO{n} çözümüyle karşılaştırın.}

\cozum
Böl-yönet kurgusunda diziyi ortadan ikiye böleriz. Maksimum alt dizi üç yerden birinde olur: (1) tamamen sol yarıda, (2) tamamen sağ yarıda, (3) orta çizgiyi \emph{geçen} bir alt dizide. İlk ikisini özyinelemeyle buluruz. Üçüncüsü için, ortadan başlayıp sola doğru en iyi eki ve sağa doğru en iyi eki ayrı ayrı \bigO{n}'de hesaplayıp toplarız. Yineleme bağıntısı $T(n) = 2T(n/2) + \bigO{n}$ olduğundan (aynı merge sort gibi) sonuç \bigO{n \log n}'dir.

Kadane algoritması ise aynı problemi tek geçişte, dinamik programlama fikriyle çözer: ``şu ana kadar biten en iyi alt dizi toplamı''nı taşır ve negatifleşince sıfırlar. Karşılaştırmada Kadane hem daha basit hem de \bigO{n} ile daha hızlıdır; böl-yönet çözümü ise böl-yönet düşünme biçimini öğretmek için değerlidir.

\begin{lstlisting}
#include <vector>
#include <algorithm>
#include <climits>
using namespace std;

// --- Bol-yonet: O(n log n) ---
long long crossSum(vector<int>& a, int left, int mid, int right) {
    long long leftBest = LLONG_MIN, sum = 0;
    for (int i = mid; i >= left; --i) { sum += a[i]; leftBest = max(leftBest, sum); }
    long long rightBest = LLONG_MIN; sum = 0;
    for (int i = mid + 1; i <= right; ++i) { sum += a[i]; rightBest = max(rightBest, sum); }
    return leftBest + rightBest;
}
long long maxSubDC(vector<int>& a, int left, int right) {
    if (left == right) return a[left];
    int mid = left + (right - left) / 2;
    return max({ maxSubDC(a, left, mid),
                 maxSubDC(a, mid + 1, right),
                 crossSum(a, left, mid, right) });
}

// --- Kadane: O(n) ---
long long maxSubKadane(vector<int>& a) {
    long long best = a[0], cur = a[0];
    for (size_t i = 1; i < a.size(); ++i) {
        cur = max((long long)a[i], cur + a[i]); // yeni basla ya da uzat
        best = max(best, cur);
    }
    return best;
}
\end{lstlisting}

\uyari{Tüm elemanlar negatif olabileceği için başlangıç değerini 0 değil, \code{a[0]} (ya da \code{LLONG\_MIN}) yapın. \code{best = 0} ile başlamak, tümü negatif dizilerde yanlışlıkla boş alt diziyi (toplam 0) döndürür.}

\soru{Hızlı üs alma (fast exponentiation) ile matris kuvveti alarak $n$. Fibonacci sayısını \bigO{\log n} zamanda hesaplayın. (Yeni konu, kısa)}

\cozum
Fibonacci için $\begin{psmallmatrix}F_{n+1}\\F_n\end{psmallmatrix} = \begin{psmallmatrix}1&1\\1&0\end{psmallmatrix}^{n} \begin{psmallmatrix}F_1\\F_0\end{psmallmatrix}$ özdeşliği geçerlidir. Yani $2\times2$ ``dönüşüm matrisinin'' $n$. kuvvetini alırsak $F_n$'i buluruz. Bir sayının $n$. kuvvetini hesaplarken $x^n = (x^{n/2})^2$ (çift $n$) özelliğini kullanan ``binary exponentiation'', çarpma sayısını $n$'den $\log n$'e düşürür. Aynı fikir matrislerde de çalışır: skaler çarpma yerine matris çarpımı kullanırız. Böylece Fibonacci \bigO{\log n} matris çarpımıyla (her biri sabit $2^3$ işlem) hesaplanır.

\begin{lstlisting}
#include <array>
using namespace std;
typedef array<array<long long,2>,2> Mat;

Mat mul(const Mat& A, const Mat& B) {
    Mat C = {{{0,0},{0,0}}};
    for (int i = 0; i < 2; ++i)
        for (int j = 0; j < 2; ++j)
            for (int k = 0; k < 2; ++k)
                C[i][j] += A[i][k] * B[k][j];   // (mod almak gerekirse burada)
    return C;
}

long long fib(long long n) {
    if (n == 0) return 0;
    Mat result = {{{1,0},{0,1}}};               // birim matris
    Mat base   = {{{1,1},{1,0}}};
    while (n > 0) {                              // binary exponentiation
        if (n & 1) result = mul(result, base);  // bit set ise carp
        base = mul(base, base);                 // tabani karele
        n >>= 1;
    }
    return result[0][1];                        // F_n bu hucrede
}
\end{lstlisting}

\bilgi{Aynı ``binary exponentiation'' iskeleti, modüler üs alma ($a^b \bmod m$) için kriptografide, doğrusal yineleme bağıntılarını hızlandırmak için yarışma programlamada yaygın kullanılır. Matris çarpımını \code{mod} ile yaparsanız çok büyük Fibonacci sayılarını taşmadan hesaplarsınız.}

\gsub{Açgözlü (Greedy)}

\soru{Aktivite seçimi (activity selection): Başlangıç ve bitiş zamanları verilen aktivitelerden, hiçbiri çakışmayacak şekilde en fazla kaç tanesini seçebileceğinizi bulun.}

\cozum
Açgözlü sezgi: her adımda \emph{en erken biten} aktiviteyi seçmek, geriye kalan zamanı maksimize eder ve en fazla aktiviteye yer bırakır. Bunu kanıtlamak için ``değiştirme argümanı'' (exchange argument) kullanılır: herhangi bir optimal çözümdeki ilk aktiviteyi, en erken biten aktiviteyle değiştirsek çözüm hâlâ geçerli ve en az o kadar büyük kalır. Dolayısıyla aktiviteleri bitiş zamanına göre sıralayıp, son seçilenle çakışmayan ilk aktiviteyi almak optimaldir. Sıralama \bigO{n \log n}, tarama \bigO{n}.

\begin{lstlisting}
#include <vector>
#include <algorithm>
using namespace std;

int maxActivities(vector<pair<int,int>> acts) { // {start, finish}
    // Bitis zamanina gore sirala
    sort(acts.begin(), acts.end(),
         [](auto& x, auto& y){ return x.second < y.second; });
    int count = 0, lastFinish = -1;
    for (auto& [start, finish] : acts) {
        if (start >= lastFinish) {   // son secilenle cakismiyor
            count++;
            lastFinish = finish;
        }
    }
    return count;
}
\end{lstlisting}

\notk[Greedy ne zaman doğru?]{Açgözlü yaklaşım ancak problem ``greedy choice property'' ve ``optimal substructure'' özelliklerini taşıyorsa doğru sonuç verir. Bunu her zaman ya bir değiştirme argümanıyla kanıtlamalı ya da küçük karşı örneklerle test etmelisiniz; ``sezgisel olarak doğru görünmesi'' yeterli değildir.}

\soru{Jump Game: \code{nums[i]}, $i$ indeksinden ileri atlayabileceğiniz maksimum adımı gösterir. 0. indeksten son indekse ulaşılabilir mi?}

\cozum
Açgözlü ölçü, o ana kadar ulaşılabilen \emph{en uzak indekstir} (\code{reach}). Diziyi soldan sağa tararken, eğer bulunduğumuz indeks \code{reach}'i aşmışsa oraya hiç ulaşamamışız demektir, \code{false} döneriz. Aksi halde \code{reach}'i \code{i + nums[i]} ile güncelleriz. En baştan \code{reach}, son indekse eşit veya büyük olursa ulaşılabilir. Tek geçiş, \bigO{n} zaman, \bigO{1} bellek. Bu, dinamik programlama çözümünden (her indekse ulaşılabilir mi tablosu) hem daha hızlı hem daha basittir.

\begin{lstlisting}
#include <vector>
using namespace std;

bool canJump(vector<int>& nums) {
    int reach = 0;                       // ulasabildigimiz en uzak indeks
    for (int i = 0; i < (int)nums.size(); ++i) {
        if (i > reach) return false;     // buraya hic ulasilamadi
        reach = max(reach, i + nums[i]); // menzili genislet
    }
    return true;                         // dongu bitti demek son indekse ulasildi
}
\end{lstlisting}

\soru{Para üstü problemi: Verilen madeni para türleriyle bir tutarı en az sayıda parayla ödemek istiyoruz. Açgözlü (her seferinde en büyük parayı seç) yaklaşımının neden \emph{genelde} yanlış olduğunu bir örnekle gösterin.}

\cozum
Açgözlü yaklaşım, standart para sistemlerinde (örn. 1, 5, 10, 25) doğru çalışır ama her para kümesinde değil. Klasik karşı örnek: paralar \code{[1, 3, 4]} ve tutar 6 olsun. Açgözlü önce 4'ü seçer, kalan 2 için iki tane 1 kullanır: toplam 3 para (4+1+1). Oysa optimal çözüm 3+3, yani sadece 2 paradır. Açgözlü ``şimdiki en iyi'' seçimi yaptı ama bu, ileride daha kötü bir duruma sürükledi -- greedy choice property burada geçerli değil.

Doğru genel çözüm dinamik programlamadır (bkz. \emph{Dinamik Programlama} bölümündeki ``madeni para -- min'' sorusu): \code{dp[x]} = $x$ tutarını ödemek için gereken minimum para. Her tutar için tüm para türlerini deneyip en iyisini alırız; karmaşıklık \bigO{\text{amount} \times \text{coins}}. Bu, açgözlünün göremediği ``geriye dönük'' en iyi bileşimleri de değerlendirir.

\begin{lstlisting}
#include <vector>
#include <algorithm>
using namespace std;

// Yanlis: acgozlu (karsi ornek uretmek icin)
int coinGreedy(vector<int> coins, int amount) {
    sort(coins.rbegin(), coins.rend());   // buyukten kucuge
    int count = 0;
    for (int c : coins)
        while (amount >= c) { amount -= c; count++; }
    return amount == 0 ? count : -1;      // coins={1,3,4}, amount=6 -> 3 (yanlis!)
}
\end{lstlisting}

\uyari{Açgözlü para üstü yalnızca ``canonical'' (kurallı) para sistemlerinde optimaldir. Rastgele bir para kümesi verildiğinde her zaman DP kullanın; aksi halde yukarıdaki gibi optimalin altında sonuçlar alırsınız.}

\soru{Aralık birleştirme (merge intervals): Örtüşen aralıklardan oluşan bir listeyi, örtüşenleri birleştirerek sadeleştirin. Örn. $[[1,3],[2,6],[8,10]] \to [[1,6],[8,10]]$.}

\cozum
Aralıkları başlangıç noktasına göre sıralarsak, örtüşen aralıklar mutlaka listede yan yana gelir. Sonucu tutan listeye ilk aralığı ekleriz; her yeni aralık için, son eklenen aralığın bitişi yeni aralığın başlangıcından büyük veya eşitse örtüşme var demektir -- son aralığın bitişini \code{max} ile genişletiriz. Aksi halde araya boşluk vardır, yeni aralığı olduğu gibi ekleriz. Sıralama \bigO{n \log n} baskın maliyettir; tarama \bigO{n}.

\begin{lstlisting}
#include <vector>
#include <algorithm>
using namespace std;

vector<vector<int>> mergeIntervals(vector<vector<int>> intervals) {
    if (intervals.empty()) return {};
    sort(intervals.begin(), intervals.end());   // baslangica gore sirala
    vector<vector<int>> merged;
    merged.push_back(intervals[0]);
    for (size_t i = 1; i < intervals.size(); ++i) {
        if (intervals[i][0] <= merged.back()[1]) // ortusuyor
            merged.back()[1] = max(merged.back()[1], intervals[i][1]);
        else
            merged.push_back(intervals[i]);       // yeni ayri aralik
    }
    return merged;
}
\end{lstlisting}

\gsub{Geri İzleme (Backtracking)}

\soru{Tekrarsız bir dizinin tüm alt kümelerini (subsets / power set) üretin. $n$ elemanlı bir kümenin $2^n$ alt kümesi vardır.}

\cozum
Geri izleme (backtracking) bir karar ağacını sistematik gezmektir. Her eleman için iki karar veririz: alt kümeye \emph{ekle} ya da \emph{ekleme}. Bir \code{start} indeksi taşıyarak her adımda mevcut kısmi çözümü (partial solution) kaydeder, sonra kalan her elemanı sırayla ekleyip özyineleniriz; dönüşte eklediğimizi geri alırız (\code{pop\_back}) -- işte ``geri izleme'' budur. Toplam $2^n$ alt küme ürettiğimiz için karmaşıklık \bigO{n \cdot 2^n} (her alt kümeyi kopyalamak dahil).

\begin{lstlisting}
#include <vector>
using namespace std;

void backtrack(vector<int>& nums, int start,
               vector<int>& current, vector<vector<int>>& result) {
    result.push_back(current);              // her dugum gecerli bir alt kume
    for (int i = start; i < (int)nums.size(); ++i) {
        current.push_back(nums[i]);         // sec
        backtrack(nums, i + 1, current, result);
        current.pop_back();                 // geri al (backtrack)
    }
}

vector<vector<int>> subsets(vector<int>& nums) {
    vector<vector<int>> result;
    vector<int> current;
    backtrack(nums, 0, current, result);
    return result;
}
\end{lstlisting}

\ipucu{Alt kümeleri bit maskeleriyle de üretebilirsiniz: 0'dan $2^n - 1$'e her sayının ikili gösterimindeki set bitler, o alt kümeye hangi elemanların girdiğini söyler (bkz. \emph{Bit Manipülasyonu} bölümü). Backtracking versiyonu ise ``kısıt ekleme'' (örn. tekrarları atlama) gerektiren durumlarda daha esnektir.}

\soru{Tekrarsız bir dizinin tüm permütasyonlarını üretin. $n$ eleman için $n!$ permütasyon vardır.}

\cozum
Permütasyonda sıra önemlidir, bu yüzden her pozisyona kalan tüm elemanları deneriz. Hangi elemanların kullanıldığını bir \code{used} boolean dizisiyle izleriz. Kısmi permütasyon tam uzunluğa ulaştığında sonuca ekleriz. Alternatif olarak yerinde takas (swap) yöntemiyle ekstra bellek olmadan da yapılabilir. Karmaşıklık \bigO{n \cdot n!}.

\begin{lstlisting}
#include <vector>
using namespace std;

void permute(vector<int>& nums, vector<int>& current,
             vector<bool>& used, vector<vector<int>>& result) {
    if (current.size() == nums.size()) {    // tam permutasyon hazir
        result.push_back(current);
        return;
    }
    for (int i = 0; i < (int)nums.size(); ++i) {
        if (used[i]) continue;              // zaten kullanildi
        used[i] = true;
        current.push_back(nums[i]);
        permute(nums, current, used, result);
        current.pop_back();                 // geri al
        used[i] = false;
    }
}

vector<vector<int>> permutations(vector<int>& nums) {
    vector<vector<int>> result;
    vector<int> current;
    vector<bool> used(nums.size(), false);
    permute(nums, current, used, result);
    return result;
}
\end{lstlisting}

\soru{Combination Sum: Tekrarsız pozitif \code{candidates} dizisi ve bir \code{target} verildiğinde, toplamı hedefe eşit olan tüm kombinasyonları bulun. Her aday \emph{sınırsız} kez kullanılabilir.}

\cozum
Her adayı sınırsız kullanabildiğimiz için, bir adayı seçtikten sonra özyinelemeli çağrıda \code{start} indeksini \emph{ilerletmeyiz} (\code{i}'de kalırız); böylece aynı aday tekrar seçilebilir. Aday listesini indeksle sınırlamak, aynı kombinasyonun farklı sıralamalarını (örn. [2,3] ve [3,2]) saymamızı engeller -- yalnızca artan sırada üretiriz. Kalan hedef 0 olunca geçerli bir kombinasyon buluruz; negatife düşerse dalı buda (prune). Erken budama için adayları sıralayıp \code{remaining - candidates[i] < 0} olunca döngüyü kırabiliriz.

\begin{lstlisting}
#include <vector>
#include <algorithm>
using namespace std;

void findCombos(vector<int>& candidates, int remaining, int start,
                vector<int>& current, vector<vector<int>>& result) {
    if (remaining == 0) { result.push_back(current); return; }
    for (int i = start; i < (int)candidates.size(); ++i) {
        if (candidates[i] > remaining) break;   // sirali oldugu icin kalanlar da buyuk
        current.push_back(candidates[i]);
        // i ile devam: ayni aday tekrar kullanilabilir
        findCombos(candidates, remaining - candidates[i], i, current, result);
        current.pop_back();                     // geri al
    }
}

vector<vector<int>> combinationSum(vector<int> candidates, int target) {
    sort(candidates.begin(), candidates.end()); // budama icin sirala
    vector<vector<int>> result;
    vector<int> current;
    findCombos(candidates, target, 0, current, result);
    return result;
}
\end{lstlisting}

\soru{N-Queens: $N \times N$ satranç tahtasına, hiçbiri birbirini tehdit etmeyecek şekilde $N$ vezir yerleştirmenin kaç farklı yolu olduğunu bulun.}

\cozum
Vezirler aynı satır, sütun ve çaprazları paylaşamaz. Satır satır ilerleyerek her satıra tam bir vezir koyarız; böylece satır çakışması otomatik önlenir. Her sütun ve iki çapraz yönü (sol-üst ve sağ-üst) için birer boolean dizi tutarak bir konumun güvenli olup olmadığını \bigO{1}'de kontrol ederiz. Çaprazları tanımlarken \code{row + col} sabiti bir çapraz ailesini, \code{row - col} sabiti diğerini belirler (ikincisini negatiften kaçınmak için \code{row - col + N} ile kaydırırız). Bir satıra hiçbir güvenli sütun yoksa geri izleriz. Karmaşıklık üst sınırı \bigO{N!}'dir ama budama sayesinde pratikte çok daha hızlıdır.

\begin{lstlisting}
#include <vector>
using namespace std;

int solveQueens(int row, int n, vector<bool>& cols,
                vector<bool>& diag1, vector<bool>& diag2) {
    if (row == n) return 1;                  // tum satirlara yerlestik: 1 cozum
    int count = 0;
    for (int col = 0; col < n; ++col) {
        int d1 = row + col;                  // sol-alt / sag-ust caprazi
        int d2 = row - col + n;              // sol-ust / sag-alt caprazi (kaydirmali)
        if (cols[col] || diag1[d1] || diag2[d2]) continue; // tehdit var
        cols[col] = diag1[d1] = diag2[d2] = true;          // yerlestir
        count += solveQueens(row + 1, n, cols, diag1, diag2);
        cols[col] = diag1[d1] = diag2[d2] = false;         // geri al
    }
    return count;
}

int totalNQueens(int n) {
    vector<bool> cols(n, false), diag1(2*n, false), diag2(2*n, false);
    return solveQueens(0, n, cols, diag1, diag2);
}
\end{lstlisting}

\bilgi{Çapraz indekslemede $row+col$ değeri $0$ ile $2n-2$ arasında, $row-col$ değeri ise $-(n-1)$ ile $n-1$ arasında değişir; ikincisini negatif indeksten kurtarmak için \code{+ n} ile kaydırırız. Bu, N-Queens'in klasik ve öğretici bit numarasıdır.}

\gsub{Dinamik Programlama}

\soru{Ev Soyguncusu (House Robber): Her evde belli miktarda para var. Yan yana iki evi soyamazsınız (alarm çalar). Soyabileceğiniz maksimum toplam parayı bulun.}

\cozum
Bu, dinamik programlamanın en temiz örneklerinden biridir. Her ev için iki seçenek var: \emph{soy} ya da \emph{atla}. \code{dp[i]}, ilk $i$ evi göz önüne aldığımızdaki maksimum kazanç olsun. $i$. evi soyarsak $i-1$'i soyamayız, yani \code{dp[i-2] + nums[i]} kazanırız; soymazsak \code{dp[i-1]} kadar kalırız. İkisinin maksimumunu alırız: \code{dp[i] = max(dp[i-1], dp[i-2] + nums[i])}. Bu bağıntı ``optimal substructure'' içerir: büyük problemin çözümü küçük alt problemlerin çözümünden kurulur. Sadece son iki değer gerektiğinden \bigO{1} bellekle, \bigO{n} zamanda çözülür.

\begin{lstlisting}
#include <vector>
#include <algorithm>
using namespace std;

int rob(vector<int>& nums) {
    int prev2 = 0;   // dp[i-2]
    int prev1 = 0;   // dp[i-1]
    for (int money : nums) {
        int cur = max(prev1, prev2 + money); // soyma vs. soy
        prev2 = prev1;
        prev1 = cur;
    }
    return prev1;
}
\end{lstlisting}

\soru{Benzersiz Yollar (Unique Paths): $m \times n$ bir gridde sol üst köşeden sağ alt köşeye, sadece \emph{sağa} ve \emph{aşağı} giderek kaç farklı yol vardır?}

\cozum
Bir hücreye yalnızca yukarıdan ya da soldan gelinebilir. Dolayısıyla o hücreye varan yol sayısı, üstteki ve soldaki hücrelere varan yolların toplamıdır: \code{dp[i][j] = dp[i-1][j] + dp[i][j-1]}. İlk satır ve ilk sütundaki her hücreye tek yol vardır (hep düz git). Tabloyu satır satır doldururuz. 2B tablo \bigO{mn} bellek kullanır ama sadece bir önceki satıra ihtiyaç olduğundan tek boyutlu diziyle \bigO{n} belleğe indirilebilir (aşağıdaki gibi). Zaman \bigO{mn}.

\begin{lstlisting}
#include <vector>
using namespace std;

int uniquePaths(int m, int n) {
    vector<int> dp(n, 1);                 // ilk satir: her hucreye 1 yol
    for (int i = 1; i < m; ++i)
        for (int j = 1; j < n; ++j)
            dp[j] += dp[j - 1];           // ust(dp[j]) + sol(dp[j-1])
    return dp[n - 1];
}
\end{lstlisting}

\bilgi{Bu problemin kapalı formu da vardır: sağ alt köşeye ulaşmak için toplam $(m-1)+(n-1)$ hamlenin $(m-1)$'ini ``aşağı'' seçmek gerektiğinden cevap $\binom{m+n-2}{m-1}$'dir. DP çözümü ise engelli grid gibi varyasyonlara kolayca genişler.}

\soru{Kelime Kırma (Word Break): Bir \code{s} dizgesi ve bir sözlük verildiğinde, \code{s}'in sözlükteki kelimelerin (tekrarlı) birleşimiyle bölünüp bölünemeyeceğini belirleyin.}

\cozum
\code{dp[i]}, \code{s}'in ilk $i$ karakterinin sözlük kelimeleriyle bölünebilir olup olmadığını tutsun. \code{dp[0] = true} (boş dizge her zaman bölünür). Her $i$ için, ondan önceki her $j$ noktasına bakarız: eğer \code{dp[j]} doğruysa (yani $[0, j)$ bölünüyorsa) ve \code{s[j..i)} alt dizgesi sözlükteyse, o zaman \code{dp[i]} de doğrudur. Sözlüğü \code{unordered\_set}'e koyarak kelime kontrolünü \bigO{1}'e yaklaştırırız (kelime uzunluğu sabit varsayımıyla). Karmaşıklık \bigO{n^2} alt dizge çifti $\times$ kontrol maliyeti.

\begin{lstlisting}
#include <string>
#include <vector>
#include <unordered_set>
using namespace std;

bool wordBreak(string s, vector<string>& wordDict) {
    unordered_set<string> dict(wordDict.begin(), wordDict.end());
    int n = (int)s.size();
    vector<bool> dp(n + 1, false);
    dp[0] = true;                              // bos dizge bolunebilir
    for (int i = 1; i <= n; ++i) {
        for (int j = 0; j < i; ++j) {
            if (dp[j] && dict.count(s.substr(j, i - j))) {
                dp[i] = true;                  // [0,j) bolunuyor ve [j,i) sozlukte
                break;
            }
        }
    }
    return dp[n];
}
\end{lstlisting}

\soru{Madeni Para -- Minimum (Coin Change): Sınırsız sayıda madeni para türü ve bir \code{amount} verildiğinde, tutarı oluşturmak için gereken minimum para sayısını bulun; mümkün değilse -1 döndürün.}

\cozum
\emph{Açgözlü} bölümünde gördüğümüz gibi, greedy bu problemde genelde yanılır; doğru çözüm DP'dir. \code{dp[x]}, $x$ tutarını ödemek için gereken minimum para olsun. \code{dp[0] = 0}. Her $x$ tutarı için tüm para türlerini deneriz: \code{dp[x] = min(dp[x], dp[x - coin] + 1)}. Yani ``$x - coin$'i optimal ödedikten sonra bir para daha ekle''. Ulaşılamayan tutarları ``sonsuz'' (örn. \code{amount + 1}) ile işaretleriz; sonuç hâlâ sonsuzsa -1 döneriz. Bu ``bottom-up'' (aşağıdan yukarı) yaklaşım tüm alt tutarları küçükten büyüğe çözer. Karmaşıklık \bigO{\text{amount} \times \text{coins}}.

\begin{lstlisting}
#include <vector>
#include <algorithm>
using namespace std;

int coinChange(vector<int>& coins, int amount) {
    vector<int> dp(amount + 1, amount + 1);   // "sonsuz" isaretci
    dp[0] = 0;                                 // 0 tutar icin 0 para
    for (int x = 1; x <= amount; ++x)
        for (int coin : coins)
            if (coin <= x)
                dp[x] = min(dp[x], dp[x - coin] + 1);
    return dp[amount] > amount ? -1 : dp[amount];
}
\end{lstlisting}

\soru{En Uzun Artan Alt Dizi (Longest Increasing Subsequence, LIS): Bir dizideki en uzun kesin artan alt dizinin uzunluğunu \bigO{n \log n} zamanda bulun.}

\cozum
Basit DP çözümü \bigO{n^2}'dir (\code{dp[i]} = $i$'de biten LIS). Daha hızlı yöntem, ``patience sorting'' fikrine dayanır: \code{tails[k]}, uzunluğu $k+1$ olan artan alt dizilerin bitebileceği \emph{en küçük son elemanı} tutar. Bu dizi her zaman sıralıdır. Her yeni eleman için, \code{tails} içinde onu yerleştirebileceğimiz ilk konumu ikili aramayla (\code{lower\_bound}) buluruz: o konumdaki değeri günceller (daha küçük bir sonla aynı uzunluğu koru) ya da diziyi bir uzatırız. \code{tails}'in nihai boyutu LIS uzunluğudur. İkili arama sayesinde \bigO{n \log n}.

Dikkat: \code{tails} dizisi \emph{gerçek} LIS'i içermez, yalnızca uzunluğunu verir; gerçek diziyi geri kurmak için ek işaretçiler gerekir.

\begin{lstlisting}
#include <vector>
#include <algorithm>
using namespace std;

int lengthOfLIS(vector<int>& nums) {
    vector<int> tails;                          // tails[k]: uzunluk k+1 icin en kucuk son
    for (int x : nums) {
        // x'i yerlestirebilecegimiz ilk konum (kesin artan icin lower_bound)
        auto it = lower_bound(tails.begin(), tails.end(), x);
        if (it == tails.end())
            tails.push_back(x);                 // LIS'i bir uzat
        else
            *it = x;                            // daha kucuk sonla ayni uzunlugu koru
    }
    return (int)tails.size();
}
\end{lstlisting}

\uyari{Kesin artan (strictly increasing) için \code{lower\_bound}, artmayan/azalmayan (non-decreasing) için \code{upper\_bound} kullanılır. Yanlış olanı seçmek eşit elemanlarda bir birim hataya yol açar.}

\gsub{String}

\soru{KMP (Knuth-Morris-Pratt) algoritmasıyla bir metinde bir örüntünün (pattern) tüm eşleşme başlangıçlarını \bigO{n + m} zamanda bulun.}

\cozum
Kaba kuvvet arama, uyumsuzlukta metinde bir adım geri kayar ve tekrar dener; bu \bigO{nm}'dir. KMP, örüntünün kendi iç yapısından yararlanır: uyumsuzluk olduğunda, örüntünün başına dönmek yerine, o ana kadar eşleşen ön ekin \emph{en uzun uygun ön-son eki} (longest proper prefix which is also suffix, LPS) kadar geri atlar. LPS dizisini önceden \bigO{m}'de hesaplarız; sonra metni tek geçişte tararız, metin işaretçisi hiç geri gitmez. Toplam \bigO{n + m}.

LPS[i], örüntünün $[0..i]$ ön ekinin, hem baş hem son olan en uzun (kendisi hariç) parçasının uzunluğudur. Bu, ``uyumsuzlukta örüntünün neresinden devam edebilirim'' sorusunun cevabıdır.

\begin{lstlisting}
#include <string>
#include <vector>
using namespace std;

vector<int> buildLPS(const string& p) {
    vector<int> lps(p.size(), 0);
    int len = 0;                                // onceki en uzun on-son eki
    for (int i = 1; i < (int)p.size(); ) {
        if (p[i] == p[len]) lps[i++] = ++len;   // uzat
        else if (len > 0)   len = lps[len - 1]; // geri atla (yeniden dene)
        else                lps[i++] = 0;       // eslesme yok
    }
    return lps;
}

vector<int> kmpSearch(const string& text, const string& pattern) {
    vector<int> lps = buildLPS(pattern), matches;
    int i = 0, j = 0;                           // i: metin, j: oruntu
    while (i < (int)text.size()) {
        if (text[i] == pattern[j]) { i++; j++; }
        if (j == (int)pattern.size()) {         // tam eslesme
            matches.push_back(i - j);
            j = lps[j - 1];                     // sonraki eslesmeyi ara
        } else if (i < (int)text.size() && text[i] != pattern[j]) {
            if (j > 0) j = lps[j - 1];          // oruntude geri atla
            else i++;                            // metinde ilerle
        }
    }
    return matches;
}
\end{lstlisting}

\soru{Rabin-Karp / rolling hash: Bir metinde bir örüntüyü, kayan (rolling) hash kullanarak ortalama \bigO{n + m} zamanda arayın.}

\cozum
Fikir, örüntünün hash değerini bir kez hesaplayıp, metin üzerinde $m$ uzunluklu her pencerenin hash'ini karşılaştırmaktır. Her pencereyi sıfırdan hash'lemek \bigO{m} olurdu; ``rolling hash'' ise pencereyi bir karakter kaydırırken hash'i \bigO{1}'de günceller: baştaki karakterin katkısını çıkar, hepsini tabana (base) çarp, yeni karakteri ekle. Bu bir polinom hash'idir: $h = \sum c_i \cdot b^{m-1-i} \bmod M$. Hash'ler eşitse (olası çakışmaya -- ``spurious hit'' -- karşı) karakterleri doğrulamak için birebir karşılaştırma yaparız. Büyük bir asal \code{MOD} ve iyi bir \code{BASE} ile çakışma nadirdir; ortalama \bigO{n + m}, en kötü (kötü hash) \bigO{nm}.

\begin{lstlisting}
#include <string>
#include <vector>
using namespace std;

vector<int> rabinKarp(const string& text, const string& pattern) {
    const long long BASE = 256, MOD = 1000000007;
    int n = (int)text.size(), m = (int)pattern.size();
    vector<int> matches;
    if (m > n) return matches;

    long long patHash = 0, winHash = 0, power = 1;
    for (int i = 0; i < m; ++i) {                    // ilk pencere + oruntu hash'i
        patHash = (patHash * BASE + pattern[i]) % MOD;
        winHash = (winHash * BASE + text[i]) % MOD;
        if (i < m - 1) power = (power * BASE) % MOD;  // BASE^(m-1)
    }
    for (int i = 0; i + m <= n; ++i) {
        if (patHash == winHash &&
            text.compare(i, m, pattern) == 0)         // dogrula (cakismaya karsi)
            matches.push_back(i);
        if (i + m < n) {                              // pencereyi kaydir
            winHash = (winHash - text[i] * power % MOD + MOD) % MOD; // basi cikar
            winHash = (winHash * BASE + text[i + m]) % MOD;         // yeni karakteri ekle
        }
    }
    return matches;
}
\end{lstlisting}

\bilgi{Tek modüllü hash'e karşı özel olarak tasarlanmış girdiler çakışma saldırısı yapabilir. Yarışma programlamada güvenlik için genelde iki farklı \code{(BASE, MOD)} çifti kullanılıp her iki hash de eşleşince kabul edilir (``double hashing'').}

\soru{En uzun palindromik alt dizge (longest palindromic substring): Merkezden genişleme (expand around center) yöntemiyle bulun.}

\cozum
Her palindromun bir merkezi vardır. Tek uzunluklu palindromların merkezi tek bir karakter, çift uzunluklu palindromların merkezi iki karakter arasıdır. $n$ karakterlik dizgede $2n - 1$ olası merkez vardır. Her merkezden başlayıp sol ve sağ işaretçileri dışa doğru açarak, karakterler eşit olduğu sürece palindromu genişletiriz. En uzun bulunanı saklarız. Her genişleme \bigO{n}, toplam merkez $2n-1$ olduğu için \bigO{n^2} zaman, \bigO{1} bellek. (Manacher algoritması \bigO{n}'dir ama expand-around-center hem yazması hem anlaması çok daha kolaydır.)

\begin{lstlisting}
#include <string>
using namespace std;

// left, right merkezinden disa dogru genisle; palindrom uzunlugunu dondur
int expand(const string& s, int left, int right) {
    while (left >= 0 && right < (int)s.size() && s[left] == s[right]) {
        left--; right++;
    }
    return right - left - 1;   // son gecerli uzunluk
}

string longestPalindrome(const string& s) {
    if (s.empty()) return "";
    int start = 0, maxLen = 1;
    for (int i = 0; i < (int)s.size(); ++i) {
        int len1 = expand(s, i, i);       // tek uzunluklu merkez
        int len2 = expand(s, i, i + 1);   // cift uzunluklu merkez
        int len = max(len1, len2);
        if (len > maxLen) {
            maxLen = len;
            start = i - (len - 1) / 2;    // baslangic indeksini geri hesapla
        }
    }
    return s.substr(start, maxLen);
}
\end{lstlisting}

\ipucu{\code{start = i - (len - 1) / 2} formülü hem tek hem çift uzunluklu palindromlarda doğru başlangıcı verir. Küçük bir örnekle (örn. \code{"babad"} ve \code{"cbbd"}) elle doğrulamak, bu indeks aritmetiğini içselleştirmenin en iyi yoludur.}

\gsub{Bit Manipülasyonu}

\soru{Bir tam sayının ikili gösteriminde kaç bitin set (1) olduğunu (popcount / Hamming weight) sayın. Bit sayısı kadar dönen verimli bir yöntem kullanın.}

\cozum
Naif yöntem her bit için 32 tur döner. Brian Kernighan numarası ise yalnızca \emph{set olan bit sayısı} kadar döner: \code{n \& (n - 1)} ifadesi, \code{n}'in en sağdaki (en düşük anlamlı) set bitini temizler. Neden? \code{n - 1}, en sağdaki 1'i 0 yapar ve ondan sağdaki tüm 0'ları 1'e çevirir; \code{n} ile AND'lediğimizde o en sağdaki 1 ve sağı sıfırlanır. Bu işlemi \code{n} sıfır olana dek tekrarlayıp kaç kez yaptığımızı sayarız. Karmaşıklık: set bit sayısıyla orantılı, en kötü \bigO{\text{bit sayısı}}.

\begin{lstlisting}
int countSetBits(unsigned int n) {
    int count = 0;
    while (n) {
        n &= (n - 1);   // en sagdaki set biti temizle
        count++;
    }
    return count;
    // Alternatif: C++20'de std::popcount(n), derleyici built-in'i __builtin_popcount(n)
}
\end{lstlisting}

\soru{Tek Eleman (Single Number): Bir dizide her eleman iki kez geçiyor, sadece biri bir kez. Sabit bellekle ve tek geçişte o tek elemanı bulun.}

\cozum
XOR işleminin iki temel özelliğini kullanırız: (1) $x \oplus x = 0$ (bir sayı kendisiyle XOR'lanınca sıfırlanır), (2) $x \oplus 0 = x$. Ayrıca XOR birleşmeli (associative) ve değişmelidir (commutative), yani sıra önemsizdir. Tüm elemanları XOR'ladığımızda, iki kez geçen her eleman birbirini götürüp 0 olur; geriye yalnızca tek geçen eleman kalır. \bigO{n} zaman, \bigO{1} bellek -- hash tablosuna bile gerek yok.

\begin{lstlisting}
#include <vector>
using namespace std;

int singleNumber(vector<int>& nums) {
    int result = 0;
    for (int x : nums)
        result ^= x;    // eslesenler birbirini goturur, tek kalan doner
    return result;
}
\end{lstlisting}

\bilgi{Aynı fikir varyantlara genişler: iki tek eleman varsa, önce hepsini XOR'layıp sonucun herhangi bir set bitine göre diziyi iki gruba ayırırsınız. ``Her eleman üç kez, biri bir kez'' varyantı ise bit bazında $\bmod 3$ sayımıyla çözülür.}

\soru{Bir kümenin tüm alt kümelerini bitmask (bit maskesi) ile üretin. $n$ elemanlı küme için $2^n$ alt kümeyi döngüyle listeleyin.}

\cozum
$n$ elemanlı bir kümenin her alt kümesini, $n$ bitlik bir sayının ikili gösterimiyle birebir eşleyebiliriz: $i$. bit set ise $i$. eleman alt kümededir. 0'dan $2^n - 1$'e kadar tüm sayıları gezersek tüm alt kümeleri üretmiş oluruz. Bir \code{mask} için hangi elemanların dahil olduğunu, her bit pozisyonunu \code{mask \& (1 << i)} ile kontrol ederek buluruz. Bu, backtracking'e göre daha az kod ve daha az özyineleme yükü demektir. Karmaşıklık \bigO{n \cdot 2^n}. Bit sayacı olarak \code{unsigned}/\code{long long} kullanın; $n \le 63$ ile sınırlıdır.

\begin{lstlisting}
#include <vector>
using namespace std;

vector<vector<int>> allSubsets(vector<int>& nums) {
    int n = (int)nums.size();
    vector<vector<int>> result;
    for (int mask = 0; mask < (1 << n); ++mask) {  // 0 .. 2^n - 1
        vector<int> subset;
        for (int i = 0; i < n; ++i)
            if (mask & (1 << i))                    // i. bit set mi?
                subset.push_back(nums[i]);
        result.push_back(subset);
    }
    return result;
}
\end{lstlisting}

\ipucu{Bir maskenin \emph{tüm alt maskelerini} (submask) gezmek için özel bir kalıp vardır: \code{for (int sub = mask; sub > 0; sub = (sub - 1) \& mask)}. Bu, alt küme-toplamı DP'lerinde çok işe yarar ve tüm mask/submask çiftlerini \bigO{3^n} toplamda gezer.}

\soru{İki'nin Kuvveti mi? Bir pozitif tam sayının tam olarak 2'nin bir kuvveti olup olmadığını tek bir bit işlemiyle kontrol edin.}

\cozum
2'nin kuvvetleri (1, 2, 4, 8, ...) ikilik tabanda tam olarak \emph{bir} set bite sahiptir: \code{100...0}. Böyle bir sayıdan 1 çıkarınca, o tek set bit 0 olur ve sağındaki tüm bitler 1'e döner: örn. $8 = 1000$, $7 = 0111$. Bu ikisini AND'lersek sonuç 0 olur. Birden fazla set biti olan sayılarda ise AND sıfır çıkmaz. Sıfır ve negatifleri elemek için \code{n > 0} kontrolünü eklemek şarttır (aksi halde \code{n = 0} yanlışlıkla ``kuvvet'' sanılır). Tek karşılaştırma, \bigO{1}.

\begin{lstlisting}
bool isPowerOfTwo(int n) {
    // n > 0 olmali; tek set bitli sayilar icin n & (n-1) == 0
    return n > 0 && (n & (n - 1)) == 0;
}
\end{lstlisting}

\notk[Bit numaraları özeti]{\code{n \& (n-1)} en sağdaki set biti siler (popcount, 2-kuvvet testi). \code{n \& (-n)} en sağdaki set biti \emph{izole eder} (Fenwick/BIT ağaçlarında kullanılır). \code{n | (n+1)} en sağdaki 0'ı set eder. Bu küçük özdeşlikleri ezberlemek, düşük seviyeli optimizasyonlarda ve yarışma problemlerinde büyük kolaylık sağlar.}

\gsec{Kapanış: Doğru Aracı Seçmek}
% =====================================================================

Bu rehber boyunca veri yapılarını ve algoritmaları tek tek gördük. Gerçek
problemlerde asıl beceri, hangi aracın ne zaman kullanılacağını sezmektir.
Aşağıdaki tablo hızlı bir başvuru sunar.

\begin{center}
\begin{tabularx}{\textwidth}{@{}>{\bfseries\color{anaMavi}}l X@{}}
\toprule
\rowcolor{acikMavi}\color{anaMavi}\textbf{İhtiyaç} & \color{anaMavi}\textbf{Uygun araç}\\
\midrule
Hızlı indeksli erişim & Dizi / \code{vector}\\
\rowcolor{acikGri} Sık baş/orta ekleme-silme & Bağlı liste / \code{list}\\
Anahtarla \bigO{1} arama & Hash tablosu / \code{unordered\_map}\\
\rowcolor{acikGri} Sıralı anahtar + aralık sorgusu & Dengeli ağaç / \code{map}\\
En büyük/küçüğe hızlı erişim & Heap / \code{priority\_queue}\\
\rowcolor{acikGri} Önek arama, otomatik tamamlama & Trie\\
Aralık toplamı + güncelleme & Segment / Fenwick ağacı\\
\rowcolor{acikGri} En kısa yol (ağırlıksız) & BFS\\
En kısa yol (pozitif ağırlık) & Dijkstra\\
\rowcolor{acikGri} Bağlantı/döngü sorgusu & Union-Find (DSU)\\
Örtüşen alt problemler & Dinamik programlama\\
\rowcolor{acikGri} Tüm kombinasyonlar + budama & Geri izleme\\
\bottomrule
\end{tabularx}
\end{center}

\ipucu{Bir probleme yaklaşırken sıra: (1) \emph{Girdi boyutuna} bak ---
hedeflenen karmaşıklığı söyler. (2) Problemi bildiğin bir kalıba benzet.
(3) Basit (brute-force) çözümü yaz, sonra veri yapısı veya DP ile hızlandır.
(4) Kenar durumlarını (boş girdi, tek eleman, taşma) test et. Bol pratik,
bu sezgiyi geliştiren tek yoldur.}

\vfill
\begin{center}
{\color{ortaGri}\rule{0.5\linewidth}{0.6pt}}\\[6pt]
{\color{ortaGri}\small Rehberin sonu --- İyi kodlamalar!}
\end{center}

\end{document}
